\documentclass[12pt,a4paper,oneside]{report}

% =========================================================
% MAIN SOURCE - AI-EDU PATHWAYS
% Phần III (SRS) là LaTeX source thực, không phải PDF wrapper.
% Các sơ đồ 3.1.1--3.1.10 + Context + Use Case lấy artifact gốc GitHub.
% Phần 3.2 đã tích hợp 25 ảnh chụp màn hình bổ sung ngày 2026-09-10;
% các ảnh ghép được tách thành vùng hiển thị riêng nhưng vẫn giữ file gốc.
% Phần I-II và IV-V được giữ từ PDF hiện tại vì source .tex tương ứng
% chưa có trong gói nguồn. Xem README_BUILD.md để biết cấu trúc và cách build.
% =========================================================

\usepackage{fontspec}
\IfFontExistsTF{Times New Roman}{\setmainfont{Times New Roman}}{\setmainfont{DejaVu Serif}}
\usepackage[provide=*,vietnamese]{babel}
\usepackage[top=2.5cm,bottom=2.5cm,left=3cm,right=2cm]{geometry}
\usepackage{setspace}
\onehalfspacing
\usepackage{graphicx}
\usepackage{pdfpages}
\usepackage{longtable}
\usepackage{array}
\usepackage{tabularx}
\usepackage{booktabs}
\usepackage[table]{xcolor}
\usepackage{caption}
\usepackage{float}
\usepackage{needspace}
\usepackage{enumitem}
\usepackage{amsmath}
\usepackage[hidelinks]{hyperref}
\renewcommand{\contentsname}{Mục lục}
\setcounter{secnumdepth}{3}
\setcounter{tocdepth}{3}

% GitHub artifacts đã được đóng gói sẵn trong thư mục figures/github.
\newcommand{\GitHubMissingBox}[1]{%
  \fbox{\parbox{0.88\textwidth}{\centering\small
  \textbf{Thiếu artifact GitHub:} \texttt{\detokenize{#1}}\\[3pt]
  Chạy \texttt{PREPARE\_GITHUB\_ASSETS.bat} rồi compile lại bằng XeLaTeX.}}%
}
\newcommand{\GitHubIncludeGraphics}[2][]{%
  \IfFileExists{#2}{\includegraphics[#1]{#2}}{\GitHubMissingBox{#2}}%
}

% Nếu một ảnh giao diện bị thiếu khỏi gói nguồn, hiển thị khung cảnh báo ngay
% trong PDF để việc lệch giữa main.tex và artifact được phát hiện khi build.
\newcommand{\SafeIncludeGraphics}[2][]{%
  \IfFileExists{#2}{\includegraphics[#1]{#2}}{%
    \fbox{\parbox{0.86\textwidth}{\centering\small
      Ảnh giao diện chưa có trong source package: \texttt{\detokenize{#2}}}}%
  }%
}


\newcommand{\ScreenLayoutItem}{%
  \Needspace{0.43\textheight}%
  \item \textbf{Screen layout:}%
}

% Mỗi mục Screen layout hiển thị ảnh trực tiếp của đúng chức năng. Ảnh chỉ
% được dùng chung khi các thao tác thực sự cùng xuất hiện trong một màn hình.
\newcommand{\ScreenLayoutFigure}[2]{%
  \begin{figure}[H]
    \centering
    \SafeIncludeGraphics[width=0.94\textwidth,height=0.70\textheight,keepaspectratio]{#1}
    \caption{Giao diện dùng cho chức năng \textit{#2}.}
  \end{figure}
}

\begin{document}

% ---------------------------------------------------------
% I + II: giữ nguyên từ báo cáo hiện có vì source .tex gốc chưa có.
% ---------------------------------------------------------
\includepdf[pages=1-21,fitpaper=true,pagecommand={\thispagestyle{empty}}]{base/AI_EDU_PATHWAYS_FINAL_289P_DIAGRAM_QA.pdf}
\clearpage

% =========================================================
% III. SOFTWARE REQUIREMENT SPECIFICATION - PHẦN 1, 2 + 3.1
% Phần 1 và 2 lấy từ nội dung KAN-15 + Use Case trên GitHub.
% Phần 3.1 giữ theo bản tích hợp đã thống nhất trước đó.
% =========================================================

\chapter*{III. Software Requirement Specification}
\addcontentsline{toc}{chapter}{III. Software Requirement Specification}

% =========================================================
% 1. PRODUCT OVERVIEW
% =========================================================
\section*{1. Product Overview}
\addcontentsline{toc}{section}{1. Product Overview}

\newpage
\begin{figure}[htbp]
    \centering
    \GitHubIncludeGraphics[width=0.92\textwidth,height=0.76\textheight,keepaspectratio]{figures/github/context_diagram.png}
    \caption{Context Diagram tổng quan hệ thống AI-Edu Pathways.}
    \label{fig:context-diagram-ai-edu}
\end{figure}

Hệ thống \textbf{Cá nhân hóa lộ trình học và Luyện thi Đánh giá năng lực tích hợp AI} là một nền tảng giáo dục trên nền web (web-based) toàn diện. Hệ thống được thiết kế nhằm số hóa và tối ưu hóa quá trình ôn luyện, hỗ trợ đánh giá dựa trên năng lực chuẩn và tích hợp sâu các công nghệ AI tiên tiến (LLM + RAG) để đồng hành cùng Học sinh, Giáo viên, Phụ huynh và Quản trị viên trong suốt quá trình học tập.

\begin{itemize}
    \item \textbf{Môi trường hoạt động (Operational Environment):}
    \begin{itemize}
        \item \textbf{Nền tảng (Platform):} Ứng dụng Web-based đa nền tảng, truy cập mượt mà qua các trình duyệt web hiện đại (Chrome, Edge, Firefox, Safari).
        \item \textbf{Kết nối (Connectivity):} Yêu cầu kết nối Internet liên tục và ổn định (không hỗ trợ chế độ offline) để xử lý các luồng dữ liệu đánh giá và truy vấn AI theo thời gian thực.
        \item \textbf{Hạ tầng (Infrastructure):} Triển khai trên nền tảng điện toán đám mây (Cloud-based) với kiến trúc linh hoạt, hỗ trợ tốt các API dịch vụ AI, Cơ sở dữ liệu Vector (Vector Database) và hệ thống Gateway gửi thông báo.
    \end{itemize}

    \item \textbf{Người dùng dự kiến (Anticipated Users):}
    \begin{itemize}
        \item \textbf{Người dùng nội bộ (Internal Users):} Quản trị viên Học vụ (Academic Administrator), Giáo viên / Giảng viên chuyên môn (Teacher / Lecturer).
        \item \textbf{Người dùng bên ngoài (External Users):} Học sinh / Sinh viên (Khách hàng chính), Phụ huynh (Parent).
        \item \textbf{Tác nhân hệ thống ngoài (External Systems):} Các dịch vụ API AI, Cổng thanh toán, Tổ chức Giáo dục.
    \end{itemize}

    \item \textbf{Ràng buộc \& Giả định (Constraints \& Assumptions):}
    \begin{itemize}
        \item Hệ thống bám sát phương pháp học tập và xây dựng lộ trình cá nhân hóa dựa trên chuẩn năng lực cốt lõi (competency-based).
        \item Các phân tích và đề xuất từ AI chỉ đóng vai trò như một ``Tutor'' hỗ trợ (Learning Assistance), tuyệt đối không thay thế các quyết định mang tính chuyên môn, học thuật của Giáo viên.
        \item Độ chính xác của kết quả đánh giá phụ thuộc chặt chẽ vào chất lượng ngân hàng câu hỏi và khung năng lực do các Tổ chức giáo dục đối tác cung cấp.
        \item Mọi giao dịch thanh toán trực tuyến phải được mã hóa và xử lý thông qua các cổng thanh toán nội địa uy tín.
    \end{itemize}

    \item \textbf{Tích hợp (Integrations):}
    \begin{itemize}
        \item Dịch vụ Trí tuệ nhân tạo (OpenAI / Gemini API).
        \item Cơ sở dữ liệu Vector chuyên dụng (Qdrant / Pinecone).
        \item Cổng thanh toán điện tử (VNPay / MoMo / ZaloPay).
        \item Dịch vụ truyền thông đa kênh (Email / SMS / Push Notification).
        \item Hệ thống dữ liệu của các Tổ chức Giáo dục cung cấp đề thi ĐGNL.
    \end{itemize}

    \item \textbf{Phụ thuộc (Dependencies):}
    \begin{itemize}
        \item Sự ổn định, tốc độ phản hồi (latency) và giới hạn tải (rate limits) của các API dịch vụ AI.
        \item Độ trễ và băng thông của Cơ sở dữ liệu Vector khi thực hiện tìm kiếm ngữ nghĩa (Semantic Search).
        \item Tính sẵn sàng của các API Cổng thanh toán và Dịch vụ Thông báo.
    \end{itemize}
\end{itemize}

% =========================================================
% 2. USER REQUIREMENTS
% =========================================================
\section*{2. User Requirements}
\addcontentsline{toc}{section}{2. User Requirements}

\subsection*{2.1 System Actors}
\addcontentsline{toc}{subsection}{2.1 System Actors}

\renewcommand{\arraystretch}{1.15}
\begin{longtable}{|c|p{3.8cm}|p{9.5cm}|}
\hline
\rowcolor{orange!15}
\textbf{\#} & \textbf{Tác nhân (Actor)} & \textbf{Mô tả vai trò \& Nhiệm vụ cốt lõi} \\
\hline
\endfirsthead
\hline
\rowcolor{orange!15}
\textbf{\#} & \textbf{Tác nhân (Actor)} & \textbf{Mô tả vai trò \& Nhiệm vụ cốt lõi} \\
\hline
\endhead

1 & \textbf{Quản trị viên Học vụ} \newline \textit{(Academic Admin)} & Chịu trách nhiệm vận hành toàn diện hệ thống: quản lý tài khoản, phân quyền truy cập (RBAC), quản lý nội dung học tập và thiết lập khung năng lực. Thực hiện cấu hình các chính sách và siêu tham số (hyperparameters) cho thuật toán AI. Giám sát luồng đồng bộ dữ liệu, theo dõi logs hoạt động và đánh giá hiệu suất tổng thể qua các bảng điều khiển (Dashboards) để xử lý sự cố. \\ \hline
2 & \textbf{Giáo viên / Giảng viên} \newline \textit{(Teacher)} & Đóng vai trò chuyên gia học thuật. Trực tiếp tạo lập, cấu hình và quản lý tài liệu học tập, ngân hàng câu hỏi, khung năng lực và kế hoạch giảng dạy. Có nhiệm vụ theo dõi báo cáo phân tích hiệu suất lớp học, kiểm duyệt các nội dung/lộ trình do AI tự động sinh ra và trực tiếp can thiệp (giao thêm bài, điều chỉnh mục tiêu, đánh giá cuối cùng) để đảm bảo chất lượng. \\ \hline
3 & \textbf{Học sinh / Sinh viên} \newline \textit{(Student)} & Đối tượng sử dụng cốt lõi (Khách hàng). Sử dụng hệ thống để đăng ký, thực hiện bài kiểm tra chẩn đoán năng lực ban đầu. Truy cập và học tập theo lộ trình cá nhân hóa, thực hành luyện đề thi thích ứng (Adaptive Test). Tương tác trực tiếp với AI Tutor để được hướng dẫn giải bài, nhận đề xuất học tập, xem báo cáo tiến độ và theo dõi các thành tựu (Achievements). \\ \hline
4 & \textbf{Phụ huynh} \newline \textit{(Parent)} & Người đồng hành và bảo trợ. Sử dụng hệ thống để nhận các thông báo, báo cáo định kỳ nhằm theo dõi tiến trình học tập của con em. Đặc biệt, có quyền thiết lập mục tiêu học tập, cấu hình giới hạn thời gian học (Time limits). Xem xét các đề xuất lộ trình từ AI và thực hiện giao dịch thanh toán (gia hạn/nâng cấp gói dịch vụ). \\ \hline
5 & \textbf{Dịch vụ AI} \newline \textit{(OpenAI / Gemini / Qdrant)} & Động cơ xử lý ngôn ngữ tự nhiên (NLP) và phân tích dữ liệu. Ứng dụng kỹ thuật truy xuất kiến thức theo ngữ cảnh (RAG - kết hợp Vector DB) để cung cấp tính năng AI Tutor. Có khả năng: xây dựng đồ thị lộ trình cá nhân hóa, cung cấp lời giải thích chi tiết từng bước, chấm điểm tự luận, đánh giá mức độ hiểu bài và dự đoán kết quả học tập. \\ \hline
6 & \textbf{Cổng Thanh toán} \newline \textit{(Payment Gateway)} & Tác nhân xử lý tài chính độc lập. Nhận yêu cầu giao dịch trực tuyến, thực hiện quy trình bảo mật (Tokenization) với ngân hàng, hỗ trợ xuất hóa đơn điện tử và tự động phản hồi qua Webhook/IPN về trạng thái giao dịch (Thành công/Thất bại) để hệ thống tức thời cấp quyền khóa học. \\ \hline
7 & \textbf{Tổ chức Giáo dục} \newline \textit{(Edu Organization)} & Nguồn cung cấp và chuẩn hóa dữ liệu. Cung cấp bài thi ĐGNL mới nhất, khung năng lực, ngân hàng câu hỏi gốc và các siêu dữ liệu (Metadata/Knowledge Tags). Đồng thời, tiếp nhận ngược lại các báo cáo phân tích, thống kê tổng hợp về phổ điểm và chất lượng câu hỏi từ hệ thống. \\ \hline
8 & \textbf{Dịch vụ Thông báo} \newline \textit{(Notification Service)} & Hệ thống trung gian tự động xử lý truyền thông đa kênh (Email, SMS, Push Notification). Chịu trách nhiệm quản lý lịch trình gửi, theo dõi trạng thái phân phối tin nhắn và thực hiện gửi các cảnh báo học tập, nhắc nhở hoạt động/lịch thi thử, mã OTP đăng nhập đến người dùng cuối. \\ \hline
\end{longtable}

\subsection*{2.2 Use Cases}
\addcontentsline{toc}{subsection}{2.2 Use Cases}

\subsubsection*{2.2.1 Use Case Diagram}
\addcontentsline{toc}{subsubsection}{2.2.1 Use Case Diagram}

\begin{figure}[htbp]
    \centering
    \GitHubIncludeGraphics[width=0.99\textwidth,height=0.80\textheight,keepaspectratio]{figures/github/usecase_diagram_tong_hop_latest.png}
    \caption{Use Case Diagram tổng hợp của hệ thống.}
    \label{fig:usecase-diagram-tong-hop}
\end{figure}

% =========================================================
% USE CASE DESCRIPTIONS - BẢNG MÔ TẢ TÓM TẮT
% Nguồn: branch develop / latex/chapters/03_use_case_descriptions.tex
% =========================================================

\subsubsection*{2.2.2 Use Case Descriptions}\addcontentsline{toc}{subsubsection}{2.2.2 Use Case Descriptions}
\label{subsubsec:use-case-descriptions}

\renewcommand{\arraystretch}{1.25}
\setlength{\tabcolsep}{5pt}
\small

\begin{longtable}{|p{0.09\textwidth}|p{0.25\textwidth}|p{0.20\textwidth}|p{0.38\textwidth}|}
\caption{Use Case Descriptions của hệ thống}
\label{tab:use-case-descriptions}\\
\hline
\rowcolor{orange!15}
\textbf{ID} & \textbf{Use Case} & \textbf{Actors} & \textbf{Use Case Description} \\
\hline
\endfirsthead

\hline
\rowcolor{orange!15}
\textbf{ID} & \textbf{Use Case} & \textbf{Actors} & \textbf{Use Case Description} \\
\hline
\endhead

\hline
\multicolumn{4}{r}{\textit{Tiếp tục ở trang sau}}\\
\endfoot

\hline
\endlastfoot

STU-01 &
Quản lý thông tin đăng ký và hồ sơ &
Học sinh / Sinh viên &
Cho phép người học đăng ký tài khoản, khai báo và cập nhật thông tin cá nhân, trường học và hồ sơ học tập. \\
\hline

STU-02 &
Thiết lập mục tiêu học tập và kỳ thi đích &
Học sinh / Sinh viên &
Cho phép người học thiết lập điểm mục tiêu, kỳ thi cần tham gia, thời gian chuẩn bị và thời lượng học dự kiến. \\
\hline

STU-03 &
Tương tác với AI Tutor &
Học sinh / Sinh viên &
Cho phép người học gửi câu hỏi, trao đổi với AI Tutor và nhận hỗ trợ giải thích kiến thức trong quá trình học. \\
\hline

STU-04 &
Gửi phản hồi và đánh giá nội dung &
Học sinh / Sinh viên &
Cho phép người học đánh giá bài học, câu hỏi, lời giải và gửi phản hồi về chất lượng nội dung. \\
\hline

STU-05 &
Yêu cầu điều chỉnh lộ trình &
Học sinh / Sinh viên &
Cho phép người học đề nghị hệ thống điều chỉnh lộ trình khi mục tiêu, thời gian học hoặc năng lực hiện tại thay đổi. \\
\hline

STU-06 &
Xem lộ trình học cá nhân hóa &
Học sinh / Sinh viên &
Hiển thị kế hoạch học tập cá nhân hóa theo năng lực, mục tiêu, mức độ ưu tiên và thời gian còn lại trước kỳ thi. \\
\hline

STU-07 &
Làm đề luyện thi thích ứng &
Học sinh / Sinh viên &
Cung cấp bài luyện tập hoặc đề thi có độ khó và nội dung được điều chỉnh theo năng lực hiện tại của người học. \\
\hline

STU-08 &
Xem kết quả và phân tích năng lực &
Học sinh / Sinh viên &
Hiển thị điểm số, kết quả theo nhóm kiến thức, điểm mạnh, điểm yếu và mức độ tiến bộ của người học. \\
\hline

STU-09 &
Nhận phản hồi AI Tutor (LLM + RAG) &
Học sinh / Sinh viên &
Cung cấp câu trả lời và lời giải có căn cứ từ AI Tutor bằng mô hình ngôn ngữ kết hợp truy xuất tài liệu RAG. \\
\hline

STU-10 &
Xem dự đoán điểm thi và xếp hạng &
Học sinh / Sinh viên &
Hiển thị điểm thi dự đoán, xác suất đạt mục tiêu và vị trí tương đối dựa trên dữ liệu học tập hiện có. \\
\hline

STU-11 &
Xem gợi ý tài liệu bổ trợ &
Học sinh / Sinh viên &
Đề xuất bài học, tài liệu, câu hỏi và nội dung ôn tập phù hợp với phần kiến thức còn yếu. \\
\hline

STU-12 &
Xem báo cáo tiến độ và nhắc nhở &
Học sinh / Sinh viên &
Hiển thị tiến độ hoàn thành lộ trình, lịch học, nhiệm vụ sắp tới và các nhắc nhở học tập cần thiết. \\
\hline

TEA-01 &
Quản lý nội dung bài giảng và tài liệu &
Giáo viên / Giảng viên &
Cho phép giáo viên tạo, cập nhật, phân loại và quản lý bài giảng, tài liệu học tập dùng trong hệ thống. \\
\hline

TEA-02 &
Quản lý ngân hàng câu hỏi và đáp án &
Giáo viên / Giảng viên &
Cho phép giáo viên thêm, sửa, xóa, duyệt và phân loại câu hỏi, đáp án và lời giải trong ngân hàng câu hỏi. \\
\hline

TEA-03 &
Đánh giá và nhận xét học sinh &
Giáo viên / Giảng viên &
Cho phép giáo viên xem kết quả, đưa ra nhận xét và đánh giá mức độ tiến bộ của từng học sinh. \\
\hline

TEA-04 &
Quản lý ma trận kiến thức và tiêu chí &
Giáo viên / Giảng viên &
Cho phép giáo viên xây dựng và cập nhật ma trận kiến thức, tiêu chí đánh giá và mức độ khó của nội dung. \\
\hline

TEA-05 &
Cập nhật chương trình giảng dạy &
Giáo viên / Giảng viên &
Cho phép giáo viên điều chỉnh cấu trúc môn học, nội dung, thứ tự bài học và mục tiêu giảng dạy. \\
\hline

TEA-06 &
Xem Dashboard Learning Analytics &
Giáo viên / Giảng viên &
Cung cấp dashboard tổng hợp dữ liệu học tập, mức độ tham gia, kết quả và xu hướng tiến bộ của học sinh. \\
\hline

TEA-07 &
Xem báo cáo tiến độ học sinh &
Giáo viên / Giảng viên &
Hiển thị tiến độ lộ trình, tần suất học, kết quả luyện tập và mức độ hoàn thành mục tiêu của học sinh. \\
\hline

TEA-08 &
Xem thống kê hiệu quả giảng dạy &
Giáo viên / Giảng viên &
Cung cấp số liệu về mức độ sử dụng tài liệu, tỷ lệ hoàn thành, kết quả học tập và hiệu quả nội dung giảng dạy. \\
\hline

TEA-09 &
Xem kết quả Knowledge Tracing &
Giáo viên / Giảng viên &
Hiển thị mức độ thành thạo kiến thức được mô hình Knowledge Tracing ước lượng cho từng học sinh. \\
\hline

TEA-10 &
Xem đề xuất nội dung bởi AI &
Giáo viên / Giảng viên &
Cung cấp các đề xuất của AI về tài liệu, câu hỏi, độ khó và nội dung cần bổ sung hoặc điều chỉnh. \\
\hline

PAR-01 &
Đăng ký tài khoản giám sát &
Phụ huynh &
Cho phép phụ huynh đăng ký tài khoản và liên kết với hồ sơ của học sinh cần theo dõi. \\
\hline

PAR-02 &
Yêu cầu báo cáo tiến độ &
Phụ huynh &
Cho phép phụ huynh yêu cầu hệ thống tạo hoặc gửi báo cáo học tập của học sinh theo khoảng thời gian lựa chọn. \\
\hline

PAR-03 &
Thiết lập giới hạn thời gian &
Phụ huynh &
Cho phép phụ huynh thiết lập thời lượng học, khung giờ học và giới hạn sử dụng hệ thống của học sinh. \\
\hline

PAR-04 &
Xem báo cáo tiến độ con em &
Phụ huynh &
Hiển thị tiến độ học tập, mức độ hoàn thành lộ trình, thời gian học và kết quả gần đây của học sinh. \\
\hline

PAR-05 &
Xem kết quả thi và dự đoán điểm &
Phụ huynh &
Hiển thị kết quả bài kiểm tra, đề thi thử và điểm thi dự đoán của học sinh. \\
\hline

PAR-06 &
Xem khuyến nghị hỗ trợ học tập &
Phụ huynh &
Cung cấp các khuyến nghị giúp phụ huynh hỗ trợ học sinh về thời gian, nội dung và phương pháp học. \\
\hline

PAR-07 &
Nhận cảnh báo suy giảm kết quả &
Phụ huynh &
Gửi cảnh báo khi kết quả, tần suất học hoặc mức độ hoàn thành lộ trình của học sinh có dấu hiệu suy giảm. \\
\hline

PAR-08 &
Xem lịch học và hoạt động gần đây &
Phụ huynh &
Hiển thị lịch học, nhiệm vụ sắp tới và các hoạt động học tập gần đây của học sinh. \\
\hline

NOT-01 &
Theo dõi trạng thái gửi thông báo &
Dịch vụ thông báo &
Theo dõi trạng thái đã gửi, đang xử lý, gửi thành công hoặc thất bại của từng thông báo. \\
\hline

NOT-02 &
Xem báo cáo lỗi gửi &
Dịch vụ thông báo &
Ghi nhận và cung cấp thông tin lỗi phát sinh trong quá trình gửi Email, SMS hoặc Push Notification. \\
\hline

NOT-03 &
Gửi nội dung Email / SMS / Push &
Dịch vụ thông báo &
Gửi thông báo học tập, cảnh báo và nhắc nhở đến người dùng qua Email, SMS hoặc Push Notification. \\
\hline

NOT-04 &
Lập lịch gửi thông báo &
Dịch vụ thông báo &
Thiết lập thời điểm và tần suất gửi tự động cho các loại thông báo và nhắc nhở. \\
\hline

ADM-01 &
Quản lý tài khoản và phân quyền &
Quản trị viên học vụ &
Cho phép quản trị viên tạo, cập nhật, khóa tài khoản và gán quyền truy cập theo từng vai trò. \\
\hline

ADM-02 &
Cấu hình hệ thống và chính sách &
Quản trị viên học vụ &
Cho phép quản trị viên thiết lập tham số hệ thống, chính sách sử dụng, bảo mật và quy tắc vận hành. \\
\hline

ADM-03 &
Quản lý danh mục nội dung &
Quản trị viên học vụ &
Cho phép quản trị viên quản lý môn học, chủ đề, khối lớp, loại tài liệu và các danh mục dùng chung. \\
\hline

ADM-04 &
Thiết lập quy tắc AI &
Quản trị viên học vụ &
Cho phép quản trị viên cấu hình quy tắc đề xuất, ngưỡng cảnh báo, giới hạn sử dụng và hành vi của các chức năng AI. \\
\hline

ADM-05 &
Xem Dashboard tổng quan &
Quản trị viên học vụ &
Hiển thị số liệu tổng quan về người dùng, nội dung, hoạt động học tập, hệ thống và các dịch vụ AI. \\
\hline

ADM-06 &
Xem báo cáo phân tích xu hướng &
Quản trị viên học vụ &
Cung cấp báo cáo về xu hướng sử dụng, kết quả học tập, hành vi người dùng và hiệu quả vận hành hệ thống. \\
\hline

ADM-07 &
Xem nhật ký truy cập và bảo mật &
Quản trị viên học vụ &
Cho phép quản trị viên kiểm tra lịch sử đăng nhập, hành động người dùng, sự kiện bảo mật và truy cập bất thường. \\
\hline

ADM-08 &
Theo dõi hiệu suất mô hình AI &
Quản trị viên học vụ &
Theo dõi độ trễ, tỷ lệ lỗi, mức độ sử dụng và các chỉ số hiệu suất của các mô hình AI. \\
\hline

AI-01 &
Gửi Prompt và ngữ cảnh AI Tutor &
Dịch vụ AI &
Gửi câu hỏi, lịch sử hội thoại và ngữ cảnh học tập từ hệ thống đến dịch vụ AI để xử lý. \\
\hline

AI-02 &
Nhận phản hồi AI Tutor &
Dịch vụ AI &
Tiếp nhận câu trả lời, hướng dẫn và nội dung hỗ trợ học tập do AI Tutor sinh ra. \\
\hline

AI-03 &
Gửi yêu cầu sinh câu hỏi &
Dịch vụ AI &
Gửi chủ đề, mức độ khó, dạng câu hỏi và yêu cầu nội dung đến dịch vụ AI để tạo câu hỏi. \\
\hline

AI-04 &
Nhận câu hỏi sinh tự động &
Dịch vụ AI &
Tiếp nhận câu hỏi, đáp án và lời giải do dịch vụ AI tạo để giáo viên hoặc hệ thống xem xét. \\
\hline

AI-05 &
Gửi Context tài liệu RAG &
Dịch vụ AI &
Cung cấp truy vấn và phạm vi tài liệu cần truy xuất để dịch vụ RAG tìm kiếm ngữ cảnh phù hợp. \\
\hline

AI-06 &
Nhận kết quả tìm kiếm RAG &
Dịch vụ AI &
Tiếp nhận các đoạn tài liệu, nguồn tham chiếu và mức độ liên quan từ hệ thống RAG. \\
\hline

AI-07 &
Nhận giải thích chi tiết lời giải &
Dịch vụ AI &
Tiếp nhận phần giải thích từng bước, kiến thức liên quan và nguyên nhân đúng hoặc sai của đáp án. \\
\hline

AI-08 &
Nhận dự đoán kết quả thi &
Dịch vụ AI &
Tiếp nhận điểm dự đoán, xác suất đạt mục tiêu và các yếu tố rủi ro do mô hình AI tính toán. \\
\hline

AI-09 &
Tạo Embedding và truy vấn Vector DB &
Dịch vụ AI &
Chuyển đổi nội dung thành vector và thực hiện truy vấn trên cơ sở dữ liệu vector để tìm tài liệu liên quan. \\
\hline

AI-10 &
Gửi dữ liệu huấn luyện mô hình &
Dịch vụ AI &
Cung cấp dữ liệu học tập đã được kiểm soát để phục vụ huấn luyện hoặc cập nhật mô hình AI. \\
\hline

PAY-01 &
Nhận hóa đơn và biên lai điện tử &
Cổng thanh toán &
Tiếp nhận và hiển thị hóa đơn hoặc biên lai điện tử sau khi giao dịch được hoàn tất. \\
\hline

PAY-02 &
Xác nhận thanh toán &
Cổng thanh toán &
Nhận kết quả xác nhận thành công hoặc thất bại của giao dịch từ cổng thanh toán. \\
\hline

PAY-03 &
Kiểm tra trạng thái giao dịch &
Cổng thanh toán &
Kiểm tra trạng thái đang xử lý, thành công, thất bại hoặc hoàn tiền của một giao dịch. \\
\hline

PAY-04 &
Gửi yêu cầu thanh toán &
Cổng thanh toán &
Gửi thông tin đơn hàng, số tiền và phương thức thanh toán đến VNPay, MoMo hoặc ZaloPay. \\
\hline

PAY-05 &
Xem thông tin gói dịch vụ &
Học sinh / Phụ huynh &
Hiển thị tên gói, thời hạn, phạm vi chức năng và chi phí của các gói dịch vụ. \\
\hline

PAY-06 &
Gửi yêu cầu hoàn tiền &
Cổng thanh toán &
Gửi yêu cầu hoàn trả một phần hoặc toàn bộ giá trị giao dịch theo chính sách của hệ thống. \\
\hline

EDU-01 &
Đồng bộ đề thi mẫu ĐGNL &
Tổ chức giáo dục / Đơn vị ra đề ĐGNL &
Tiếp nhận và cập nhật các đề thi mẫu ĐGNL từ đơn vị cung cấp dữ liệu chính thức. \\
\hline

EDU-02 &
Đồng bộ cấu trúc và ma trận đề thi &
Tổ chức giáo dục / Đơn vị ra đề ĐGNL &
Cập nhật cấu trúc đề, số lượng câu hỏi, phân bố nội dung và ma trận kiến thức của kỳ thi. \\
\hline

EDU-03 &
Đồng bộ chuẩn đầu ra và khung năng lực &
Tổ chức giáo dục / Đơn vị ra đề ĐGNL &
Cập nhật chuẩn đầu ra, nhóm năng lực và tiêu chí đánh giá được áp dụng cho kỳ thi ĐGNL. \\
\hline

EDU-04 &
Đồng bộ thống kê kết quả thi &
Tổ chức giáo dục / Đơn vị ra đề ĐGNL &
Cập nhật số liệu thống kê kết quả thi để phục vụ so sánh, phân tích và dự đoán. \\
\hline

EDU-05 &
Đồng bộ phân tích độ khó câu hỏi &
Tổ chức giáo dục / Đơn vị ra đề ĐGNL &
Cập nhật dữ liệu về độ khó, độ phân biệt và hiệu quả đánh giá của từng câu hỏi. \\
\hline

EDU-06 &
Đồng bộ phân bố điểm và năng lực &
Tổ chức giáo dục / Đơn vị ra đề ĐGNL &
Cập nhật phân bố điểm, mức năng lực và dữ liệu tham chiếu dùng trong phân tích kết quả. \\
\hline

\end{longtable}

\normalsize


% =========================================================
% 3. FUNCTIONAL REQUIREMENTS
% =========================================================
\section*{3. Functional Requirements}
\addcontentsline{toc}{section}{3. Functional Requirements}

\subsection*{3.1 System Functional Overview}
\addcontentsline{toc}{subsection}{3.1 System Functional Overview}


% =========================================================
% 3.1.1--3.1.10: DÙNG ARTIFACT GỐC TỪ GITHUB
% Không dùng TikZ, flowstep, fbox hoặc sơ đồ tự dựng.
% =========================================================
\newcommand{\GitHubFlowImage}[4]{%
  \subsubsection*{#1}%
  \addcontentsline{toc}{subsubsection}{#1}%
  \begin{figure}[H]
    \centering
    \GitHubIncludeGraphics[width=#3\textwidth,height=0.77\textheight,keepaspectratio]{figures/github/flow_#2.png}
    \caption{#4}
  \end{figure}
  \clearpage
}

\newcommand{\GitHubFlowPng}[3]{%
  \subsubsection*{#1}%
  \addcontentsline{toc}{subsubsection}{#1}%
  \begin{figure}[H]
    \centering
    \GitHubIncludeGraphics[width=0.96\textwidth,height=0.77\textheight,keepaspectratio]{#2}
    \caption{#3}
  \end{figure}
  \clearpage
}

\GitHubFlowPng{3.1.1 Academic Administrator Screen Flow}{figures/github/flow_admin_latest.png}{Screen Flow Quản trị viên Học vụ mới nhất từ GitHub.}
\GitHubFlowPng{3.1.2 Teacher Screen Flow}{figures/github/flow_teacher_latest.png}{Screen Flow Giáo viên / Giảng viên mới nhất từ GitHub.}
\GitHubFlowPng{3.1.3 Parent Screen Flow}{figures/github/flow_parent_latest.png}{Screen Flow Phụ huynh mới nhất từ GitHub.}
\GitHubFlowPng{3.1.4 Student Screen Flow}{figures/github/flow_student_latest.png}{Screen Flow Học sinh / Sinh viên mới nhất từ GitHub.}
\GitHubFlowPng{3.1.5 AI Tutor and Personalized Learning Roadmap Screen Flow}{figures/github/flow_ai_personalized_latest.png}{Screen Flow AI Tutor và Cá nhân hóa Lộ trình mới nhất từ GitHub.}

\subsubsection*{3.1.6 Screen Authorization}
\addcontentsline{toc}{subsubsection}{3.1.6 Screen Authorization}
\renewcommand{\arraystretch}{1.15}
\begin{longtable}{|p{4.7cm}|c|c|c|c|}
\hline
\rowcolor{orange!15}
\textbf{Screen / Module} & \textbf{Student} & \textbf{Teacher} & \textbf{Parent} & \textbf{Admin} \\
\hline
Login / Logout & X & X & X & X \\
\hline
Student Profile \& Goals & CRU & R & R* & R \\
\hline
AI Tutor & X & -- & -- & -- \\
\hline
Personalized Roadmap & R/X & R & R & R \\
\hline
Practice / Examination & X & R/A & R* & R \\
\hline
Results \& Competency Analytics & R & R* & R* & R \\
\hline
Class / Assignment / Grading & -- & CRUD/A & R* & R \\
\hline
Question Bank / Examination Management & -- & CRUD/A & -- & A \\
\hline
Parent Monitoring & -- & R* & R/X & -- \\
\hline
User / Role / System Configuration & -- & -- & -- & CRUD/A \\
\hline
Payment \& Invoice & X* & -- & X* & R* \\
\hline
Notification / Reporting & R & R* & R* & R/A \\
\hline
\end{longtable}
\noindent\textit{Ký hiệu: C--Create, R--Read, U--Update, D--Delete, A--Approve/Review, X--thao tác nghiệp vụ; dấu * biểu thị dữ liệu bị giới hạn theo phạm vi liên kết/lớp/tài khoản.}

\subsubsection*{3.1.7 Non-screen Functions}
\addcontentsline{toc}{subsubsection}{3.1.7 Non-screen Functions}
\begin{longtable}{|p{4.2cm}|p{9.2cm}|}
\hline
\rowcolor{orange!15}\textbf{System Function} & \textbf{Description} \\
\hline
LLM + RAG Processing & Xây dựng prompt, truy xuất Vector Database, kết hợp context và gọi LLM để sinh phản hồi AI Tutor. \\
\hline
Embedding / Vector Search & Tạo embedding cho tài liệu và truy vấn ngữ nghĩa phục vụ RAG. \\
\hline
Knowledge Tracing \& Prediction & Ước lượng mức thành thạo kiến thức, dự đoán kết quả thi và hỗ trợ điều chỉnh lộ trình. \\
\hline
Adaptive Practice Generation & Chọn/sinh câu hỏi theo chủ đề, độ khó và năng lực hiện tại của người học. \\
\hline
Payment Confirmation & Trao đổi với Payment Gateway, nhận trạng thái giao dịch và kích hoạt quyền lợi sau xác nhận thành công. \\
\hline
Notification Delivery & Lập lịch và gửi Email/SMS/Push, theo dõi trạng thái gửi và lỗi theo từng người nhận. \\
\hline
Education Organization Synchronization & Đồng bộ đề mẫu ĐGNL, ma trận đề, chuẩn đầu ra, khung năng lực, thống kê và phân tích độ khó. \\
\hline
Audit / Security Logging & Ghi nhận đăng nhập, thay đổi cấu hình AI và các sự kiện quản trị/bảo mật cần kiểm tra. \\
\hline
\end{longtable}

\subsubsection*{3.1.8 Entity Relationship Diagram}
\addcontentsline{toc}{subsubsection}{3.1.8 Entity Relationship Diagram}
\begin{figure}[htbp]
    \centering
    \includegraphics[width=0.99\textwidth,height=0.72\textheight,keepaspectratio]{figures/erd/erd_ai_edu_pathways.png}
    \caption{Entity Relationship Diagram cập nhật của hệ thống AI-Edu Pathways (trích từ báo cáo hiện tại; chưa tìm thấy artifact ERD riêng trên GitHub).}
    \label{fig:erd-ai-edu-pathways}
\end{figure}

Sơ đồ ERD thể hiện các nhóm dữ liệu chính: người dùng và phân quyền/liên kết phụ huynh; lớp học, buổi học, điểm danh và giao bài tập; kỳ thi, đề thi, câu hỏi, bài làm và chi tiết bài làm; lộ trình, tiến độ, bài học và tài liệu; AI Tutor và cấu hình AI; thanh toán, mã giảm giá, gói dịch vụ và hóa đơn điện tử; thông báo theo từng người nhận; hội thoại trao đổi giữa Phụ huynh và Giáo viên.

% =========================================================
% 3.2 WEB APPLICATION - source thật từ Library
% =========================================================
\setcounter{chapter}{3}
\setcounter{section}{1}
\setcounter{subsection}{0}
\setcounter{subsubsection}{0}
% =========================================================
% 3.2 WEB APPLICATION - FINAL VERSION
% =========================================================
% File này được \input từ main.tex.
% Không thêm \documentclass, \begin{document}, \end{document}.
% Screen layout sử dụng ảnh giao diện đã đóng gói; thao tác dùng chung màn hình
% chỉ dùng chung ảnh khi cùng xuất hiện trong đúng giao diện đó.
% Nội dung được đồng bộ theo Use Case, Screen Flow và ERD/logical additions hiện có.
% =========================================================

\section{Web Application}
\label{sec:web-application}

Phần này đặc tả chi tiết các chức năng Web Application theo cấu trúc thống nhất:
\textit{Function trigger, Function description, Actors/Roles, Purpose,
Data processing, Business rules, Screen layout}. Các xử lý service thuần
như LLM/RAG, Payment Service và Notification Service được mô tả trong
\textit{Data processing} của chức năng Web tương ứng.

% =========================================================
% AUTHENTICATE WEB USERS
% =========================================================

\subsection{Authenticate Web Users}
\label{subsec:authenticate-web-users}

Phân hệ xác thực kiểm soát quyền truy cập và điều hướng người dùng tới Dashboard phù hợp với vai trò.

\subsubsection{Log in to the web using email/phone number and password}
\label{subsubsec:web-login}

\begin{itemize}

    \item \textbf{Function trigger:}

    Khi người dùng truy cập Web Application mà chưa được xác thực, hệ thống hiển thị màn hình đăng nhập.

    Người dùng cũng được chuyển về màn hình đăng nhập khi phiên hiện tại hết hạn hoặc không còn hợp lệ.

    \item \textbf{Function description:}

    \begin{itemize}

        \item \textbf{Actors/Roles:}
        \begin{itemize}
            \item Học sinh / Sinh viên.
            \item Giáo viên / Giảng viên.
            \item Phụ huynh.
            \item Quản trị viên học vụ.
        \end{itemize}

        \item \textbf{Purpose:}
        \begin{enumerate}
            \item Xác minh danh tính trước khi truy cập chức năng được bảo vệ.
            \item Xác định vai trò sử dụng hệ thống và phạm vi quyền tương ứng.
            \item Kiểm tra trạng thái tài khoản trước khi tạo phiên đăng nhập.
            \item Điều hướng người dùng tới Dashboard tương ứng sau khi đăng nhập thành công.
        \end{enumerate}

        \item \textbf{Data processing:}
        \begin{enumerate}
            \item Hệ thống hiển thị các trường Email hoặc Số điện thoại, Mật khẩu và lựa chọn vai trò sử dụng hệ thống.
            \item Người dùng nhập thông tin đăng nhập và chọn vai trò Học sinh, Giáo viên, Phụ huynh hoặc Quản trị viên.
            \item Khi chọn ``Đăng nhập'', hệ thống kiểm tra các trường bắt buộc; nếu thiếu dữ liệu thì giữ nguyên màn hình và hiển thị lỗi tại trường tương ứng.
            \item Hệ thống tìm \texttt{NguoiDung} theo \texttt{Email} hoặc \texttt{SoDienThoai}; nếu không tồn tại tài khoản thì từ chối đăng nhập.
            \item Nếu tìm thấy tài khoản, hệ thống xác minh mật khẩu bằng cơ chế bảo mật phù hợp và không xử lý mật khẩu dưới dạng văn bản thuần.
            \item Hệ thống kiểm tra \texttt{TrangThai}; tài khoản bị khóa hoặc không hoạt động không được tạo phiên đăng nhập.
            \item Hệ thống kiểm tra vai trò đã chọn có phù hợp với \texttt{VaiTro} của tài khoản; nếu không phù hợp thì từ chối truy cập theo vai trò đó.
            \item Khi hợp lệ, hệ thống tạo phiên đăng nhập, lấy \texttt{MaNguoiDung}, \texttt{VaiTro} và áp dụng Role Authorization.
            \item Hệ thống điều hướng tới Student Dashboard, Teacher Dashboard, Parent Dashboard hoặc Admin Dashboard theo vai trò.
            \item Nếu thông tin đăng nhập sai, hệ thống hiển thị thông báo lỗi chung và cho phép nhập lại mà không tiết lộ thông tin nhạy cảm.
        \end{enumerate}

    \end{itemize}

    \item \textbf{Business rules:}

    BR-01, BR-02, BR-03, BR-04, BR-05, BR-07.

    \ScreenLayoutItem

    \begin{figure}[H]
        \centering
        \SafeIncludeGraphics[width=0.94\textwidth,height=0.78\textheight,keepaspectratio]{figures/03_02/login.png}
        \caption{Màn hình đăng nhập hệ thống AI-Edu Pathways.}
        \label{fig:screen-login}
    \end{figure}

\end{itemize}

\subsubsection{Log out of the web application}
\label{subsubsec:web-logout}

\begin{itemize}

    \item \textbf{Function trigger:}

    Người dùng đã đăng nhập chọn ``Đăng xuất'' từ menu tài khoản hoặc Dashboard.

    \item \textbf{Function description:}

    \begin{itemize}

        \item \textbf{Actors/Roles:}
        \begin{itemize}
            \item Học sinh / Sinh viên.
            \item Giáo viên / Giảng viên.
            \item Phụ huynh.
            \item Quản trị viên học vụ.
        \end{itemize}

        \item \textbf{Purpose:}
        \begin{enumerate}
            \item Kết thúc phiên làm việc hiện tại một cách an toàn.
            \item Ngăn phiên cũ tiếp tục được dùng để truy cập chức năng bảo vệ.
        \end{enumerate}

        \item \textbf{Data processing:}
        \begin{enumerate}
            \item Hệ thống nhận yêu cầu đăng xuất từ người dùng đã xác thực.
            \item Hệ thống xác định và vô hiệu hóa phiên đăng nhập hiện tại theo cơ chế xác thực được triển khai.
            \item Dữ liệu quyền tạm của phiên không còn được sử dụng cho các yêu cầu sau.
            \item Hệ thống chuyển người dùng về màn hình đăng nhập.
            \item Nếu người dùng truy cập lại chức năng yêu cầu xác thực, hệ thống yêu cầu đăng nhập lại.
        \end{enumerate}

    \end{itemize}

    \item \textbf{Business rules:}

    BR-01, BR-04, BR-07.

    \ScreenLayoutItem

    \ScreenLayoutFigure{figures/03_02/student_dashboard.png}{Log out of the web application}

\end{itemize}


% =========================================================
% MANAGE STUDENT PROFILE & LEARNING GOALS
% =========================================================

\subsection{Manage Student Profile \& Learning Goals}
\label{subsec:student-profile-learning-goals}

Phân hệ cho phép người học quản lý hồ sơ của chính mình và thiết lập mục tiêu làm đầu vào cho cá nhân hóa.

\subsubsection{View Student Dashboard}
\label{subsubsec:view-student-dashboard}

\begin{itemize}

    \item \textbf{Function trigger:}

    Sau khi Học sinh/Sinh viên đăng nhập thành công, hệ thống điều hướng tới Student Dashboard; người học cũng có thể quay về Dashboard từ thanh điều hướng chính.

    \item \textbf{Function description:}

    \begin{itemize}

        \item \textbf{Actors/Roles:}
        \begin{itemize}
            \item Học sinh / Sinh viên.
        \end{itemize}

        \item \textbf{Purpose:}
        \begin{enumerate}
            \item Cung cấp điểm truy cập tập trung tới các chức năng học tập chính của người học.
            \item Hiển thị thông tin tổng quan về lộ trình, tiến độ, nhiệm vụ và thông báo liên quan.
        \end{enumerate}

        \item \textbf{Data processing:}
        \begin{enumerate}
            \item Hệ thống xác minh phiên đăng nhập và vai trò Học sinh/Sinh viên.
            \item Hệ thống tải dữ liệu tổng quan thuộc chính người học, gồm lộ trình đang hoạt động, tiến độ gần nhất, nhiệm vụ/bài được giao, kết quả gần đây và thông báo chưa đọc theo dữ liệu hiện có.
            \item Dashboard cung cấp lối truy cập tới Lộ trình học cá nhân hóa, Tài liệu học tập, Luyện tập, Thi ĐGNL, AI Tutor, Kết quả/Phân tích năng lực, Thông báo và Cài đặt cá nhân theo Screen Flow Học sinh.
            \item Các số liệu hiển thị phải được giới hạn theo \texttt{MaNguoiDung} của phiên hiện tại; người học không được xem Dashboard dữ liệu của tài khoản khác.
            \item Nếu một chức năng chưa có dữ liệu, hệ thống hiển thị trạng thái rỗng phù hợp thay vì tự tạo dữ liệu.
            \item Khi người học chọn một lối tắt, hệ thống tiếp tục kiểm tra quyền trước khi điều hướng tới chức năng tương ứng.
        \end{enumerate}

    \end{itemize}

    \item \textbf{Business rules:}

    BR-04, BR-05, BR-07, BR-24.

    \ScreenLayoutItem

    \begin{figure}[H]
        \centering
        \SafeIncludeGraphics[width=0.94\textwidth,height=0.78\textheight,keepaspectratio]{figures/03_02/student_dashboard.png}
        \caption{Màn hình Dashboard tổng quan của Học sinh/Sinh viên.}
        \label{fig:screen-student-dashboard}
    \end{figure}

\end{itemize}


\subsubsection{Register Student Account}
\label{subsubsec:register-student-account}

\begin{itemize}

    \item \textbf{Function trigger:}

    Người dùng chưa có tài khoản chọn ``Đăng ký'' tại màn hình xác thực.

    \item \textbf{Function description:}

    \begin{itemize}

        \item \textbf{Actors/Roles:}
        \begin{itemize}
            \item Học sinh / Sinh viên.
            \item Hệ thống.
        \end{itemize}

        \item \textbf{Purpose:}
        \begin{enumerate}
            \item Tạo tài khoản người học để sử dụng hệ thống.
            \item Ghi nhận thông tin cá nhân ban đầu và khởi tạo vai trò Học sinh/Sinh viên.
        \end{enumerate}

        \item \textbf{Data processing:}
        \begin{enumerate}
            \item Hệ thống hiển thị biểu mẫu đăng ký với các trường hồ sơ phù hợp.
            \item Người học nhập tối thiểu họ tên, Email, số điện thoại và thông tin xác thực theo yêu cầu triển khai.
            \item Hệ thống kiểm tra trường bắt buộc, định dạng Email/Số điện thoại và tính duy nhất.
            \item Nếu dữ liệu không hợp lệ hoặc trùng với tài khoản đã có, hệ thống không tạo tài khoản và hiển thị lỗi.
            \item Nếu hợp lệ, hệ thống tạo \texttt{NguoiDung} với \texttt{VaiTro = HocSinh}, trạng thái mặc định và \texttt{NgayTao}.
            \item Các thông tin trường học/hồ sơ học tập của STU-01 được ghi nhận theo phạm vi dữ liệu hồ sơ mà hệ thống triển khai.
            \item Hệ thống thông báo đăng ký thành công và đưa người dùng tới bước đăng nhập hoặc hoàn thiện hồ sơ.
        \end{enumerate}

    \end{itemize}

    \item \textbf{Business rules:}

    BR-04, BR-08.

    \ScreenLayoutItem

    \ScreenLayoutFigure{figures/03_02/new_2026_09_10/register_student.png}{Register Student Account}

\end{itemize}

\subsubsection{View Student Profile}
\label{subsubsec:view-student-profile}

\begin{itemize}

    \item \textbf{Function trigger:}

    Học sinh/Sinh viên chọn ``Hồ sơ cá nhân'' hoặc ``Cài đặt cá nhân'' từ Student Dashboard.

    \item \textbf{Function description:}

    \begin{itemize}

        \item \textbf{Actors/Roles:}
        \begin{itemize}
            \item Học sinh / Sinh viên.
        \end{itemize}

        \item \textbf{Purpose:}
        \begin{enumerate}
            \item Hiển thị thông tin tài khoản và hồ sơ của người học hiện tại.
            \item Cung cấp điểm truy cập tới chức năng cập nhật hồ sơ.
        \end{enumerate}

        \item \textbf{Data processing:}
        \begin{enumerate}
            \item Hệ thống lấy \texttt{MaNguoiDung} từ phiên đăng nhập, không tin vào định danh do client tự nhập.
            \item Hệ thống truy xuất hồ sơ của chính người dùng đang đăng nhập.
            \item Hệ thống hiển thị họ tên, Email, số điện thoại, vai trò, trạng thái và các dữ liệu hồ sơ học tập đã triển khai.
            \item Các trường quản trị như vai trò và trạng thái tài khoản được hiển thị ở chế độ chỉ đọc đối với Học sinh/Sinh viên.
            \item Nếu phiên hết hạn hoặc hồ sơ không thuộc tài khoản hiện tại, hệ thống từ chối truy cập.
        \end{enumerate}

    \end{itemize}

    \item \textbf{Business rules:}

    BR-07, BR-08.

    \ScreenLayoutItem

    \ScreenLayoutFigure{figures/03_02/student_goals.png}{View Student Profile}

\end{itemize}

\subsubsection{Update Student Profile}
\label{subsubsec:update-student-profile}

\begin{itemize}

    \item \textbf{Function trigger:}

    Học sinh/Sinh viên chọn ``Chỉnh sửa hồ sơ'' từ màn hình Hồ sơ cá nhân.

    \item \textbf{Function description:}

    \begin{itemize}

        \item \textbf{Actors/Roles:}
        \begin{itemize}
            \item Học sinh / Sinh viên.
        \end{itemize}

        \item \textbf{Purpose:}
        \begin{enumerate}
            \item Cho phép người học cập nhật dữ liệu cá nhân và hồ sơ thuộc quyền của mình.
            \item Duy trì dữ liệu chính xác cho các chức năng học tập và thông báo.
        \end{enumerate}

        \item \textbf{Data processing:}
        \begin{enumerate}
            \item Hệ thống tải hồ sơ hiện tại và chỉ hiển thị các trường được phép chỉnh sửa.
            \item Người học thay đổi dữ liệu và chọn ``Lưu''.
            \item Hệ thống kiểm tra trường bắt buộc, định dạng và tính duy nhất của Email/Số điện thoại nếu thay đổi.
            \item Hệ thống xác minh bản ghi được cập nhật thuộc chính người dùng đang đăng nhập.
            \item Học sinh/Sinh viên không được tự thay đổi \texttt{VaiTro}, quyền truy cập hoặc trạng thái khóa/mở.
            \item Nếu hợp lệ, hệ thống cập nhật hồ sơ; nếu không hợp lệ, dữ liệu cũ được giữ nguyên và lỗi được hiển thị.
            \item Sau khi lưu thành công, hệ thống tải lại dữ liệu mới nhất.
        \end{enumerate}

    \end{itemize}

    \item \textbf{Business rules:}

    BR-07, BR-08.

    \ScreenLayoutItem

    \ScreenLayoutFigure{figures/03_02/student_goals.png}{Update Student Profile}

\end{itemize}

\subsubsection{Set or Update Learning Goal and Target Examination}
\label{subsubsec:set-update-learning-goal}

\begin{itemize}

    \item \textbf{Function trigger:}

    Học sinh/Sinh viên chọn ``Thiết lập mục tiêu'' lần đầu hoặc ``Điều chỉnh mục tiêu'' từ khu vực lộ trình.

    \item \textbf{Function description:}

    \begin{itemize}

        \item \textbf{Actors/Roles:}
        \begin{itemize}
            \item Học sinh / Sinh viên.
            \item Hệ thống AI.
        \end{itemize}

        \item \textbf{Purpose:}
        \begin{enumerate}
            \item Ghi nhận điểm mục tiêu, kỳ thi đích, thời gian chuẩn bị và thời lượng học dự kiến.
            \item Cung cấp đầu vào cho tạo hoặc điều chỉnh lộ trình cá nhân hóa.
        \end{enumerate}

        \item \textbf{Data processing:}
        \begin{enumerate}
            \item Hệ thống tải danh sách kỳ thi mục tiêu từ dữ liệu \texttt{KyThi}.
            \item Người học nhập hoặc chỉnh sửa điểm mục tiêu, kỳ thi đích, thời gian chuẩn bị và thời lượng học dự kiến.
            \item Hệ thống kiểm tra kỳ thi hợp lệ và các giá trị mục tiêu thuộc miền cho phép.
            \item Nếu dữ liệu không hợp lệ, hệ thống không lưu và yêu cầu điều chỉnh.
            \item Nếu hợp lệ, hệ thống lưu mục tiêu gắn với người học.
            \item Hệ thống kết hợp mục tiêu với năng lực hiện tại, kết quả đánh giá đầu vào và tiến độ để tạo/điều chỉnh Roadmap.
            \item Nếu mục tiêu thay đổi khi đã có lộ trình, hệ thống đánh giá lại lộ trình nhưng không xóa lịch sử tiến độ đã hoàn thành.
            \item Hệ thống hiển thị tóm tắt mục tiêu mới.
        \end{enumerate}

    \end{itemize}

    \item \textbf{Business rules:}

    BR-17, BR-18, BR-19, BR-22, BR-23.

    \ScreenLayoutItem

    \begin{figure}[H]
        \centering
        \SafeIncludeGraphics[width=0.94\textwidth,height=0.78\textheight,keepaspectratio]{figures/03_02/student_goals.png}
        \caption{Màn hình Mục tiêu và Thử thách của Học sinh/Sinh viên.}
        \label{fig:screen-student-goals}
    \end{figure}

\end{itemize}


% =========================================================
% AI TUTOR
% =========================================================

\subsection{AI Tutor}
\label{subsec:ai-tutor}

Phân hệ AI Tutor hỗ trợ người học đặt câu hỏi, truy xuất ngữ cảnh bằng RAG, nhận phản hồi từ LLM và lưu lịch sử hội thoại.

\subsubsection{Start or Open an AI Tutor Conversation}
\label{subsubsec:start-ai-tutor-conversation}

\begin{itemize}

    \item \textbf{Function trigger:}

    Học sinh/Sinh viên chọn ``AI Tutor'' từ Student Dashboard hoặc từ ngữ cảnh học tập liên quan.

    \item \textbf{Function description:}

    \begin{itemize}

        \item \textbf{Actors/Roles:}
        \begin{itemize}
            \item Học sinh / Sinh viên.
        \end{itemize}

        \item \textbf{Purpose:}
        \begin{enumerate}
            \item Mở phiên AI Tutor để đặt câu hỏi.
            \item Cho phép tiếp tục hội thoại cũ hoặc tạo hội thoại mới.
        \end{enumerate}

        \item \textbf{Data processing:}
        \begin{enumerate}
            \item Hệ thống xác minh người dùng có vai trò Học sinh/Sinh viên.
            \item Hệ thống tải danh sách \texttt{HoiThoaiAITutor} thuộc người học nếu mở lịch sử.
            \item Nếu tạo mới, hệ thống khởi tạo \texttt{HoiThoaiAITutor} với thông tin trạng thái phù hợp.
            \item Nếu mở hội thoại cũ, hệ thống kiểm tra hội thoại thuộc chính người học.
            \item Hệ thống tải \texttt{TinNhanAITutor} theo thứ tự thời gian để khôi phục bối cảnh.
            \item Hệ thống không cho phép truy cập hội thoại AI Tutor của tài khoản khác.
        \end{enumerate}

    \end{itemize}

    \item \textbf{Business rules:}

    BR-07, BR-26, BR-28, BR-29, BR-32.

    \ScreenLayoutItem

    \begin{figure}[H]
        \centering
        \SafeIncludeGraphics[width=0.94\textwidth,height=0.78\textheight,keepaspectratio]{figures/03_02/ai_tutor.png}
        \caption{Màn hình AI Tutor với hội thoại, kiến thức liên quan và tài liệu đề xuất.}
        \label{fig:screen-ai-tutor}
    \end{figure}

\end{itemize}

\subsubsection{Ask AI Tutor and Receive an LLM + RAG Response}
\label{subsubsec:ask-ai-tutor}

\begin{itemize}

    \item \textbf{Function trigger:}

    Người học nhập câu hỏi trong hội thoại AI Tutor và chọn gửi.

    \item \textbf{Function description:}

    \begin{itemize}

        \item \textbf{Actors/Roles:}
        \begin{itemize}
            \item Học sinh / Sinh viên.
            \item AI Service.
        \end{itemize}

        \item \textbf{Purpose:}
        \begin{enumerate}
            \item Cung cấp lời giải thích có ngữ cảnh cho câu hỏi học tập.
            \item Kết hợp ngữ cảnh học tập với tài liệu RAG trước khi LLM sinh phản hồi.
        \end{enumerate}

        \item \textbf{Data processing:}
        \begin{enumerate}
            \item Hệ thống kiểm tra câu hỏi không rỗng và hội thoại thuộc người học.
            \item Hệ thống lưu câu hỏi như tin nhắn do User gửi trong \texttt{TinNhanAITutor}.
            \item Hệ thống tạo prompt từ câu hỏi và ngữ cảnh phù hợp như lịch sử hội thoại, môn/chủ đề, cấp độ hoặc dữ liệu học tập liên quan.
            \item Hệ thống thực hiện RAG trên kho tài liệu/Vector Database để tìm ngữ cảnh liên quan trước khi gọi LLM.
            \item Ngữ cảnh truy xuất được kết hợp với prompt theo cấu hình AI đang áp dụng.
            \item AI Service sinh phản hồi dựa trên prompt và ngữ cảnh.
            \item Hệ thống lưu phản hồi AI vào \texttt{TinNhanAITutor}, phân biệt \texttt{NguoiGui = AI}.
            \item Hệ thống hiển thị câu trả lời và duy trì lịch sử.
            \item Nếu RAG hoặc LLM thất bại, hệ thống không ghi phản hồi giả và cho phép người học thử lại.
        \end{enumerate}

    \end{itemize}

    \item \textbf{Business rules:}

    BR-26, BR-27, BR-29, BR-30, BR-32, BR-33.

    \ScreenLayoutItem

    \ScreenLayoutFigure{figures/03_02/ai_tutor.png}{Ask AI Tutor and Receive an LLM + RAG Response}

\end{itemize}

\subsubsection{Rate AI Tutor Response and Send Feedback}
\label{subsubsec:rate-ai-tutor-response}

\begin{itemize}

    \item \textbf{Function trigger:}

    Sau khi nhận câu trả lời, Học sinh/Sinh viên đánh giá mức độ hài lòng hoặc gửi phản hồi.

    \item \textbf{Function description:}

    \begin{itemize}

        \item \textbf{Actors/Roles:}
        \begin{itemize}
            \item Học sinh / Sinh viên.
        \end{itemize}

        \item \textbf{Purpose:}
        \begin{enumerate}
            \item Thu thập đánh giá về chất lượng câu trả lời.
            \item Cho phép tiếp tục hỏi lại khi phản hồi chưa đáp ứng nhu cầu.
        \end{enumerate}

        \item \textbf{Data processing:}
        \begin{enumerate}
            \item Hệ thống xác định phản hồi AI đang được đánh giá thuộc hội thoại của người học.
            \item Người học chọn hài lòng/không hài lòng hoặc nhập nội dung phản hồi theo chức năng triển khai.
            \item Hệ thống ghi nhận phản hồi gắn với nội dung được đánh giá.
            \item Nếu không hài lòng, hệ thống giữ hội thoại hiện tại để người học đặt lại câu hỏi hoặc làm rõ.
            \item Nếu hài lòng, hội thoại vẫn được lưu để tham khảo.
            \item Dữ liệu phản hồi có thể phục vụ kiểm soát chất lượng nhưng không tự động sửa nội dung học thuật đã lưu mà không qua cơ chế kiểm soát.
        \end{enumerate}

    \end{itemize}

    \item \textbf{Business rules:}

    BR-31, BR-32, BR-44.

    \ScreenLayoutItem

    \ScreenLayoutFigure{figures/03_02/ai_tutor.png}{Rate AI Tutor Response and Send Feedback}

\end{itemize}


% =========================================================
% PERSONALIZED LEARNING ROADMAP
% =========================================================

\subsection{Personalized Learning Roadmap}
\label{subsec:personalized-learning-roadmap}

Phân hệ lộ trình cá nhân hóa dùng kết quả đánh giá năng lực và mục tiêu học tập để xây dựng, theo dõi và điều chỉnh Roadmap.

\subsubsection{Complete Entry Assessment and Generate Initial Roadmap}
\label{subsubsec:generate-initial-roadmap}

\begin{itemize}

    \item \textbf{Function trigger:}

    Người học bắt đầu xây dựng lộ trình mới hoặc chưa có dữ liệu năng lực đầu vào phù hợp.

    \item \textbf{Function description:}

    \begin{itemize}

        \item \textbf{Actors/Roles:}
        \begin{itemize}
            \item Học sinh / Sinh viên.
            \item Hệ thống AI.
        \end{itemize}

        \item \textbf{Purpose:}
        \begin{enumerate}
            \item Đánh giá năng lực hiện tại.
            \item Xác định điểm mạnh, điểm yếu và khoảng trống kiến thức.
            \item Tạo lộ trình ban đầu phù hợp với mục tiêu.
        \end{enumerate}

        \item \textbf{Data processing:}
        \begin{enumerate}
            \item Hệ thống yêu cầu người học thực hiện bài đánh giá năng lực đầu vào theo luồng thi/đánh giá.
            \item Sau khi bài được nộp và chấm, hệ thống tổng hợp điểm theo nhóm kiến thức.
            \item AI phân tích kết quả để xác định điểm mạnh, điểm yếu và khoảng trống kiến thức.
            \item Hệ thống kết hợp kết quả với mục tiêu/kỳ thi đích.
            \item Hệ thống tạo \texttt{LoTrinhHocTap} thuộc người học.
            \item Các \texttt{BaiHoc} được đưa vào lộ trình qua \texttt{ChiTietLoTrinh} với thứ tự và trạng thái hoàn thành.
            \item Hệ thống hiển thị Roadmap ban đầu.
            \item Nếu thiếu dữ liệu mục tiêu cần thiết, hệ thống yêu cầu thiết lập mục tiêu trước khi tạo lộ trình hoàn chỉnh.
        \end{enumerate}

    \end{itemize}

    \item \textbf{Business rules:}

    BR-18, BR-19, BR-20, BR-21.

    \ScreenLayoutItem

    \ScreenLayoutFigure{figures/03_02/exam.png}{Complete Entry Assessment and Generate Initial Roadmap}

    \ScreenLayoutFigure{figures/03_02/student_learning_path.png}{Complete Entry Assessment and Generate Initial Roadmap}

\end{itemize}

\subsubsection{View Personalized Learning Roadmap}
\label{subsubsec:view-personalized-roadmap}

\begin{itemize}

    \item \textbf{Function trigger:}

    Học sinh/Sinh viên chọn ``Lộ trình học cá nhân'' từ Student Dashboard.

    \item \textbf{Function description:}

    \begin{itemize}

        \item \textbf{Actors/Roles:}
        \begin{itemize}
            \item Học sinh / Sinh viên.
        \end{itemize}

        \item \textbf{Purpose:}
        \begin{enumerate}
            \item Hiển thị lộ trình hiện tại theo thứ tự ưu tiên.
            \item Cho phép truy cập bài học, luyện tập và tiến độ tương ứng.
        \end{enumerate}

        \item \textbf{Data processing:}
        \begin{enumerate}
            \item Hệ thống lấy các \texttt{LoTrinhHocTap} thuộc người học.
            \item Hệ thống xác định lộ trình đang hoạt động và tải \texttt{ChiTietLoTrinh} theo thứ tự học.
            \item Hệ thống kết hợp \texttt{BaiHoc}, tiến độ và tài liệu liên quan để trình bày Roadmap.
            \item Mỗi mục thể hiện tối thiểu trạng thái hoàn thành và thứ tự học; hạn hoàn thành hiển thị nếu có.
            \item Người học có thể chọn bài học hoặc hoạt động để tiếp tục.
            \item Hệ thống không cho phép truy cập lộ trình của tài khoản khác.
        \end{enumerate}

    \end{itemize}

    \item \textbf{Business rules:}

    BR-20, BR-21, BR-24, BR-25.

    \ScreenLayoutItem

    \begin{figure}[H]
        \centering
        \SafeIncludeGraphics[width=0.94\textwidth,height=0.78\textheight,keepaspectratio]{figures/03_02/student_learning_path.png}
        \caption{Màn hình Lộ trình học tập cá nhân hóa.}
        \label{fig:screen-learning-path}
    \end{figure}

\end{itemize}

\subsubsection{Track Roadmap Progress}
\label{subsubsec:track-roadmap-progress}

\begin{itemize}

    \item \textbf{Function trigger:}

    Người học hoàn thành bài học/luyện tập hoặc mở khu vực theo dõi tiến độ.

    \item \textbf{Function description:}

    \begin{itemize}

        \item \textbf{Actors/Roles:}
        \begin{itemize}
            \item Học sinh / Sinh viên.
            \item Hệ thống.
        \end{itemize}

        \item \textbf{Purpose:}
        \begin{enumerate}
            \item Cập nhật và hiển thị mức độ hoàn thành lộ trình.
            \item Theo dõi thời gian học và trạng thái từng bài học.
        \end{enumerate}

        \item \textbf{Data processing:}
        \begin{enumerate}
            \item Hệ thống xác định bài học và người học đang thực hiện.
            \item Hệ thống cập nhật \texttt{TienDoHocTap} với mức hoàn thành, thời gian học tích lũy và thời điểm cập nhật.
            \item Hệ thống cập nhật trạng thái \texttt{ChiTietLoTrinh} khi điều kiện hoàn thành bài học được đáp ứng.
            \item Hệ thống tính lại mức độ hoàn thành của lộ trình.
            \item Nếu giai đoạn hoàn thành, hệ thống hiển thị báo cáo tiến độ và gợi ý bước tiếp theo.
            \item Nếu chưa đạt mục tiêu, dữ liệu tiến độ được đưa vào quá trình điều chỉnh Roadmap.
        \end{enumerate}

    \end{itemize}

    \item \textbf{Business rules:}

    BR-21, BR-23, BR-24.

    \ScreenLayoutItem

    \begin{figure}[H]
        \centering
        \SafeIncludeGraphics[width=0.94\textwidth,height=0.78\textheight,keepaspectratio]{figures/03_02/student_roadmap_progress.png}
        \caption{Màn hình theo dõi và cập nhật tiến độ lộ trình học tập.}
        \label{fig:screen-roadmap-progress}
    \end{figure}

\end{itemize}

\subsubsection{Request or Apply Roadmap Adjustment}
\label{subsubsec:adjust-roadmap}

\begin{itemize}

    \item \textbf{Function trigger:}

    Người học chọn ``Yêu cầu điều chỉnh lộ trình'' hoặc hệ thống phát hiện chưa đạt mục tiêu của giai đoạn hiện tại.

    \item \textbf{Function description:}

    \begin{itemize}

        \item \textbf{Actors/Roles:}
        \begin{itemize}
            \item Học sinh / Sinh viên.
            \item Hệ thống AI.
        \end{itemize}

        \item \textbf{Purpose:}
        \begin{enumerate}
            \item Điều chỉnh lộ trình khi mục tiêu, thời gian học hoặc năng lực thay đổi.
            \item Duy trì lộ trình phù hợp với kết quả thực tế.
        \end{enumerate}

        \item \textbf{Data processing:}
        \begin{enumerate}
            \item Hệ thống thu thập mục tiêu hiện tại, tiến độ, kết quả luyện tập/thi và trạng thái bài học.
            \item Nếu do người học yêu cầu, hệ thống ghi nhận thay đổi mục tiêu hoặc lý do liên quan.
            \item AI đánh giá phần kiến thức chưa đạt, mức độ khó hiện tại và thời gian còn lại.
            \item Hệ thống có thể thay đổi độ khó, thứ tự chủ đề hoặc ưu tiên nội dung.
            \item Hệ thống áp dụng thay đổi vào lộ trình theo cơ chế dữ liệu triển khai.
            \item Hệ thống giữ lịch sử kết quả và tiến độ đã phát sinh.
            \item Roadmap cập nhật được hiển thị lại cho người học.
        \end{enumerate}

    \end{itemize}

    \item \textbf{Business rules:}

    BR-18, BR-22, BR-23, BR-24.

    \ScreenLayoutItem

    \ScreenLayoutFigure{figures/03_02/student_learning_path.png}{Request or Apply Roadmap Adjustment}

\end{itemize}

% =========================================================
% ADAPTIVE PRACTICE & EXAMINATION
% =========================================================

\subsection{Adaptive Practice \& Examination}
\label{subsec:adaptive-practice-examination}

Phân hệ luyện tập và thi cho phép người học chọn loại đề, làm bài trong thời gian quy định, nộp bài và chuyển dữ liệu sang quy trình chấm/đánh giá.

\subsubsection{Generate an Adaptive Practice Set}
\label{subsubsec:generate-adaptive-practice-set}

\begin{itemize}

    \item \textbf{Function trigger:}

    Học sinh/Sinh viên chọn bắt đầu luyện tập thích ứng hoặc hệ thống đề xuất một phiên luyện tập dựa trên năng lực hiện tại.

    \item \textbf{Function description:}

    \begin{itemize}

        \item \textbf{Actors/Roles:}
        \begin{itemize}
            \item Học sinh / Sinh viên.
            \item Hệ thống AI.
        \end{itemize}

        \item \textbf{Purpose:}
        \begin{enumerate}
            \item Tạo bài luyện tập có nội dung và độ khó phù hợp với năng lực hiện tại theo STU-07.
            \item Ưu tiên các chủ đề còn yếu và điều chỉnh bài luyện khi năng lực thay đổi.
        \end{enumerate}

        \item \textbf{Data processing:}
        \begin{enumerate}
            \item Hệ thống xác định người học, mục tiêu hiện tại và phạm vi môn/chủ đề luyện tập.
            \item Hệ thống tổng hợp dữ liệu năng lực từ kết quả bài làm, Knowledge Tracing, tiến độ học tập và các chủ đề còn yếu.
            \item Hệ thống xác định mức độ khó và phân bố chủ đề phù hợp với năng lực hiện tại.
            \item Hệ thống lựa chọn câu hỏi từ \texttt{CauHoi} theo ma trận kiến thức/độ khó hoặc yêu cầu AI Service sinh câu hỏi khi chức năng AI-03/AI-04 được sử dụng.
            \item Các câu hỏi sinh bởi AI phải đi qua cơ chế kiểm soát chất lượng trước khi được dùng trong bài luyện nếu chính sách hệ thống yêu cầu.
            \item Hệ thống tạo/cấu hình một \texttt{DeThi} hoặc phiên luyện tập tương ứng theo cơ chế triển khai và liên kết các câu hỏi qua \texttt{DeThi\_CauHoi}.
            \item Sau mỗi phiên luyện tập, kết quả mới được sử dụng để cập nhật năng lực và làm đầu vào cho lần tạo bài thích ứng tiếp theo.
            \item Nếu dữ liệu năng lực chưa đủ, hệ thống sử dụng mức khởi đầu/đánh giá đầu vào thay vì suy đoán không có căn cứ.
        \end{enumerate}

    \end{itemize}

    \item \textbf{Business rules:}

    BR-35, BR-36, BR-42, BR-44.

    \ScreenLayoutItem

    \ScreenLayoutFigure{figures/03_02/new_2026_09_10/adaptive_practice_list.png}{Generate an Adaptive Practice Set}

\end{itemize}


\subsubsection{View Available Practice and Examination List}
\label{subsubsec:view-practice-exam-list}

\begin{itemize}

    \item \textbf{Function trigger:}

    Học sinh/Sinh viên chọn ``Luyện tập'' hoặc ``Thi ĐGNL'' từ Student Dashboard.

    \item \textbf{Function description:}

    \begin{itemize}

        \item \textbf{Actors/Roles:}
        \begin{itemize}
            \item Học sinh / Sinh viên.
        \end{itemize}

        \item \textbf{Purpose:}
        \begin{enumerate}
            \item Hiển thị các bài luyện tập, thi thử và kỳ thi phù hợp.
            \item Phân biệt rõ loại kỳ thi trước khi bắt đầu.
        \end{enumerate}

        \item \textbf{Data processing:}
        \begin{enumerate}
            \item Hệ thống tải các \texttt{KyThi}, \texttt{DeThi} và các phiên/bài luyện thích ứng mà người học được phép truy cập.
            \item Hệ thống phân loại nội dung theo loại như Thi thử, Thi chính thức hoặc Đánh giá đầu vào.
            \item Nếu đề được giao qua lớp, hệ thống kiểm tra người học thuộc lớp của \texttt{GiaoBaiTap} và assignment còn hiệu lực.
            \item Hệ thống hiển thị tên đề, loại đề, thời gian làm bài và trạng thái khả dụng.
            \item Người học chọn một đề để xem thông tin và bắt đầu.
            \item Đề ngoài phạm vi hoặc chưa tới thời điểm cho phép không được mở.
        \end{enumerate}

    \end{itemize}

    \item \textbf{Business rules:}

    BR-35, BR-38, BR-63.

    \ScreenLayoutItem

    \ScreenLayoutFigure{figures/03_02/new_2026_09_10/adaptive_practice_list.png}{View Available Practice and Examination List}

\end{itemize}

\subsubsection{Start a Practice or Examination Attempt}
\label{subsubsec:start-exam-attempt}

\begin{itemize}

    \item \textbf{Function trigger:}

    Người học chọn ``Bắt đầu'' trên một đề luyện tập/đề thi hợp lệ.

    \item \textbf{Function description:}

    \begin{itemize}

        \item \textbf{Actors/Roles:}
        \begin{itemize}
            \item Học sinh / Sinh viên.
        \end{itemize}

        \item \textbf{Purpose:}
        \begin{enumerate}
            \item Khởi tạo một lần làm bài hợp lệ.
            \item Ghi nhận thời điểm bắt đầu và gắn bài làm với người học/đề thi.
        \end{enumerate}

        \item \textbf{Data processing:}
        \begin{enumerate}
            \item Hệ thống kiểm tra quyền truy cập đề và điều kiện thời gian trước khi bắt đầu.
            \item Hệ thống lấy cấu trúc đề từ \texttt{DeThi} và danh sách câu hỏi qua \texttt{DeThi\_CauHoi}.
            \item Hệ thống tạo \texttt{BaiLam} với \texttt{MaHocSinh}, \texttt{MaDeThi}, thời gian bắt đầu và \texttt{MaGiaoBaiTap} nếu bài bắt nguồn từ assignment.
            \item Hệ thống tải câu hỏi theo thứ tự được cấu hình.
            \item Bộ đếm thời gian được khởi tạo theo thời gian làm bài của đề.
            \item Nếu không đủ điều kiện bắt đầu, hệ thống không tạo bài làm mới.
        \end{enumerate}

    \end{itemize}

    \item \textbf{Business rules:}

    BR-36, BR-37, BR-38, BR-39, BR-63.

    \ScreenLayoutItem

    \ScreenLayoutFigure{figures/03_02/exam.png}{Start a Practice or Examination Attempt}

\end{itemize}

\subsubsection{Answer Examination Questions}
\label{subsubsec:answer-exam-questions}

\begin{itemize}

    \item \textbf{Function trigger:}

    Người học đang trong một \texttt{BaiLam} hợp lệ và nhập/chọn câu trả lời.

    \item \textbf{Function description:}

    \begin{itemize}

        \item \textbf{Actors/Roles:}
        \begin{itemize}
            \item Học sinh / Sinh viên.
        \end{itemize}

        \item \textbf{Purpose:}
        \begin{enumerate}
            \item Ghi nhận câu trả lời cho từng câu hỏi.
            \item Theo dõi dữ liệu phục vụ chấm điểm và phân tích năng lực.
        \end{enumerate}

        \item \textbf{Data processing:}
        \begin{enumerate}
            \item Hệ thống xác định \texttt{BaiLam} đang hoạt động và câu hỏi hiện tại.
            \item Người học nhập hoặc chọn câu trả lời theo dạng câu hỏi.
            \item Hệ thống lưu/cập nhật \texttt{ChiTietBaiLam} gắn với \texttt{MaBaiLam} và \texttt{MaCauHoi}.
            \item Hệ thống có thể ghi nhận thời gian suy nghĩ/xử lý câu hỏi.
            \item Người học có thể chuyển giữa các câu khi quy tắc đề cho phép.
            \item Hệ thống không chấp nhận câu trả lời gửi sau khi bài đã nộp hoặc hết thời gian.
        \end{enumerate}

    \end{itemize}

    \item \textbf{Business rules:}

    BR-36, BR-38, BR-41.

    \ScreenLayoutItem

    \begin{figure}[H]
        \centering
        \SafeIncludeGraphics[width=0.94\textwidth,height=0.78\textheight,keepaspectratio]{figures/03_02/exam.png}
        \caption{Màn hình làm bài thi/luyện tập và điều hướng câu hỏi.}
        \label{fig:screen-exam}
    \end{figure}

\end{itemize}

\subsubsection{Submit or Automatically Submit an Examination}
\label{subsubsec:submit-examination}

\begin{itemize}

    \item \textbf{Function trigger:}

    Người học chọn ``Nộp bài'' hoặc thời gian làm bài hết.

    \item \textbf{Function description:}

    \begin{itemize}

        \item \textbf{Actors/Roles:}
        \begin{itemize}
            \item Học sinh / Sinh viên.
            \item Hệ thống.
        \end{itemize}

        \item \textbf{Purpose:}
        \begin{enumerate}
            \item Kết thúc lần làm bài và chuyển sang quy trình chấm.
            \item Đảm bảo bài được tự động nộp khi hết thời gian.
        \end{enumerate}

        \item \textbf{Data processing:}
        \begin{enumerate}
            \item Khi người học chủ động nộp, hệ thống yêu cầu xác nhận trước khi khóa bài.
            \item Khi hết thời gian, hệ thống tự động kết thúc và nộp theo dữ liệu đã lưu.
            \item Hệ thống cập nhật \texttt{ThoiGianNop} và \texttt{TongThoiGianLam}.
            \item Hệ thống khóa việc chỉnh sửa câu trả lời sau khi nộp.
            \item Hệ thống xác định hình thức chấm: Tự động, AI hỗ trợ hoặc Giáo viên.
            \item Phần có thể chấm tự động được đối chiếu đáp án và cập nhật điểm chi tiết.
            \item Phần cần giáo viên/AI hỗ trợ được chuyển trạng thái chờ chấm.
            \item Sau khi chấm hoàn tất, \texttt{TongDiem} và dữ liệu phân tích được cập nhật.
        \end{enumerate}

    \end{itemize}

    \item \textbf{Business rules:}

    BR-38, BR-39, BR-40, BR-41, BR-42.

    \ScreenLayoutItem

    \ScreenLayoutFigure{figures/03_02/exam.png}{Submit or Automatically Submit an Examination}

\end{itemize}


% =========================================================
% LEARNING RESULTS & ANALYTICS
% =========================================================

\subsection{Learning Results \& Analytics}
\label{subsec:learning-results-analytics}

Phân hệ kết quả và phân tích hiển thị điểm, năng lực, Knowledge Tracing, dự đoán và khuyến nghị học tập.

\subsubsection{View Examination Result and Competency Analysis}
\label{subsubsec:view-exam-result-analysis}

\begin{itemize}

    \item \textbf{Function trigger:}

    Sau khi bài được chấm hoàn tất, người học chọn kết quả của bài thi/luyện tập.

    \item \textbf{Function description:}

    \begin{itemize}

        \item \textbf{Actors/Roles:}
        \begin{itemize}
            \item Học sinh / Sinh viên.
        \end{itemize}

        \item \textbf{Purpose:}
        \begin{enumerate}
            \item Hiển thị điểm và kết quả theo câu/nhóm kiến thức.
            \item Trình bày điểm mạnh, điểm yếu và phân tích năng lực.
        \end{enumerate}

        \item \textbf{Data processing:}
        \begin{enumerate}
            \item Hệ thống kiểm tra \texttt{BaiLam} thuộc người học hiện tại.
            \item Hệ thống lấy \texttt{TongDiem}, hình thức chấm và các \texttt{ChiTietBaiLam}.
            \item Hệ thống tổng hợp đúng/sai, điểm thành phần và dữ liệu theo câu hỏi/nhóm kiến thức.
            \item Dữ liệu câu hỏi được liên hệ với \texttt{MaTranKienThuc} để phục vụ phân tích.
            \item Hệ thống hiển thị điểm mạnh, điểm yếu và phần kiến thức cần cải thiện.
            \item Kết quả được lưu trong lịch sử học tập và dùng để cập nhật lộ trình/tiến độ.
            \item Nếu chưa đạt mục tiêu, hệ thống có thể đề xuất luyện tập thêm hoặc thi lại.
        \end{enumerate}

    \end{itemize}

    \item \textbf{Business rules:}

    BR-41, BR-42, BR-43.

    \ScreenLayoutItem

    \begin{figure}[H]
        \centering
        \SafeIncludeGraphics[width=0.94\textwidth,height=0.78\textheight,keepaspectratio]{figures/03_02/student_results.png}
        \caption{Màn hình Kết quả luyện tập của Học sinh/Sinh viên.}
        \label{fig:screen-student-results}
    \end{figure}

\end{itemize}

\subsubsection{View Knowledge Tracing Result}
\label{subsubsec:view-knowledge-tracing}

\begin{itemize}

    \item \textbf{Function trigger:}

    Người học hoặc Giáo viên trong phạm vi được phép mở kết quả Knowledge Tracing.

    \item \textbf{Function description:}

    \begin{itemize}

        \item \textbf{Actors/Roles:}
        \begin{itemize}
            \item Học sinh / Sinh viên.
            \item Giáo viên / Giảng viên.
        \end{itemize}

        \item \textbf{Purpose:}
        \begin{enumerate}
            \item Hiển thị mức độ thành thạo kiến thức ước lượng từ dữ liệu học tập.
            \item Hỗ trợ xác định chủ đề cần ưu tiên.
        \end{enumerate}

        \item \textbf{Data processing:}
        \begin{enumerate}
            \item Hệ thống xác định học sinh cần xem và kiểm tra phạm vi quyền.
            \item Hệ thống tổng hợp lịch sử bài làm, tiến độ và ma trận kiến thức liên quan.
            \item Mô hình Knowledge Tracing xử lý dữ liệu để ước lượng mức độ thành thạo theo chủ đề.
            \item Hệ thống hiển thị kết quả theo nhóm kiến thức và xu hướng khi đủ dữ liệu.
            \item Giáo viên chỉ xem học sinh thuộc lớp/phạm vi phụ trách; Học sinh chỉ xem chính mình.
        \end{enumerate}

    \end{itemize}

    \item \textbf{Business rules:}

    BR-07, BR-11, BR-42.

    \ScreenLayoutItem

    \begin{figure}[H]
        \centering
        \SafeIncludeGraphics[width=0.94\textwidth,height=0.78\textheight,keepaspectratio]{figures/03_02/student_competency.png}
        \caption{Màn hình Năng lực học thuật và phân tích năng lực tổng hợp.}
        \label{fig:screen-student-competency}
    \end{figure}

\end{itemize}

\subsubsection{View Predicted Examination Score and Ranking}
\label{subsubsec:view-predicted-score-ranking}

\begin{itemize}

    \item \textbf{Function trigger:}

    Người dùng có quyền mở khu vực dự đoán điểm thi/xếp hạng của một học sinh.

    \item \textbf{Function description:}

    \begin{itemize}

        \item \textbf{Actors/Roles:}
        \begin{itemize}
            \item Học sinh / Sinh viên.
            \item Giáo viên / Giảng viên.
            \item Phụ huynh.
        \end{itemize}

        \item \textbf{Purpose:}
        \begin{enumerate}
            \item Hiển thị điểm thi dự đoán và xác suất đạt mục tiêu.
            \item Cung cấp vị trí tương đối khi hệ thống có dữ liệu hỗ trợ.
        \end{enumerate}

        \item \textbf{Data processing:}
        \begin{enumerate}
            \item Hệ thống xác định học sinh cần phân tích và kiểm tra quyền theo vai trò.
            \item Hệ thống lấy lịch sử kết quả, tiến độ, Knowledge Tracing và mục tiêu hiện tại.
            \item Mô hình dự đoán xử lý dữ liệu và trả về điểm dự kiến cùng chỉ số liên quan.
            \item Hệ thống so sánh dự đoán với điểm mục tiêu.
            \item Nếu có xếp hạng, hệ thống chỉ hiển thị dữ liệu tổng hợp phù hợp chính sách bảo mật.
            \item Phụ huynh chỉ xem học sinh đã liên kết; Giáo viên chỉ xem phạm vi phụ trách.
        \end{enumerate}

    \end{itemize}

    \item \textbf{Business rules:}

    BR-11, BR-13, BR-17, BR-42.

    \ScreenLayoutItem

    \ScreenLayoutFigure{figures/03_02/student_dashboard.png}{View Predicted Examination Score and Ranking}

\end{itemize}

\subsubsection{View Recommended Learning Materials}
\label{subsubsec:view-recommended-materials}

\begin{itemize}

    \item \textbf{Function trigger:}

    Người học mở khu vực gợi ý tài liệu hoặc hệ thống hiển thị gợi ý sau phân tích kết quả.

    \item \textbf{Function description:}

    \begin{itemize}

        \item \textbf{Actors/Roles:}
        \begin{itemize}
            \item Học sinh / Sinh viên.
        \end{itemize}

        \item \textbf{Purpose:}
        \begin{enumerate}
            \item Đề xuất bài học/tài liệu phù hợp phần kiến thức còn yếu.
            \item Kết nối phân tích năng lực với nội dung học tiếp theo.
        \end{enumerate}

        \item \textbf{Data processing:}
        \begin{enumerate}
            \item Hệ thống lấy kết quả năng lực, tiến độ và các chủ đề chưa đạt.
            \item Hệ thống tìm \texttt{BaiHoc} và \texttt{TaiLieu} liên quan.
            \item Nếu dùng AI đề xuất, hệ thống áp dụng quy tắc AI và dữ liệu đầu vào phù hợp.
            \item Hệ thống sắp xếp danh sách gợi ý theo mức liên quan/ưu tiên.
            \item Người học có thể mở bài học hoặc tài liệu từ danh sách.
            \item Tài liệu không còn khả dụng không được đề xuất để truy cập.
        \end{enumerate}

    \end{itemize}

    \item \textbf{Business rules:}

    BR-24, BR-25, BR-33.

    \ScreenLayoutItem

    \ScreenLayoutFigure{figures/03_02/ai_tutor.png}{View Recommended Learning Materials}

\end{itemize}

\subsubsection{View Progress Report and Learning Reminders}
\label{subsubsec:view-progress-report-reminders}

\begin{itemize}

    \item \textbf{Function trigger:}

    Học sinh/Sinh viên chọn ``Báo cáo tiến độ'' hoặc khu vực nhắc nhở.

    \item \textbf{Function description:}

    \begin{itemize}

        \item \textbf{Actors/Roles:}
        \begin{itemize}
            \item Học sinh / Sinh viên.
        \end{itemize}

        \item \textbf{Purpose:}
        \begin{enumerate}
            \item Hiển thị tiến độ lộ trình, thời gian học và kết quả gần đây.
            \item Hiển thị nhiệm vụ, hạn hoàn thành và nhắc nhở.
        \end{enumerate}

        \item \textbf{Data processing:}
        \begin{enumerate}
            \item Hệ thống tổng hợp \texttt{TienDoHocTap}, \texttt{ChiTietLoTrinh}, kết quả bài làm và nhiệm vụ đang áp dụng.
            \item Hệ thống tính mức độ hoàn thành và thời gian học tích lũy.
            \item Hệ thống xác định nhiệm vụ sắp đến hạn hoặc nội dung chưa hoàn thành.
            \item Hệ thống hiển thị báo cáo theo khoảng thời gian phù hợp.
            \item Nhắc nhở có thể được chuyển tới Notification Service theo cấu hình người dùng.
            \item Nếu giai đoạn hoàn thành, báo cáo được dùng để cung cấp gợi ý bước tiếp theo.
        \end{enumerate}

    \end{itemize}

    \item \textbf{Business rules:}

    BR-23, BR-24, BR-56, BR-59.

    \ScreenLayoutItem

    \begin{figure}[H]
        \centering
        \SafeIncludeGraphics[width=0.94\textwidth,height=0.78\textheight,keepaspectratio]{figures/03_02/new_2026_09_10/reporting_hub_student_summary.png}
        \caption{Màn hình Reporting Hub tổng hợp tiến độ và kết quả học tập của Học sinh / Sinh viên.}
        \label{fig:screen-student-progress-report}
    \end{figure}

\end{itemize}


% =========================================================
% MANAGE LEARNING CONTENT
% =========================================================

\subsection{Manage Learning Content}
\label{subsec:manage-learning-content}

Phân hệ dành cho Giáo viên/Quản trị viên quản lý bài học, tài liệu, chương trình giảng dạy và ma trận kiến thức.

\subsubsection{Rate Learning Content and Send Feedback}
\label{subsubsec:rate-learning-content-feedback}

\begin{itemize}

    \item \textbf{Function trigger:}

    Học sinh/Sinh viên đang xem một bài học, câu hỏi hoặc lời giải và chọn chức năng đánh giá/phản hồi nội dung.

    \item \textbf{Function description:}

    \begin{itemize}

        \item \textbf{Actors/Roles:}
        \begin{itemize}
            \item Học sinh / Sinh viên.
        \end{itemize}

        \item \textbf{Purpose:}
        \begin{enumerate}
            \item Cho phép người học đánh giá chất lượng bài học, câu hỏi và lời giải theo STU-04.
            \item Ghi nhận phản hồi để Giáo viên/Quản trị viên có cơ sở rà soát và cải thiện nội dung.
        \end{enumerate}

        \item \textbf{Data processing:}
        \begin{enumerate}
            \item Hệ thống xác định nội dung đang được đánh giá và kiểm tra người dùng đã đăng nhập hợp lệ.
            \item Người học chọn mức đánh giá và/hoặc nhập nội dung phản hồi theo loại nội dung mà giao diện hỗ trợ.
            \item Hệ thống kiểm tra phản hồi không tham chiếu tới tài nguyên nằm ngoài phạm vi truy cập của người học.
            \item Nếu dữ liệu hợp lệ, hệ thống ghi nhận đánh giá/phản hồi gắn với người học và nội dung tương ứng theo mô hình lưu trữ phản hồi được triển khai.
            \item Phản hồi của người học không tự động thay đổi bài học, đáp án đúng hoặc lời giải đang công bố.
            \item Giáo viên/Quản trị viên có thể sử dụng dữ liệu phản hồi trong quy trình kiểm duyệt/chỉnh sửa nội dung theo quyền được phân công.
            \item Nếu hệ thống chưa triển khai bảng phản hồi riêng trong ERD hiện tại, phần lưu trữ phải được bổ sung ở thiết kế dữ liệu trước khi hiện thực chức năng.
        \end{enumerate}

    \end{itemize}

    \item \textbf{Business rules:}

    BR-31, BR-44, BR-47.

    \ScreenLayoutItem

    \ScreenLayoutFigure{figures/03_02/admin_academic_management.png}{Rate Learning Content and Send Feedback}

\end{itemize}


\subsubsection{View Learning Content and Material List}
\label{subsubsec:view-learning-content-list}

\begin{itemize}

    \item \textbf{Function trigger:}

    Giáo viên/Giảng viên hoặc Quản trị viên mở khu vực quản lý nội dung.

    \item \textbf{Function description:}

    \begin{itemize}

        \item \textbf{Actors/Roles:}
        \begin{itemize}
            \item Giáo viên / Giảng viên.
            \item Quản trị viên học vụ.
        \end{itemize}

        \item \textbf{Purpose:}
        \begin{enumerate}
            \item Hiển thị danh sách bài học và tài liệu.
            \item Hỗ trợ tìm kiếm/lọc để truy cập chức năng quản lý.
        \end{enumerate}

        \item \textbf{Data processing:}
        \begin{enumerate}
            \item Hệ thống xác minh quyền người dùng.
            \item Hệ thống tải \texttt{BaiHoc} cùng \texttt{TaiLieu} liên quan.
            \item Thông tin hiển thị có thể gồm tên bài, chủ đề, thời lượng, loại tài liệu và trạng thái.
            \item Giáo viên chỉ thao tác nội dung trong phạm vi được phân quyền; Quản trị viên theo phạm vi toàn hệ thống.
            \item Người dùng chọn một bài học/tài liệu để xem chi tiết hoặc chuyển sang tạo/cập nhật.
        \end{enumerate}

    \end{itemize}

    \item \textbf{Business rules:}

    BR-07, BR-44, BR-48.

    \ScreenLayoutItem

    \begin{figure}[H]
        \centering
        \SafeIncludeGraphics[width=0.94\textwidth,height=0.78\textheight,keepaspectratio]{figures/03_02/admin_academic_management.png}
        \caption{Màn hình Quản lý Học thuật và tài liệu học tập.}
        \label{fig:screen-academic-management}
    \end{figure}

\end{itemize}

\subsubsection{Create or Update Learning Content}
\label{subsubsec:create-update-learning-content}

\begin{itemize}

    \item \textbf{Function trigger:}

    Người có quyền chọn tạo bài học mới hoặc chỉnh sửa bài học hiện có.

    \item \textbf{Function description:}

    \begin{itemize}

        \item \textbf{Actors/Roles:}
        \begin{itemize}
            \item Giáo viên / Giảng viên.
            \item Quản trị viên học vụ.
        \end{itemize}

        \item \textbf{Purpose:}
        \begin{enumerate}
            \item Tạo/cập nhật nội dung bài học.
            \item Duy trì chủ đề, mô tả và thời lượng dự kiến.
        \end{enumerate}

        \item \textbf{Data processing:}
        \begin{enumerate}
            \item Hệ thống tải dữ liệu hiện tại khi chỉnh sửa.
            \item Người dùng nhập/cập nhật tên bài học, mô tả, chủ đề và thời lượng dự kiến.
            \item Hệ thống kiểm tra trường bắt buộc và phạm vi quyền.
            \item Nếu hợp lệ, hệ thống tạo mới hoặc cập nhật \texttt{BaiHoc}.
            \item Nếu nội dung nằm trong quy trình kiểm duyệt, trạng thái xuất bản chỉ đổi sau bước phê duyệt.
            \item Hệ thống không cho phép người dùng ngoài phạm vi chỉnh sửa.
        \end{enumerate}

    \end{itemize}

    \item \textbf{Business rules:}

    BR-44, BR-47.

    \ScreenLayoutItem

    \ScreenLayoutFigure{figures/03_02/new_2026_09_10/teacher_content_management.png}{Create or Update Learning Content}

\end{itemize}

\subsubsection{Manage Learning Materials}
\label{subsubsec:manage-learning-materials}

\begin{itemize}

    \item \textbf{Function trigger:}

    Người có quyền mở phần tài liệu của một bài học.

    \item \textbf{Function description:}

    \begin{itemize}

        \item \textbf{Actors/Roles:}
        \begin{itemize}
            \item Giáo viên / Giảng viên.
            \item Quản trị viên học vụ.
        \end{itemize}

        \item \textbf{Purpose:}
        \begin{enumerate}
            \item Thêm, cập nhật hoặc loại bỏ tài liệu gắn với bài học.
            \item Phân loại tài liệu theo dạng hệ thống hỗ trợ.
        \end{enumerate}

        \item \textbf{Data processing:}
        \begin{enumerate}
            \item Hệ thống xác định \texttt{BaiHoc} đang được quản lý.
            \item Hệ thống tải các \texttt{TaiLieu} gắn với bài học.
            \item Người dùng thêm tên, loại tài liệu và đường dẫn URL hoặc dữ liệu tham chiếu.
            \item Hệ thống kiểm tra bài học tồn tại và quyền quản lý.
            \item Khi lưu, tài liệu mới được gắn với \texttt{MaBaiHoc}.
            \item Khi cập nhật/xóa, hệ thống kiểm tra ảnh hưởng tới nội dung đang sử dụng trước khi áp dụng.
        \end{enumerate}

    \end{itemize}

    \item \textbf{Business rules:}

    BR-25, BR-44.

    \ScreenLayoutItem

    \ScreenLayoutFigure{figures/03_02/new_2026_09_10/admin_learning_materials.png}{Manage Learning Materials}

\end{itemize}

\subsubsection{Manage Curriculum and Learning Sequence}
\label{subsubsec:manage-curriculum}

\begin{itemize}

    \item \textbf{Function trigger:}

    Giáo viên/Giảng viên hoặc Quản trị viên chọn cập nhật chương trình giảng dạy.

    \item \textbf{Function description:}

    \begin{itemize}

        \item \textbf{Actors/Roles:}
        \begin{itemize}
            \item Giáo viên / Giảng viên.
            \item Quản trị viên học vụ.
        \end{itemize}

        \item \textbf{Purpose:}
        \begin{enumerate}
            \item Điều chỉnh cấu trúc môn học, nội dung và thứ tự bài học.
            \item Duy trì mục tiêu giảng dạy và trật tự nội dung dùng cho lộ trình.
        \end{enumerate}

        \item \textbf{Data processing:}
        \begin{enumerate}
            \item Hệ thống tải cấu trúc chương trình/nội dung hiện tại.
            \item Người dùng thay đổi thứ tự, nhóm chủ đề, bài học hoặc thông tin chương trình theo quyền.
            \item Hệ thống kiểm tra tham chiếu tới \texttt{BaiHoc} và lộ trình đang sử dụng.
            \item Thay đổi không được làm mất dữ liệu tiến độ lịch sử.
            \item Nếu ảnh hưởng lộ trình hiện tại, hệ thống đánh dấu dữ liệu liên quan để đánh giá/cập nhật.
            \item Hệ thống lưu cấu trúc mới khi hợp lệ.
        \end{enumerate}

    \end{itemize}

    \item \textbf{Business rules:}

    BR-21, BR-22, BR-44.

    \ScreenLayoutItem

    \ScreenLayoutFigure{figures/03_02/new_2026_09_10/admin_curriculum_sequence.png}{Manage Curriculum and Learning Sequence}

\end{itemize}

\subsubsection{View AI Content and Question Recommendations}
\label{subsubsec:view-ai-content-question-recommendations}

\begin{itemize}

    \item \textbf{Function trigger:}

    Giáo viên/Giảng viên chọn chức năng đề xuất nội dung bởi AI từ khu vực quản lý nội dung hoặc ngân hàng câu hỏi.

    \item \textbf{Function description:}

    \begin{itemize}

        \item \textbf{Actors/Roles:}
        \begin{itemize}
            \item Giáo viên / Giảng viên.
            \item AI Service.
        \end{itemize}

        \item \textbf{Purpose:}
        \begin{enumerate}
            \item Cung cấp đề xuất AI về tài liệu, câu hỏi, mức độ khó và nội dung cần bổ sung/điều chỉnh theo TEA-10.
            \item Cho phép Giáo viên xem xét đề xuất trước khi sử dụng trong nội dung chính thức.
        \end{enumerate}

        \item \textbf{Data processing:}
        \begin{enumerate}
            \item Hệ thống xác định phạm vi môn/chủ đề/lớp mà Giáo viên đang xem và kiểm tra quyền.
            \item Hệ thống tổng hợp dữ liệu phù hợp như kết quả học tập, Knowledge Tracing, độ khó câu hỏi và mức sử dụng nội dung.
            \item AI Service tạo đề xuất về tài liệu, câu hỏi, độ khó hoặc nội dung cần bổ sung; với yêu cầu sinh câu hỏi, hệ thống sử dụng luồng AI-03/AI-04.
            \item Hệ thống hiển thị nguồn/ngữ cảnh liên quan khi dữ liệu hỗ trợ để Giáo viên có thể đánh giá đề xuất.
            \item Giáo viên có thể chấp nhận, chỉnh sửa hoặc bỏ qua đề xuất theo chức năng triển khai.
            \item Nội dung do AI đề xuất không được tự động trở thành nội dung chính thức nếu chưa qua quy trình kiểm soát/duyệt của Giáo viên hoặc hệ thống theo chính sách.
            \item Nếu AI Service không trả kết quả hợp lệ, hệ thống hiển thị trạng thái lỗi/không có đề xuất và không tạo nội dung giả.
        \end{enumerate}

    \end{itemize}

    \item \textbf{Business rules:}

    BR-33, BR-44, BR-47.

    \ScreenLayoutItem

    \begin{figure}[H]
        \centering
        \SafeIncludeGraphics[width=0.94\textwidth,height=0.78\textheight,keepaspectratio]{figures/03_02/new_2026_09_10/teacher_ai_recommendations.png}
        \caption{Màn hình Giáo viên yêu cầu đề xuất nội dung và câu hỏi từ AI.}
        \label{fig:screen-question-bank}
    \end{figure}

\end{itemize}


\subsubsection{Manage Knowledge Matrix and Evaluation Criteria}
\label{subsubsec:manage-knowledge-matrix}

\begin{itemize}

    \item \textbf{Function trigger:}

    Giáo viên/Giảng viên hoặc Quản trị viên mở quản lý ma trận kiến thức/tiêu chí.

    \item \textbf{Function description:}

    \begin{itemize}

        \item \textbf{Actors/Roles:}
        \begin{itemize}
            \item Giáo viên / Giảng viên.
            \item Quản trị viên học vụ.
        \end{itemize}

        \item \textbf{Purpose:}
        \begin{enumerate}
            \item Gắn câu hỏi với chủ đề Knowledge Tracing và thông tin độ khó.
            \item Duy trì dữ liệu phục vụ phân tích năng lực.
        \end{enumerate}

        \item \textbf{Data processing:}
        \begin{enumerate}
            \item Hệ thống tải câu hỏi và \texttt{MaTranKienThuc} liên quan.
            \item Người dùng tạo/cập nhật chủ đề Knowledge Tracing, độ khó và tiêu chí được hỗ trợ.
            \item Mỗi câu hỏi chỉ có một bản ghi ma trận theo quan hệ 1--1 đã thiết kế.
            \item Hệ thống kiểm tra \texttt{MaCauHoi} tồn tại trước khi lưu.
            \item Thay đổi được dùng cho các lần phân tích tiếp theo và không xóa kết quả lịch sử.
            \item Người ngoài phạm vi không được thay đổi ma trận.
        \end{enumerate}

    \end{itemize}

    \item \textbf{Business rules:}

    BR-36, BR-42, BR-44.

    \ScreenLayoutItem

    \ScreenLayoutFigure{figures/03_02/new_2026_09_10/teacher_knowledge_matrix.png}{Manage Knowledge Matrix and Evaluation Criteria -- Teacher scope}

    \ScreenLayoutFigure{figures/03_02/new_2026_09_10/question_matrix_assignment.png}{Assign a Question to a Knowledge Matrix and Difficulty}

    \ScreenLayoutFigure{figures/03_02/new_2026_09_10/admin_global_matrix.png}{Manage Knowledge Matrix and Evaluation Criteria -- Global scope}

\end{itemize}


\subsubsection{Manage Common Content Categories}
\label{subsubsec:manage-common-content-categories}

\begin{itemize}

    \item \textbf{Function trigger:}

    Quản trị viên học vụ mở chức năng quản lý danh mục nội dung dùng chung.

    \item \textbf{Function description:}

    \begin{itemize}

        \item \textbf{Actors/Roles:}
        \begin{itemize}
            \item Quản trị viên học vụ.
        \end{itemize}

        \item \textbf{Purpose:}
        \begin{enumerate}
            \item Quản lý môn học, chủ đề, khối lớp, loại tài liệu và các danh mục dùng chung theo ADM-03.
            \item Chuẩn hóa dữ liệu phân loại được sử dụng bởi bài học, tài liệu, câu hỏi và báo cáo.
        \end{enumerate}

        \item \textbf{Data processing:}
        \begin{enumerate}
            \item Hệ thống xác minh người dùng có quyền Quản trị viên học vụ.
            \item Hệ thống tải các danh mục nội dung đã được triển khai trong hệ thống.
            \item Quản trị viên có thể tạo, cập nhật, kích hoạt/vô hiệu hóa hoặc sắp xếp các giá trị danh mục theo chính sách.
            \item Hệ thống kiểm tra tên/mã danh mục bắt buộc và tính duy nhất khi quy tắc dữ liệu yêu cầu.
            \item Trước khi xóa hoặc vô hiệu hóa một giá trị đang được tham chiếu, hệ thống kiểm tra các bài học, tài liệu, câu hỏi hoặc cấu hình liên quan để tránh làm mất tính toàn vẹn dữ liệu.
            \item Nếu ERD hiện tại chưa có các bảng danh mục tương ứng, thiết kế dữ liệu phải được bổ sung trước khi triển khai thao tác CRUD thực tế.
            \item Sau khi cập nhật, các màn hình tạo nội dung sử dụng danh mục mới trong phạm vi được cấu hình.
        \end{enumerate}

    \end{itemize}

    \item \textbf{Business rules:}

    BR-44, BR-48.

    \ScreenLayoutItem

    \ScreenLayoutFigure{figures/03_02/new_2026_09_10/admin_common_category.png}{Manage Common Content Categories}

\end{itemize}


% =========================================================
% MANAGE QUESTION BANK & EXAMINATIONS
% =========================================================

\subsection{Manage Question Bank \& Examinations}
\label{subsec:manage-question-bank-examinations}

Phân hệ hỗ trợ Giáo viên quản lý ngân hàng câu hỏi và tạo/cấu hình đề thi; Quản trị viên kiểm duyệt khi quy trình yêu cầu.

\subsubsection{View Question Bank}
\label{subsubsec:view-question-bank}

\begin{itemize}

    \item \textbf{Function trigger:}

    Giáo viên/Giảng viên mở ``Ngân hàng câu hỏi''.

    \item \textbf{Function description:}

    \begin{itemize}

        \item \textbf{Actors/Roles:}
        \begin{itemize}
            \item Giáo viên / Giảng viên.
            \item Quản trị viên học vụ.
        \end{itemize}

        \item \textbf{Purpose:}
        \begin{enumerate}
            \item Hiển thị câu hỏi, đáp án, độ khó và phân loại.
            \item Cung cấp điểm vào cho tạo, cập nhật và xây dựng đề.
        \end{enumerate}

        \item \textbf{Data processing:}
        \begin{enumerate}
            \item Hệ thống kiểm tra quyền truy cập.
            \item Hệ thống tải \texttt{CauHoi} cùng ma trận kiến thức liên quan.
            \item Danh sách có thể lọc theo chủ đề, dạng câu hỏi, độ khó hoặc tiêu chí được triển khai.
            \item Giáo viên chỉ quản lý nội dung trong phạm vi được phân quyền.
            \item Người dùng chọn câu hỏi để xem chi tiết hoặc chỉnh sửa.
        \end{enumerate}

    \end{itemize}

    \item \textbf{Business rules:}

    BR-44.

    \ScreenLayoutItem

    \ScreenLayoutFigure{figures/03_02/question_bank.png}{View Question Bank}

\end{itemize}

\subsubsection{Create or Update a Question, Answer and Solution}
\label{subsubsec:create-update-question}

\begin{itemize}

    \item \textbf{Function trigger:}

    Giáo viên chọn tạo câu hỏi mới hoặc chỉnh sửa câu hỏi hiện có.

    \item \textbf{Function description:}

    \begin{itemize}

        \item \textbf{Actors/Roles:}
        \begin{itemize}
            \item Giáo viên / Giảng viên.
        \end{itemize}

        \item \textbf{Purpose:}
        \begin{enumerate}
            \item Tạo/cập nhật câu hỏi, đáp án đúng, độ khó và phân loại.
            \item Chuẩn bị dữ liệu cho đề luyện tập/thi.
        \end{enumerate}

        \item \textbf{Data processing:}
        \begin{enumerate}
            \item Hệ thống hiển thị biểu mẫu câu hỏi; khi chỉnh sửa, tải \texttt{CauHoi} hiện tại.
            \item Giáo viên nhập nội dung, đáp án đúng, độ khó, khung năng lực và dạng câu hỏi.
            \item Hệ thống kiểm tra trường bắt buộc và định dạng theo dạng câu hỏi.
            \item Nếu hợp lệ, hệ thống tạo mới hoặc cập nhật \texttt{CauHoi}.
            \item Thông tin ma trận kiến thức được tạo/cập nhật qua chức năng tương ứng.
            \item Câu hỏi đang được tham chiếu bởi đề thi cần được xử lý thận trọng để không làm sai dữ liệu lịch sử.
            \item Nếu nội dung cần kiểm duyệt, việc xuất bản phụ thuộc trạng thái phê duyệt.
        \end{enumerate}

    \end{itemize}

    \item \textbf{Business rules:}

    BR-44, BR-47.

    \ScreenLayoutItem

    \ScreenLayoutFigure{figures/03_02/question_bank.png}{Create or Update a Question, Answer and Solution}

\end{itemize}

\subsubsection{Create an Examination}
\label{subsubsec:create-examination}

\begin{itemize}

    \item \textbf{Function trigger:}

    Giáo viên chọn ``Tạo đề thi'' từ khu vực ngân hàng câu hỏi/đề thi.

    \item \textbf{Function description:}

    \begin{itemize}

        \item \textbf{Actors/Roles:}
        \begin{itemize}
            \item Giáo viên / Giảng viên.
        \end{itemize}

        \item \textbf{Purpose:}
        \begin{enumerate}
            \item Tạo đề thi/luyện tập mới.
            \item Gắn đề với kỳ thi, lớp hoặc giáo viên khi phù hợp.
        \end{enumerate}

        \item \textbf{Data processing:}
        \begin{enumerate}
            \item Giáo viên nhập tên đề, loại đề, thời gian làm bài và cấu hình cơ bản.
            \item Hệ thống cho phép liên kết \texttt{DeThi} với \texttt{KyThi} khi phù hợp.
            \item Hệ thống ghi nhận giáo viên tạo đề và lớp liên quan nếu quy trình yêu cầu.
            \item Giáo viên chọn câu hỏi từ ngân hàng.
            \item Hệ thống tạo \texttt{DeThi\_CauHoi} cho từng câu hỏi, lưu thứ tự và điểm thành phần.
            \item Hệ thống kiểm tra cấu hình điểm/thời gian trước khi lưu.
            \item Đề mới chưa được công bố nếu quy trình yêu cầu kiểm duyệt.
        \end{enumerate}

    \end{itemize}

    \item \textbf{Business rules:}

    BR-35, BR-36, BR-45, BR-47.

    \ScreenLayoutItem

    \ScreenLayoutFigure{figures/03_02/question_bank.png}{Create an Examination}

\end{itemize}

\subsubsection{Configure Examination Answer Key, Scoring and Availability}
\label{subsubsec:configure-examination}

\begin{itemize}

    \item \textbf{Function trigger:}

    Giáo viên mở cấu hình của một đề thi đã tạo.

    \item \textbf{Function description:}

    \begin{itemize}

        \item \textbf{Actors/Roles:}
        \begin{itemize}
            \item Giáo viên / Giảng viên.
        \end{itemize}

        \item \textbf{Purpose:}
        \begin{enumerate}
            \item Thiết lập nội dung, đáp án, thang điểm, thời gian và điều kiện mở đề.
            \item Chuẩn bị đề trước khi giao/đưa vào kỳ thi.
        \end{enumerate}

        \item \textbf{Data processing:}
        \begin{enumerate}
            \item Hệ thống tải \texttt{DeThi} và \texttt{DeThi\_CauHoi}.
            \item Giáo viên cập nhật thứ tự, điểm thành phần, đáp án/giải pháp và thời gian làm bài.
            \item Hệ thống kiểm tra câu hỏi thuộc đề và dữ liệu điểm không mâu thuẫn.
            \item Thời gian mở/đóng hoặc điều kiện khả dụng được ghi nhận theo cấu hình triển khai.
            \item Nếu đề đã có bài làm, hệ thống hạn chế thay đổi làm sai lệch lịch sử.
            \item Sau khi lưu, đề chuyển tới bước kiểm duyệt/công bố theo luồng.
        \end{enumerate}

    \end{itemize}

    \item \textbf{Business rules:}

    BR-36, BR-45, BR-47.

    \ScreenLayoutItem

    \ScreenLayoutFigure{figures/03_02/question_bank.png}{Configure Examination Answer Key, Scoring and Availability}

\end{itemize}

\subsubsection{Review and Approve Examination Content}
\label{subsubsec:review-approve-exam}

\begin{itemize}

    \item \textbf{Function trigger:}

    Đề thi/nội dung được gửi vào quy trình kiểm duyệt trước khi xuất bản.

    \item \textbf{Function description:}

    \begin{itemize}

        \item \textbf{Actors/Roles:}
        \begin{itemize}
            \item Quản trị viên học vụ.
        \end{itemize}

        \item \textbf{Purpose:}
        \begin{enumerate}
            \item Kiểm tra nội dung thuộc quy trình cần phê duyệt.
            \item Ngăn nội dung chưa duyệt được công bố khi chính sách yêu cầu.
        \end{enumerate}

        \item \textbf{Data processing:}
        \begin{enumerate}
            \item Hệ thống hiển thị danh sách nội dung/đề chờ kiểm duyệt.
            \item Quản trị viên mở chi tiết đề, cấu trúc câu hỏi và cấu hình.
            \item Quản trị viên phê duyệt hoặc yêu cầu chỉnh sửa.
            \item Chỉ nội dung đáp ứng điều kiện mới được chuyển sang trạng thái cho phép xuất bản/sử dụng.
            \item Nếu bị từ chối/yêu cầu sửa, đề không được mở cho người học.
            \item Hệ thống lưu trạng thái kiểm duyệt theo cơ chế dữ liệu triển khai.
        \end{enumerate}

    \end{itemize}

    \item \textbf{Business rules:}

    BR-47, BR-48.

    \ScreenLayoutItem

    \ScreenLayoutFigure{figures/03_02/new_2026_09_10/admin_exam_review_queue.png}{Review and Approve Examination Content}

\end{itemize}

% =========================================================
% MANAGE CLASSES, ASSIGNMENTS & ATTENDANCE
% =========================================================

\subsection{Manage Classes, Assignments \& Attendance}
\label{subsec:manage-classes-assignments-attendance}

Phân hệ lớp học hỗ trợ Giáo viên xem lớp, giao đề/bài tập, theo dõi bài nộp, chấm bài và quản lý chuyên cần.

\subsubsection{View Teacher Dashboard}
\label{subsubsec:view-teacher-dashboard}

\begin{itemize}

    \item \textbf{Function trigger:}

    Sau khi Giáo viên/Giảng viên đăng nhập thành công, hệ thống điều hướng tới Teacher Dashboard.

    \item \textbf{Function description:}

    \begin{itemize}

        \item \textbf{Actors/Roles:}
        \begin{itemize}
            \item Giáo viên / Giảng viên.
        \end{itemize}

        \item \textbf{Purpose:}
        \begin{enumerate}
            \item Cung cấp điểm truy cập tập trung tới lớp học, ngân hàng câu hỏi, đề thi, giao bài, chấm bài và Learning Analytics.
            \item Hiển thị thông tin tổng quan phục vụ hoạt động giảng dạy và theo dõi học sinh.
        \end{enumerate}

        \item \textbf{Data processing:}
        \begin{enumerate}
            \item Hệ thống xác minh vai trò Giáo viên/Giảng viên.
            \item Hệ thống tải các lớp thuộc phạm vi phụ trách và các thông tin tổng quan như bài đã giao, bài cần chấm, kết quả/tiến độ lớp và thông báo liên quan theo dữ liệu hiện có.
            \item Dashboard cung cấp lối truy cập tới Quản lý lớp, Ngân hàng câu hỏi, Quản lý đề thi, Giao bài, Danh sách bài nộp, Chấm bài, Learning Analytics và báo cáo theo Screen Flow Giáo viên.
            \item Dữ liệu lớp/học sinh trên Dashboard chỉ thuộc phạm vi Giáo viên phụ trách.
            \item Nếu Giáo viên chọn một chức năng, hệ thống tiếp tục kiểm tra quyền và phạm vi trước khi hiển thị dữ liệu chi tiết.
            \item Các đề xuất AI được hiển thị như nội dung hỗ trợ và không tự thay thế quyết định chuyên môn của Giáo viên.
        \end{enumerate}

    \end{itemize}

    \item \textbf{Business rules:}

    BR-05, BR-07, BR-09, BR-11.

    \ScreenLayoutItem

    \begin{figure}[H]
        \centering
        \SafeIncludeGraphics[width=0.94\textwidth,height=0.78\textheight,keepaspectratio]{figures/03_02/new_2026_09_10/teacher_dashboard_summary.png}
        \caption{Màn hình Dashboard tổng quan lớp học đang chọn của Giáo viên/Giảng viên.}
        \label{fig:screen-teacher-dashboard}
    \end{figure}

\end{itemize}


\subsubsection{View Class and Student List}
\label{subsubsec:view-class-student-list}

\begin{itemize}

    \item \textbf{Function trigger:}

    Giáo viên mở ``Quản lý lớp'' từ Teacher Dashboard.

    \item \textbf{Function description:}

    \begin{itemize}

        \item \textbf{Actors/Roles:}
        \begin{itemize}
            \item Giáo viên / Giảng viên.
            \item Quản trị viên học vụ.
        \end{itemize}

        \item \textbf{Purpose:}
        \begin{enumerate}
            \item Hiển thị các lớp giáo viên phụ trách và danh sách học sinh.
            \item Cung cấp dữ liệu đầu vào cho giao bài, báo cáo và điểm danh.
        \end{enumerate}

        \item \textbf{Data processing:}
        \begin{enumerate}
            \item Hệ thống lấy các \texttt{LopHoc} có \texttt{MaGiaoVien} phù hợp với giáo viên đang đăng nhập.
            \item Khi chọn lớp, hệ thống tải học sinh qua \texttt{HocSinh\_LopHoc}.
            \item Chỉ học sinh có quan hệ tham gia lớp phù hợp được hiển thị trong danh sách hiện hành.
            \item Giáo viên không được xem/quản lý lớp ngoài phạm vi phụ trách bằng cách thay đổi tham số request.
            \item Quản trị viên có thể xem dữ liệu lớp theo phạm vi quản trị.
        \end{enumerate}

    \end{itemize}

    \item \textbf{Business rules:}

    BR-09, BR-10, BR-11.

    \ScreenLayoutItem

    \begin{figure}[H]
        \centering
        \SafeIncludeGraphics[width=0.94\textwidth,height=0.78\textheight,keepaspectratio]{figures/03_02/teacher_student_list.png}
        \caption{Màn hình danh sách học sinh phục vụ theo dõi lớp học.}
        \label{fig:screen-teacher-student-list}
    \end{figure}

\end{itemize}

\subsubsection{Assign an Assessment to a Class}
\label{subsubsec:assign-assessment-class}

\begin{itemize}

    \item \textbf{Function trigger:}

    Giáo viên chọn ``Giao bài'' tại một lớp hoặc từ khu vực quản lý đề.

    \item \textbf{Function description:}

    \begin{itemize}

        \item \textbf{Actors/Roles:}
        \begin{itemize}
            \item Giáo viên / Giảng viên.
            \item Notification Service.
        \end{itemize}

        \item \textbf{Purpose:}
        \begin{enumerate}
            \item Giao một đề/bài đánh giá cho một lớp cụ thể.
            \item Thiết lập thời gian giao, hạn nộp và thông báo cho học sinh.
        \end{enumerate}

        \item \textbf{Data processing:}
        \begin{enumerate}
            \item Hệ thống chỉ hiển thị các lớp thuộc phạm vi giáo viên phụ trách.
            \item Giáo viên chọn lớp, \texttt{DeThi}, nhập tiêu đề/mô tả, ngày giao và hạn nộp.
            \item Hệ thống kiểm tra giáo viên là người phụ trách lớp và đề được phép sử dụng.
            \item Hạn nộp phải hợp lệ so với thời điểm giao.
            \item Nếu hợp lệ, hệ thống tạo \texttt{GiaoBaiTap} với \texttt{MaLop}, \texttt{MaDeThi}, \texttt{MaGiaoVien} và thông tin thời hạn.
            \item Hệ thống xác định học sinh của lớp qua \texttt{HocSinh\_LopHoc}.
            \item Nếu bật thông báo, hệ thống tạo \texttt{ThongBao} và các bản ghi \texttt{ThongBaoNguoiNhan}.
            \item Bài được giao xuất hiện trong danh sách nhiệm vụ của học sinh thuộc lớp.
        \end{enumerate}

    \end{itemize}

    \item \textbf{Business rules:}

    BR-11, BR-61, BR-62, BR-63.

    \ScreenLayoutItem

    \ScreenLayoutFigure{figures/03_02/new_2026_09_10/teacher_assign_assessment.png}{Assign an Assessment to a Class}

\end{itemize}

\subsubsection{View Assigned Assessment and Submission List}
\label{subsubsec:view-assignment-submissions}

\begin{itemize}

    \item \textbf{Function trigger:}

    Giáo viên mở một bài đã giao và chọn danh sách bài nộp.

    \item \textbf{Function description:}

    \begin{itemize}

        \item \textbf{Actors/Roles:}
        \begin{itemize}
            \item Giáo viên / Giảng viên.
        \end{itemize}

        \item \textbf{Purpose:}
        \begin{enumerate}
            \item Theo dõi trạng thái Đã nộp, Chưa nộp hoặc Cần chấm.
            \item Truy cập bài làm để chấm/nhận xét.
        \end{enumerate}

        \item \textbf{Data processing:}
        \begin{enumerate}
            \item Hệ thống tải \texttt{GiaoBaiTap} thuộc giáo viên/lớp được phép.
            \item Hệ thống lấy danh sách học sinh của \texttt{MaLop}.
            \item Hệ thống đối chiếu với \texttt{BaiLam} qua \texttt{MaHocSinh}, \texttt{MaDeThi} và \texttt{MaGiaoBaiTap}.
            \item Học sinh có bài đã nộp được đánh dấu Đã nộp; bài cần chấm thủ công/AI hỗ trợ có trạng thái Cần chấm; không có bài hợp lệ là Chưa nộp.
            \item Hệ thống hiển thị thời gian nộp và kết quả nếu đã có.
            \item Giáo viên chọn một bài nộp để xem chi tiết/chấm.
        \end{enumerate}

    \end{itemize}

    \item \textbf{Business rules:}

    BR-11, BR-37, BR-39, BR-61.

    \ScreenLayoutItem

    \ScreenLayoutFigure{figures/03_02/teacher_dashboard.png}{View Assigned Assessment and Submission List}

\end{itemize}

\subsubsection{Grade and Comment on a Student Submission}
\label{subsubsec:grade-comment-submission}

\begin{itemize}

    \item \textbf{Function trigger:}

    Giáo viên chọn một bài nộp cần chấm từ danh sách submission.

    \item \textbf{Function description:}

    \begin{itemize}

        \item \textbf{Actors/Roles:}
        \begin{itemize}
            \item Giáo viên / Giảng viên.
        \end{itemize}

        \item \textbf{Purpose:}
        \begin{enumerate}
            \item Chấm phần bài cần giáo viên đánh giá.
            \item Nhập nhận xét và công bố kết quả trong phạm vi phụ trách.
        \end{enumerate}

        \item \textbf{Data processing:}
        \begin{enumerate}
            \item Hệ thống kiểm tra bài làm thuộc học sinh/lớp giáo viên phụ trách.
            \item Hệ thống tải \texttt{BaiLam} và các \texttt{ChiTietBaiLam}.
            \item Giáo viên xem câu trả lời và dữ liệu đáp án/tiêu chí liên quan.
            \item Giáo viên nhập điểm/nhận xét cho phần cần chấm thủ công.
            \item Hệ thống kiểm tra điểm không vượt phạm vi điểm thành phần.
            \item Khi xác nhận lưu, hệ thống cập nhật điểm chi tiết và tính lại \texttt{TongDiem}.
            \item Kết quả chỉ được công bố khi giáo viên xác nhận/công bố theo Screen Flow.
            \item Sau khi chấm xong, dữ liệu được dùng cho phân tích năng lực và báo cáo.
        \end{enumerate}

    \end{itemize}

    \item \textbf{Business rules:}

    BR-11, BR-40, BR-41, BR-42, BR-46.

    \ScreenLayoutItem

    \ScreenLayoutFigure{figures/03_02/teacher_dashboard.png}{Grade and Comment on a Student Submission}

\end{itemize}

\subsubsection{Evaluate and Comment on Student Progress}
\label{subsubsec:evaluate-comment-student-progress}

\begin{itemize}

    \item \textbf{Function trigger:}

    Giáo viên chọn một học sinh từ lớp phụ trách và mở chức năng đánh giá/nhận xét học sinh.

    \item \textbf{Function description:}

    \begin{itemize}

        \item \textbf{Actors/Roles:}
        \begin{itemize}
            \item Giáo viên / Giảng viên.
        \end{itemize}

        \item \textbf{Purpose:}
        \begin{enumerate}
            \item Cho phép Giáo viên xem kết quả và đưa ra nhận xét về mức độ tiến bộ của từng học sinh theo TEA-03.
            \item Cung cấp dữ liệu nhận xét để Phụ huynh/Học sinh xem trong phạm vi được phép.
        \end{enumerate}

        \item \textbf{Data processing:}
        \begin{enumerate}
            \item Hệ thống xác minh học sinh thuộc lớp/phạm vi Giáo viên phụ trách.
            \item Hệ thống tổng hợp kết quả bài làm, tiến độ lộ trình, tần suất học, Knowledge Tracing và các dữ liệu học tập có liên quan.
            \item Hệ thống hiển thị dữ liệu hiện có để Giáo viên đánh giá mức độ tiến bộ.
            \item Giáo viên nhập nhận xét và đánh giá chuyên môn theo biểu mẫu được triển khai.
            \item Nếu có phân tích AI, hệ thống hiển thị riêng phần phân tích AI để hỗ trợ tham khảo; kết luận chuyên môn của Giáo viên không bị AI tự động ghi đè.
            \item Hệ thống lưu nhận xét/đánh giá theo mô hình dữ liệu đánh giá được triển khai; nếu ERD hiện tại chưa có bảng lưu nhận xét học sinh thì cần bổ sung trước khi hiện thực.
            \item Dữ liệu được cung cấp cho Học sinh/Phụ huynh chỉ trong phạm vi quyền và liên kết hợp lệ.
        \end{enumerate}

    \end{itemize}

    \item \textbf{Business rules:}

    BR-11, BR-42, BR-46.

    \ScreenLayoutItem

    \begin{figure}[H]
        \centering
        \SafeIncludeGraphics[width=0.94\textwidth,height=0.78\textheight,keepaspectratio]{figures/03_02/teacher_class_analysis.png}
        \caption{Màn hình phân tích năng lực lớp và các học sinh cần hỗ trợ.}
        \label{fig:screen-teacher-class-analysis}
    \end{figure}

\end{itemize}


\subsubsection{Create a Class Session}
\label{subsubsec:create-class-session}

\begin{itemize}

    \item \textbf{Function trigger:}

    Giáo viên mở lịch lớp và chọn tạo buổi học mới.

    \item \textbf{Function description:}

    \begin{itemize}

        \item \textbf{Actors/Roles:}
        \begin{itemize}
            \item Giáo viên / Giảng viên.
        \end{itemize}

        \item \textbf{Purpose:}
        \begin{enumerate}
            \item Tạo một buổi học cụ thể làm cơ sở cho lịch học và điểm danh.
            \item Cung cấp dữ liệu lịch học cho Học sinh/Phụ huynh.
        \end{enumerate}

        \item \textbf{Data processing:}
        \begin{enumerate}
            \item Hệ thống chỉ cho phép chọn lớp thuộc phạm vi giáo viên phụ trách.
            \item Giáo viên nhập ngày học, thời gian bắt đầu/kết thúc, nội dung và trạng thái.
            \item Hệ thống kiểm tra thời gian và lớp hợp lệ.
            \item Nếu hợp lệ, hệ thống tạo \texttt{BuoiHoc} gắn với \texttt{MaLop}.
            \item Buổi học được đưa vào lịch của lớp và có thể dùng để gửi nhắc nhở.
            \item Nếu lớp không thuộc phạm vi giáo viên, yêu cầu bị từ chối.
        \end{enumerate}

    \end{itemize}

    \item \textbf{Business rules:}

    BR-09, BR-11, BR-64.

    \ScreenLayoutItem

    \ScreenLayoutFigure{figures/03_02/new_2026_09_10/teacher_session_attendance.png}{Create a Class Session}

\end{itemize}

\subsubsection{Record Student Attendance}
\label{subsubsec:record-attendance}

\begin{itemize}

    \item \textbf{Function trigger:}

    Giáo viên mở một buổi học và chọn chức năng điểm danh.

    \item \textbf{Function description:}

    \begin{itemize}

        \item \textbf{Actors/Roles:}
        \begin{itemize}
            \item Giáo viên / Giảng viên.
        \end{itemize}

        \item \textbf{Purpose:}
        \begin{enumerate}
            \item Ghi nhận trạng thái tham dự theo buổi học.
            \item Tạo dữ liệu chuyên cần cho báo cáo.
        \end{enumerate}

        \item \textbf{Data processing:}
        \begin{enumerate}
            \item Hệ thống tải \texttt{BuoiHoc} và xác định \texttt{MaLop}.
            \item Hệ thống lấy danh sách học sinh đang tham gia lớp.
            \item Giáo viên đánh dấu Có mặt, Vắng, Trễ hoặc Có phép và có thể nhập ghi chú.
            \item Hệ thống tạo/cập nhật \texttt{DiemDanh} cho từng cặp Buổi học--Học sinh.
            \item Hệ thống ngăn điểm danh học sinh không thuộc lớp của buổi học.
            \item Sau khi lưu, dữ liệu được dùng trong báo cáo và màn hình Phụ huynh.
        \end{enumerate}

    \end{itemize}

    \item \textbf{Business rules:}

    BR-11, BR-64, BR-65.

    \ScreenLayoutItem

    \ScreenLayoutFigure{figures/03_02/new_2026_09_10/teacher_session_attendance.png}{Record Student Attendance}

\end{itemize}

\subsubsection{View Attendance Report}
\label{subsubsec:view-attendance-report}

\begin{itemize}

    \item \textbf{Function trigger:}

    Giáo viên mở báo cáo chuyên cần của lớp hoặc Học sinh mở chuyên cần cá nhân.

    \item \textbf{Function description:}

    \begin{itemize}

        \item \textbf{Actors/Roles:}
        \begin{itemize}
            \item Giáo viên / Giảng viên.
            \item Học sinh / Sinh viên.
            \item Quản trị viên học vụ.
        \end{itemize}

        \item \textbf{Purpose:}
        \begin{enumerate}
            \item Tổng hợp chuyên cần theo lớp hoặc người học.
            \item Hiển thị số buổi có mặt, vắng, trễ và có phép.
        \end{enumerate}

        \item \textbf{Data processing:}
        \begin{enumerate}
            \item Hệ thống xác định phạm vi dữ liệu theo vai trò.
            \item Hệ thống lấy các \texttt{BuoiHoc} và \texttt{DiemDanh} tương ứng.
            \item Dữ liệu được tổng hợp theo khoảng thời gian/lớp/người học theo bộ lọc triển khai.
            \item Giáo viên chỉ xem lớp phụ trách; Học sinh chỉ xem chính mình; Quản trị viên theo phạm vi quản trị.
            \item Hệ thống hiển thị chi tiết và số liệu tổng hợp.
        \end{enumerate}

    \end{itemize}

    \item \textbf{Business rules:}

    BR-07, BR-11, BR-65.

    \ScreenLayoutItem

    \ScreenLayoutFigure{figures/03_02/new_2026_09_10/teacher_attendance_report.png}{View Attendance Report}

\end{itemize}


% =========================================================
% PARENT MONITORING
% =========================================================

\subsection{Parent Monitoring}
\label{subsec:parent-monitoring}

Phân hệ Phụ huynh cho phép đăng ký giám sát, liên kết với học sinh, xem tiến độ/kết quả/chuyên cần, thiết lập giới hạn thời gian và nhận cảnh báo.

\subsubsection{View Parent Dashboard}
\label{subsubsec:view-parent-dashboard}

\begin{itemize}

    \item \textbf{Function trigger:}

    Sau khi Phụ huynh đăng nhập thành công, hệ thống điều hướng tới Parent Dashboard.

    \item \textbf{Function description:}

    \begin{itemize}

        \item \textbf{Actors/Roles:}
        \begin{itemize}
            \item Phụ huynh.
        \end{itemize}

        \item \textbf{Purpose:}
        \begin{enumerate}
            \item Cung cấp tổng quan về các học sinh đã liên kết và những thông tin cần theo dõi.
            \item Cung cấp điểm truy cập tới tiến độ, kết quả, dự đoán, lịch học, chuyên cần, cảnh báo và trao đổi với Giáo viên.
        \end{enumerate}

        \item \textbf{Data processing:}
        \begin{enumerate}
            \item Hệ thống xác minh vai trò Phụ huynh.
            \item Hệ thống tải các \texttt{LienKetPhuHuynh} đang hợp lệ của người dùng.
            \item Nếu có nhiều học sinh liên kết, Dashboard cho phép chọn/chuyển hồ sơ con đang theo dõi.
            \item Với hồ sơ được chọn, hệ thống tổng hợp tiến độ gần đây, kết quả, lịch/nhiệm vụ, chuyên cần, cảnh báo và thông báo theo dữ liệu hiện có.
            \item Dashboard cung cấp lối truy cập tới báo cáo tiến độ, kết quả thi/dự đoán điểm, khuyến nghị hỗ trợ, giới hạn thời gian, lịch học/hoạt động, đánh giá từ Giáo viên và trao đổi với Giáo viên theo Screen Flow Phụ huynh.
            \item Trước mỗi truy vấn chi tiết, hệ thống kiểm tra lại quan hệ Phụ huynh--Học sinh; không chỉ dựa vào định danh do client gửi.
        \end{enumerate}

    \end{itemize}

    \item \textbf{Business rules:}

    BR-05, BR-13, BR-14.

    \ScreenLayoutItem

    \begin{figure}[H]
        \centering
        \SafeIncludeGraphics[width=0.94\textwidth,height=0.78\textheight,keepaspectratio]{figures/03_02/parent_dashboard.png}
        \caption{Màn hình Dashboard tổng quan của Phụ huynh.}
        \label{fig:screen-parent-dashboard}
    \end{figure}

\end{itemize}


\subsubsection{Register Parent Monitoring Account and Link a Student}
\label{subsubsec:register-parent-link-student}

\begin{itemize}

    \item \textbf{Function trigger:}

    Phụ huynh chưa có tài khoản chọn đăng ký giám sát hoặc chọn thêm liên kết học sinh.

    \item \textbf{Function description:}

    \begin{itemize}

        \item \textbf{Actors/Roles:}
        \begin{itemize}
            \item Phụ huynh.
        \end{itemize}

        \item \textbf{Purpose:}
        \begin{enumerate}
            \item Tạo tài khoản Phụ huynh và thiết lập quan hệ giám sát.
            \item Đảm bảo chỉ truy cập hồ sơ học sinh đã liên kết.
        \end{enumerate}

        \item \textbf{Data processing:}
        \begin{enumerate}
            \item Hệ thống tạo/nhận diện \texttt{NguoiDung} có \texttt{VaiTro = PhuHuynh} theo quy trình đăng ký.
            \item Phụ huynh nhập hoặc sử dụng mã xác nhận/thông tin liên kết học sinh theo chức năng triển khai.
            \item Hệ thống xác minh học sinh tồn tại và thông tin liên kết hợp lệ.
            \item Nếu thành công, hệ thống tạo \texttt{LienKetPhuHuynh} với \texttt{MaPhuHuynh}, \texttt{MaHocSinh}, mã xác nhận, trạng thái và ngày tạo.
            \item Liên kết chưa hợp lệ/không xác nhận không được dùng để xem dữ liệu học sinh.
            \item Sau khi liên kết hợp lệ, hồ sơ học sinh xuất hiện trong danh sách con em.
        \end{enumerate}

    \end{itemize}

    \item \textbf{Business rules:}

    BR-12, BR-13, BR-14.

    \ScreenLayoutItem

    \ScreenLayoutFigure{figures/03_02/new_2026_09_10/parent_link_student.png}{Link a Student to a Parent Monitoring Account}

\end{itemize}

\subsubsection{Select a Linked Child Profile}
\label{subsubsec:select-child-profile}

\begin{itemize}

    \item \textbf{Function trigger:}

    Phụ huynh đăng nhập và chọn một hồ sơ học sinh trong danh sách liên kết.

    \item \textbf{Function description:}

    \begin{itemize}

        \item \textbf{Actors/Roles:}
        \begin{itemize}
            \item Phụ huynh.
        \end{itemize}

        \item \textbf{Purpose:}
        \begin{enumerate}
            \item Xác định học sinh đang được theo dõi khi có nhiều liên kết.
            \item Áp dụng phạm vi dữ liệu đúng cho các màn hình theo dõi.
        \end{enumerate}

        \item \textbf{Data processing:}
        \begin{enumerate}
            \item Hệ thống tải các \texttt{LienKetPhuHuynh} hợp lệ thuộc Phụ huynh.
            \item Phụ huynh chọn một học sinh.
            \item Hệ thống kiểm tra lại liên kết trước mỗi truy vấn dữ liệu nhạy cảm.
            \item Hồ sơ được chọn trở thành ngữ cảnh cho các màn hình theo dõi.
            \item Phụ huynh có thể chuyển hồ sơ liên kết khác mà không làm lẫn dữ liệu.
        \end{enumerate}

    \end{itemize}

    \item \textbf{Business rules:}

    BR-13, BR-14.

    \ScreenLayoutItem

    \ScreenLayoutFigure{figures/03_02/new_2026_09_10/parent_select_child.png}{Select a Linked Child Profile}

\end{itemize}

\subsubsection{View or Request Child Progress Report}
\label{subsubsec:parent-view-request-progress}

\begin{itemize}

    \item \textbf{Function trigger:}

    Phụ huynh chọn ``Báo cáo tiến độ'' của hồ sơ con hoặc yêu cầu báo cáo theo khoảng thời gian.

    \item \textbf{Function description:}

    \begin{itemize}

        \item \textbf{Actors/Roles:}
        \begin{itemize}
            \item Phụ huynh.
        \end{itemize}

        \item \textbf{Purpose:}
        \begin{enumerate}
            \item Hiển thị tiến độ lộ trình, thời gian học và kết quả gần đây.
            \item Cho phép yêu cầu báo cáo theo khoảng thời gian được hỗ trợ.
        \end{enumerate}

        \item \textbf{Data processing:}
        \begin{enumerate}
            \item Hệ thống xác minh \texttt{LienKetPhuHuynh} với học sinh được chọn.
            \item Hệ thống tổng hợp \texttt{TienDoHocTap}, lộ trình, bài làm và hoạt động gần đây.
            \item Nếu chọn khoảng thời gian, hệ thống áp dụng bộ lọc tương ứng.
            \item Hệ thống hiển thị tiến độ hoàn thành, thời gian học, kết quả gần đây và nhiệm vụ liên quan.
            \item Nếu chức năng gửi báo cáo được kích hoạt, dữ liệu được chuyển tới Notification Service theo cấu hình kênh.
            \item Yêu cầu với học sinh không có liên kết hợp lệ bị từ chối.
        \end{enumerate}

    \end{itemize}

    \item \textbf{Business rules:}

    BR-13, BR-14, BR-24, BR-56, BR-59.

    \ScreenLayoutItem

    \begin{figure}[H]
        \centering
        \SafeIncludeGraphics[width=0.94\textwidth,height=0.78\textheight,keepaspectratio]{figures/03_02/new_2026_09_10/reporting_hub_student_summary.png}
        \caption{Màn hình tổng hợp báo cáo tiến độ của học sinh dành cho Phụ huynh.}
        \label{fig:screen-parent-child-analysis}
    \end{figure}

    \ScreenLayoutFigure{figures/03_02/new_2026_09_10/parent_progress_report_config.png}{Configure or Request a Child Progress Report}

\end{itemize}

\subsubsection{View Child Examination Results and Predicted Score}
\label{subsubsec:parent-view-exam-prediction}

\begin{itemize}

    \item \textbf{Function trigger:}

    Phụ huynh chọn khu vực kết quả thi/dự đoán điểm của hồ sơ con.

    \item \textbf{Function description:}

    \begin{itemize}

        \item \textbf{Actors/Roles:}
        \begin{itemize}
            \item Phụ huynh.
        \end{itemize}

        \item \textbf{Purpose:}
        \begin{enumerate}
            \item Hiển thị kết quả bài kiểm tra và điểm dự đoán.
            \item Giúp theo dõi khả năng đạt mục tiêu.
        \end{enumerate}

        \item \textbf{Data processing:}
        \begin{enumerate}
            \item Hệ thống xác minh quan hệ Phụ huynh--Học sinh.
            \item Hệ thống tải lịch sử \texttt{BaiLam} và kết quả phân tích đã có.
            \item Hệ thống tải dữ liệu dự đoán điểm/xác suất đạt mục tiêu nếu mô hình đã tạo.
            \item Hệ thống chỉ hiển thị dữ liệu của hồ sơ liên kết.
            \item Phụ huynh không được sửa điểm, bài làm hoặc dữ liệu năng lực.
        \end{enumerate}

    \end{itemize}

    \item \textbf{Business rules:}

    BR-13, BR-42.

    \Needspace{0.60\textheight}
    \ScreenLayoutItem

    \ScreenLayoutFigure{figures/03_02/parent_child_analysis.png}{View Child Examination Results and Predicted Score}

\end{itemize}

\subsubsection{View Learning Support Recommendations}
\label{subsubsec:parent-view-recommendations}

\begin{itemize}

    \item \textbf{Function trigger:}

    Phụ huynh mở khu vực khuyến nghị hỗ trợ học tập của hồ sơ con.

    \item \textbf{Function description:}

    \begin{itemize}

        \item \textbf{Actors/Roles:}
        \begin{itemize}
            \item Phụ huynh.
        \end{itemize}

        \item \textbf{Purpose:}
        \begin{enumerate}
            \item Cung cấp khuyến nghị về thời gian, nội dung và phương pháp hỗ trợ.
            \item Tóm tắt điểm cần quan tâm từ dữ liệu tiến độ/kết quả.
        \end{enumerate}

        \item \textbf{Data processing:}
        \begin{enumerate}
            \item Hệ thống xác minh học sinh thuộc liên kết hợp lệ.
            \item Hệ thống lấy tiến độ, kết quả, phần kiến thức còn yếu và lịch học.
            \item Hệ thống tạo/lấy các khuyến nghị đã được hệ thống hoặc AI đề xuất.
            \item Khuyến nghị chỉ mang tính hỗ trợ, không cho phép Phụ huynh sửa dữ liệu học thuật.
            \item Nếu chưa đủ dữ liệu, hệ thống hiển thị trạng thái chưa có khuyến nghị thay vì suy đoán.
        \end{enumerate}

    \end{itemize}

    \item \textbf{Business rules:}

    BR-13, BR-24, BR-33.

    \ScreenLayoutItem

    \ScreenLayoutFigure{figures/03_02/parent_child_analysis.png}{View Learning Support Recommendations}

\end{itemize}

\subsubsection{View Teacher Evaluation and AI Analysis}
\label{subsubsec:parent-view-teacher-evaluation-ai-analysis}

\begin{itemize}

    \item \textbf{Function trigger:}

    Phụ huynh chọn chức năng xem đánh giá của Giáo viên và phân tích AI đối với hồ sơ con đang theo dõi.

    \item \textbf{Function description:}

    \begin{itemize}

        \item \textbf{Actors/Roles:}
        \begin{itemize}
            \item Phụ huynh.
        \end{itemize}

        \item \textbf{Purpose:}
        \begin{enumerate}
            \item Hiển thị nhận xét chuyên môn của Giáo viên về mức độ tiến bộ của học sinh.
            \item Hiển thị phân tích AI liên quan ở dạng thông tin hỗ trợ khi có dữ liệu.
        \end{enumerate}

        \item \textbf{Data processing:}
        \begin{enumerate}
            \item Hệ thống xác minh \texttt{LienKetPhuHuynh} giữa Phụ huynh và học sinh.
            \item Hệ thống tải các đánh giá/nhận xét do Giáo viên tạo trong phạm vi học sinh được chọn.
            \item Nếu có kết quả phân tích AI, hệ thống tải dữ liệu liên quan từ kết quả năng lực, Knowledge Tracing hoặc mô hình phân tích được triển khai.
            \item Giao diện phải phân biệt rõ nhận xét của Giáo viên và phân tích/đề xuất của AI.
            \item Phụ huynh chỉ có quyền đọc; không được chỉnh sửa nhận xét chuyên môn hoặc dữ liệu năng lực.
            \item Nếu chưa có đánh giá hoặc phân tích, hệ thống hiển thị trạng thái chưa có dữ liệu.
        \end{enumerate}

    \end{itemize}

    \item \textbf{Business rules:}

    BR-13, BR-42, BR-46.

    \ScreenLayoutItem

    \ScreenLayoutFigure{figures/03_02/parent_child_analysis.png}{View Teacher Evaluation and AI Analysis}

\end{itemize}


\subsubsection{Configure Child Learning Time Limits}
\label{subsubsec:configure-child-time-limit}

\begin{itemize}

    \item \textbf{Function trigger:}

    Phụ huynh chọn chức năng giới hạn thời gian cho một học sinh đã liên kết.

    \item \textbf{Function description:}

    \begin{itemize}

        \item \textbf{Actors/Roles:}
        \begin{itemize}
            \item Phụ huynh.
        \end{itemize}

        \item \textbf{Purpose:}
        \begin{enumerate}
            \item Thiết lập khung giờ và tổng thời gian tối đa.
            \item Duy trì giới hạn riêng cho từng học sinh.
        \end{enumerate}

        \item \textbf{Data processing:}
        \begin{enumerate}
            \item Hệ thống xác minh \texttt{LienKetPhuHuynh} thuộc Phụ huynh hiện tại.
            \item Hệ thống tải \texttt{CauHinhGioiHan} hiện có nếu đã thiết lập.
            \item Phụ huynh nhập khung giờ bắt đầu, kết thúc và tổng thời gian tối đa trong ngày.
            \item Hệ thống kiểm tra định dạng và giá trị giới hạn hợp lệ.
            \item Mỗi liên kết chỉ có một cấu hình giới hạn theo quan hệ 1--1.
            \item Nếu hợp lệ, hệ thống tạo/cập nhật \texttt{CauHinhGioiHan}.
            \item Giới hạn chỉ áp dụng cho học sinh gắn với liên kết tương ứng.
        \end{enumerate}

    \end{itemize}

    \item \textbf{Business rules:}

    BR-13, BR-15, BR-16.

    \ScreenLayoutItem

    \ScreenLayoutFigure{figures/03_02/new_2026_09_10/parent_time_limit.png}{Configure Child Learning Time Limits}

\end{itemize}

\subsubsection{View Child Schedule, Recent Activity and Attendance}
\label{subsubsec:parent-view-schedule-attendance}

\begin{itemize}

    \item \textbf{Function trigger:}

    Phụ huynh chọn ``Lịch học/Hoạt động/Chuyên cần'' của hồ sơ con.

    \item \textbf{Function description:}

    \begin{itemize}

        \item \textbf{Actors/Roles:}
        \begin{itemize}
            \item Phụ huynh.
        \end{itemize}

        \item \textbf{Purpose:}
        \begin{enumerate}
            \item Hiển thị lịch học, nhiệm vụ gần đây và chuyên cần.
            \item Giúp theo dõi mức độ tham gia học tập.
        \end{enumerate}

        \item \textbf{Data processing:}
        \begin{enumerate}
            \item Hệ thống xác minh quan hệ Phụ huynh--Học sinh.
            \item Hệ thống lấy lớp đang tham gia, các \texttt{BuoiHoc}, nhiệm vụ \texttt{GiaoBaiTap} và \texttt{DiemDanh}.
            \item Hệ thống tổng hợp hoạt động gần đây từ tiến độ, bài làm và lịch học theo phạm vi triển khai.
            \item Chuyên cần được hiển thị theo từng buổi và tổng hợp theo thời gian.
            \item Phụ huynh chỉ có quyền đọc, không được sửa điểm danh.
            \item Nếu không có liên kết hợp lệ, hệ thống từ chối truy vấn.
        \end{enumerate}

    \end{itemize}

    \item \textbf{Business rules:}

    BR-13, BR-63, BR-66.

    \ScreenLayoutItem

    \ScreenLayoutFigure{figures/03_02/parent_dashboard.png}{View Child Schedule, Recent Activity and Attendance}

\end{itemize}

\subsubsection{View Performance Decline Alerts}
\label{subsubsec:parent-view-decline-alerts}

\begin{itemize}

    \item \textbf{Function trigger:}

    Phụ huynh nhận hoặc mở cảnh báo suy giảm kết quả của học sinh.

    \item \textbf{Function description:}

    \begin{itemize}

        \item \textbf{Actors/Roles:}
        \begin{itemize}
            \item Phụ huynh.
            \item Notification Service.
        \end{itemize}

        \item \textbf{Purpose:}
        \begin{enumerate}
            \item Cảnh báo khi kết quả, tần suất học hoặc mức hoàn thành suy giảm theo tiêu chí.
            \item Điều hướng tới dữ liệu chi tiết liên quan.
        \end{enumerate}

        \item \textbf{Data processing:}
        \begin{enumerate}
            \item Hệ thống đánh giá tiến độ/kết quả theo tiêu chí cảnh báo được cấu hình.
            \item Khi thỏa điều kiện, hệ thống tạo \texttt{ThongBao} cho Phụ huynh liên kết hợp lệ.
            \item Hệ thống tạo \texttt{ThongBaoNguoiNhan} theo từng người nhận/kênh gửi.
            \item Thông báo có thể chứa \texttt{DuongDanDieuHuong} tới báo cáo học sinh.
            \item Khi mở cảnh báo, hệ thống kiểm tra lại quyền hồ sơ con.
            \item Trạng thái gửi/đọc được cập nhật theo bản ghi người nhận.
        \end{enumerate}

    \end{itemize}

    \item \textbf{Business rules:}

    BR-13, BR-56, BR-58, BR-70, BR-71.

    \ScreenLayoutItem

    \ScreenLayoutFigure{figures/03_02/parent_dashboard.png}{View Performance Decline Alerts}

\end{itemize}


% =========================================================
% PARENT-TEACHER COMMUNICATION
% =========================================================

\subsection{Parent-Teacher Communication}
\label{subsec:parent-teacher-communication}

Phân hệ trao đổi cho phép Phụ huynh và Giáo viên giao tiếp về một học sinh cụ thể.

\subsubsection{View or Start a Parent-Teacher Conversation}
\label{subsubsec:start-parent-teacher-conversation}

\begin{itemize}

    \item \textbf{Function trigger:}

    Phụ huynh chọn ``Trao đổi với Giáo viên'' hoặc Giáo viên mở danh sách hội thoại liên quan.

    \item \textbf{Function description:}

    \begin{itemize}

        \item \textbf{Actors/Roles:}
        \begin{itemize}
            \item Phụ huynh.
            \item Giáo viên / Giảng viên.
        \end{itemize}

        \item \textbf{Purpose:}
        \begin{enumerate}
            \item Tạo hoặc mở hội thoại về một học sinh cụ thể.
            \item Đảm bảo chỉ đúng Phụ huynh và Giáo viên liên quan tham gia.
        \end{enumerate}

        \item \textbf{Data processing:}
        \begin{enumerate}
            \item Nếu Phụ huynh khởi tạo, hệ thống xác minh \texttt{LienKetPhuHuynh} với học sinh.
            \item Hệ thống xác định lớp của học sinh và giáo viên phụ trách hoặc danh sách giáo viên hợp lệ.
            \item Khi tạo mới, hệ thống tạo \texttt{HoiThoaiTraoDoi} với \texttt{MaPhuHuynh}, \texttt{MaGiaoVien}, \texttt{MaHocSinh}, tiêu đề và trạng thái.
            \item Khi mở hội thoại cũ, hệ thống kiểm tra người dùng là thành viên.
            \item Người không phải thành viên không được xem nội dung.
            \item Danh sách hội thoại được sắp theo hoạt động gần nhất theo dữ liệu triển khai.
        \end{enumerate}

    \end{itemize}

    \item \textbf{Business rules:}

    BR-12, BR-13, BR-67, BR-68.

    \ScreenLayoutItem

    \ScreenLayoutFigure{figures/03_02/new_2026_09_10/parent_teacher_conversation_compose.png}{View or Start a Parent-Teacher Conversation}

\end{itemize}

\subsubsection{Send a Parent-Teacher Message}
\label{subsubsec:send-parent-teacher-message}

\begin{itemize}

    \item \textbf{Function trigger:}

    Phụ huynh hoặc Giáo viên nhập nội dung và chọn gửi trong hội thoại hợp lệ.

    \item \textbf{Function description:}

    \begin{itemize}

        \item \textbf{Actors/Roles:}
        \begin{itemize}
            \item Phụ huynh.
            \item Giáo viên / Giảng viên.
        \end{itemize}

        \item \textbf{Purpose:}
        \begin{enumerate}
            \item Gửi và lưu tin nhắn trong hội thoại.
            \item Thông báo cho thành viên còn lại về tin nhắn mới.
        \end{enumerate}

        \item \textbf{Data processing:}
        \begin{enumerate}
            \item Hệ thống xác minh người gửi là \texttt{MaPhuHuynh} hoặc \texttt{MaGiaoVien} của \texttt{HoiThoaiTraoDoi}.
            \item Hệ thống kiểm tra nội dung tin nhắn hợp lệ.
            \item Hệ thống tạo \texttt{TinNhanTraoDoi} với hội thoại, người gửi, nội dung, thời gian và trạng thái đọc ban đầu.
            \item Tin nhắn được hiển thị trong lịch sử theo thời gian.
            \item Hệ thống có thể tạo thông báo cho người nhận theo cấu hình.
            \item Nếu người gửi không thuộc hội thoại, yêu cầu bị từ chối.
        \end{enumerate}

    \end{itemize}

    \item \textbf{Business rules:}

    BR-67, BR-68, BR-56, BR-59.

    \ScreenLayoutItem

    \ScreenLayoutFigure{figures/03_02/new_2026_09_10/parent_teacher_conversation_compose.png}{Send a Parent-Teacher Message}

\end{itemize}

\subsubsection{View Conversation History and Read Status}
\label{subsubsec:view-conversation-history}

\begin{itemize}

    \item \textbf{Function trigger:}

    Một thành viên mở hội thoại để xem tin nhắn trước đó.

    \item \textbf{Function description:}

    \begin{itemize}

        \item \textbf{Actors/Roles:}
        \begin{itemize}
            \item Phụ huynh.
            \item Giáo viên / Giảng viên.
        \end{itemize}

        \item \textbf{Purpose:}
        \begin{enumerate}
            \item Hiển thị lịch sử trao đổi theo thời gian.
            \item Cập nhật trạng thái đọc khi phù hợp.
        \end{enumerate}

        \item \textbf{Data processing:}
        \begin{enumerate}
            \item Hệ thống kiểm tra người dùng thuộc hội thoại.
            \item Hệ thống tải \texttt{TinNhanTraoDoi} theo thời gian gửi.
            \item Hệ thống hiển thị rõ người gửi và nội dung.
            \item Tin nhắn nhận được có thể được cập nhật đã đọc khi mở hội thoại.
            \item Lịch sử chỉ được truy cập bởi thành viên hợp lệ.
            \item Quyền xem của Quản trị viên chỉ áp dụng nếu chính sách quản trị/bảo mật cho phép; không mặc định coi Admin là thành viên.
        \end{enumerate}

    \end{itemize}

    \item \textbf{Business rules:}

    BR-67, BR-68.

    \ScreenLayoutItem

    \ScreenLayoutFigure{figures/03_02/new_2026_09_10/parent_teacher_conversation_history.png}{View Conversation History and Read Status}

\end{itemize}


% =========================================================
% MANAGE ACCOUNTS & AUTHORIZATION
% =========================================================

\subsection{Manage Accounts \& Authorization}
\label{subsec:manage-accounts-authorization}

Phân hệ Quản trị viên quản lý tài khoản, trạng thái và vai trò truy cập toàn hệ thống.

\subsubsection{View User Account List}
\label{subsubsec:view-user-account-list}

\begin{itemize}

    \item \textbf{Function trigger:}

    Quản trị viên học vụ chọn ``Quản lý tài khoản''.

    \item \textbf{Function description:}

    \begin{itemize}

        \item \textbf{Actors/Roles:}
        \begin{itemize}
            \item Quản trị viên học vụ.
        \end{itemize}

        \item \textbf{Purpose:}
        \begin{enumerate}
            \item Hiển thị danh sách tài khoản toàn hệ thống.
            \item Hỗ trợ tìm kiếm/lọc và truy cập chức năng cập nhật trạng thái/vai trò.
        \end{enumerate}

        \item \textbf{Data processing:}
        \begin{enumerate}
            \item Hệ thống xác minh người dùng có vai trò Quản trị viên.
            \item Hệ thống tải danh sách \texttt{NguoiDung}.
            \item Danh sách hiển thị họ tên, Email, số điện thoại, vai trò, trạng thái và ngày tạo.
            \item Quản trị viên có thể tìm kiếm/lọc theo tiêu chí được triển khai.
            \item Từ danh sách, Quản trị viên chọn tài khoản để xem/cập nhật hoặc khóa/mở.
        \end{enumerate}

    \end{itemize}

    \item \textbf{Business rules:}

    BR-06, BR-48.

    \ScreenLayoutItem

    \begin{figure}[H]
        \centering
        \SafeIncludeGraphics[width=0.94\textwidth,height=0.78\textheight,keepaspectratio]{figures/03_02/new_2026_09_10/admin_user_list.png}
        \caption{Màn hình danh sách tài khoản, vai trò, trạng thái và thao tác khóa của Quản trị viên.}
        \label{fig:screen-admin-users}
    \end{figure}

\end{itemize}

\subsubsection{Create or Update a User Account}
\label{subsubsec:create-update-user-account}

\begin{itemize}

    \item \textbf{Function trigger:}

    Quản trị viên chọn tạo tài khoản mới hoặc chỉnh sửa tài khoản hiện có.

    \item \textbf{Function description:}

    \begin{itemize}

        \item \textbf{Actors/Roles:}
        \begin{itemize}
            \item Quản trị viên học vụ.
        \end{itemize}

        \item \textbf{Purpose:}
        \begin{enumerate}
            \item Tạo/cập nhật tài khoản toàn hệ thống.
            \item Đảm bảo dữ liệu định danh và vai trò hợp lệ.
        \end{enumerate}

        \item \textbf{Data processing:}
        \begin{enumerate}
            \item Hệ thống hiển thị biểu mẫu hoặc tải dữ liệu hiện tại.
            \item Quản trị viên nhập/cập nhật họ tên, Email, số điện thoại, vai trò và trường quản trị được phép.
            \item Hệ thống kiểm tra trường bắt buộc, định dạng và tính duy nhất của Email/Số điện thoại.
            \item Khi tạo mới, hệ thống tạo \texttt{NguoiDung} với trạng thái và ngày tạo phù hợp.
            \item Khi cập nhật, hệ thống duy trì các quan hệ dữ liệu hiện có.
            \item Nếu dữ liệu không hợp lệ, hệ thống không lưu.
        \end{enumerate}

    \end{itemize}

    \item \textbf{Business rules:}

    BR-06, BR-48.

    \ScreenLayoutItem

    \ScreenLayoutFigure{figures/03_02/new_2026_09_10/admin_create_user.png}{Create or Update a User Account}

\end{itemize}

\subsubsection{Lock or Unlock a User Account}
\label{subsubsec:lock-unlock-account}

\begin{itemize}

    \item \textbf{Function trigger:}

    Quản trị viên chọn khóa hoặc mở khóa một tài khoản.

    \item \textbf{Function description:}

    \begin{itemize}

        \item \textbf{Actors/Roles:}
        \begin{itemize}
            \item Quản trị viên học vụ.
        \end{itemize}

        \item \textbf{Purpose:}
        \begin{enumerate}
            \item Ngăn hoặc khôi phục khả năng đăng nhập.
            \item Duy trì trạng thái theo chính sách.
        \end{enumerate}

        \item \textbf{Data processing:}
        \begin{enumerate}
            \item Hệ thống xác minh quyền Quản trị viên.
            \item Hệ thống tải tài khoản mục tiêu và trạng thái hiện tại.
            \item Quản trị viên xác nhận thao tác.
            \item Hệ thống cập nhật \texttt{TrangThai} của \texttt{NguoiDung}.
            \item Tài khoản bị khóa không thể tạo phiên đăng nhập mới.
            \item Nếu tài khoản đang có phiên hoạt động, việc vô hiệu hóa phiên xử lý theo chính sách bảo mật.
        \end{enumerate}

    \end{itemize}

    \item \textbf{Business rules:}

    BR-02, BR-06.

    \ScreenLayoutItem

    \ScreenLayoutFigure{figures/03_02/new_2026_09_10/admin_user_role_actions.png}{Lock or Unlock a User Account}

\end{itemize}

\subsubsection{Assign User Role and Access Scope}
\label{subsubsec:assign-user-role}

\begin{itemize}

    \item \textbf{Function trigger:}

    Quản trị viên chọn chức năng phân quyền/vai trò của tài khoản.

    \item \textbf{Function description:}

    \begin{itemize}

        \item \textbf{Actors/Roles:}
        \begin{itemize}
            \item Quản trị viên học vụ.
        \end{itemize}

        \item \textbf{Purpose:}
        \begin{enumerate}
            \item Gán vai trò Học sinh, Giáo viên, Phụ huynh hoặc Quản trị viên.
            \item Đảm bảo quyền truy cập theo Role Authorization.
        \end{enumerate}

        \item \textbf{Data processing:}
        \begin{enumerate}
            \item Hệ thống tải \texttt{VaiTro} hiện tại.
            \item Quản trị viên chọn vai trò mới trong tập vai trò hỗ trợ.
            \item Hệ thống kiểm tra thay đổi hợp lệ và các quan hệ nghiệp vụ liên quan.
            \item Nếu hợp lệ, hệ thống cập nhật \texttt{VaiTro}.
            \item Các yêu cầu tiếp theo áp dụng quyền theo vai trò mới.
            \item Người dùng không phải Quản trị viên không được tự gán hoặc nâng quyền.
        \end{enumerate}

    \end{itemize}

    \item \textbf{Business rules:}

    BR-04, BR-06, BR-07.

    \ScreenLayoutItem

    \ScreenLayoutFigure{figures/03_02/new_2026_09_10/admin_user_role_actions.png}{Assign User Role and Access Scope}

\end{itemize}

% =========================================================
% SYSTEM & AI CONFIGURATION
% =========================================================

\subsection{System \& AI Configuration}
\label{subsec:system-ai-configuration}

Phân hệ Quản trị viên quản lý tham số/chính sách vận hành và cấu hình AI, đồng thời theo dõi lịch sử thay đổi cấu hình AI.

\subsubsection{View or Update System Configuration and Policies}
\label{subsubsec:system-configuration-policies}

\begin{itemize}

    \item \textbf{Function trigger:}

    Quản trị viên mở ``Cấu hình hệ thống và chính sách''.

    \item \textbf{Function description:}

    \begin{itemize}

        \item \textbf{Actors/Roles:}
        \begin{itemize}
            \item Quản trị viên học vụ.
        \end{itemize}

        \item \textbf{Purpose:}
        \begin{enumerate}
            \item Xem và điều chỉnh các tham số/chính sách vận hành thuộc phạm vi quản trị.
            \item Áp dụng nhất quán chính sách sử dụng và bảo mật.
        \end{enumerate}

        \item \textbf{Data processing:}
        \begin{enumerate}
            \item Hệ thống xác minh quyền Quản trị viên trước khi hiển thị cấu hình.
            \item Hệ thống tải các tham số/chính sách hiện đang áp dụng theo thành phần dữ liệu cấu hình được triển khai.
            \item Quản trị viên chỉnh sửa các giá trị được cho phép.
            \item Hệ thống kiểm tra kiểu dữ liệu, miền giá trị và điều kiện bắt buộc.
            \item Nếu hợp lệ, hệ thống lưu cấu hình mới và áp dụng cho các chức năng liên quan.
            \item Hệ thống không cho phép vai trò khác thay đổi cấu hình toàn hệ thống.
        \end{enumerate}

    \end{itemize}

    \item \textbf{Business rules:}

    BR-48.

    \ScreenLayoutItem

    \ScreenLayoutFigure{figures/03_02/new_2026_09_10/admin_system_policy.png}{View or Update System Configuration and Policies}

\end{itemize}

\subsubsection{View AI Configuration}
\label{subsubsec:view-ai-configuration}

\begin{itemize}

    \item \textbf{Function trigger:}

    Quản trị viên mở khu vực ``Cấu hình AI''.

    \item \textbf{Function description:}

    \begin{itemize}

        \item \textbf{Actors/Roles:}
        \begin{itemize}
            \item Quản trị viên học vụ.
        \end{itemize}

        \item \textbf{Purpose:}
        \begin{enumerate}
            \item Hiển thị các cấu hình AI đang tồn tại và cấu hình đang áp dụng.
            \item Cung cấp model, System Prompt và tham số RAG cho quản trị.
        \end{enumerate}

        \item \textbf{Data processing:}
        \begin{enumerate}
            \item Hệ thống tải các \texttt{CauHinhAI}.
            \item Hệ thống hiển thị tên cấu hình, System Prompt, tham số RAG, ModelVersion và trạng thái áp dụng.
            \item Quản trị viên chọn một cấu hình để xem chi tiết hoặc cập nhật.
            \item Người dùng không có quyền quản trị không được chỉnh cấu hình AI toàn hệ thống.
        \end{enumerate}

    \end{itemize}

    \item \textbf{Business rules:}

    BR-33, BR-48.

    \ScreenLayoutItem

    \begin{figure}[H]
        \centering
        \SafeIncludeGraphics[width=0.94\textwidth,height=0.78\textheight,keepaspectratio]{figures/03_02/new_2026_09_10/admin_ai_config.png}
        \caption{Màn hình cấu hình AI Tutor Provider và RAG toàn hệ thống.}
        \label{fig:screen-admin-ai-config}
    \end{figure}

\end{itemize}

\subsubsection{Update AI Rules and Configuration}
\label{subsubsec:update-ai-configuration}

\begin{itemize}

    \item \textbf{Function trigger:}

    Quản trị viên chọn chỉnh sửa cấu hình AI hoặc thiết lập quy tắc AI.

    \item \textbf{Function description:}

    \begin{itemize}

        \item \textbf{Actors/Roles:}
        \begin{itemize}
            \item Quản trị viên học vụ.
        \end{itemize}

        \item \textbf{Purpose:}
        \begin{enumerate}
            \item Cập nhật System Prompt, tham số RAG, phiên bản mô hình và trạng thái áp dụng.
            \item Kiểm soát quy tắc đề xuất/ngưỡng và hành vi AI.
        \end{enumerate}

        \item \textbf{Data processing:}
        \begin{enumerate}
            \item Hệ thống tải \texttt{CauHinhAI} hiện tại.
            \item Quản trị viên thay đổi các trường được phép.
            \item Hệ thống kiểm tra giá trị bắt buộc và tính hợp lệ của tham số.
            \item Trước khi cập nhật, hệ thống ghi nhận giá trị cũ để audit.
            \item Hệ thống lưu giá trị mới vào \texttt{CauHinhAI}.
            \item Hệ thống tạo \texttt{LogCauHinhAI} với \texttt{MaAdmin}, hành động, giá trị cũ, giá trị mới và thời gian.
            \item Nếu thay đổi không hợp lệ, cấu hình cũ được giữ nguyên.
        \end{enumerate}

    \end{itemize}

    \item \textbf{Business rules:}

    BR-33, BR-34, BR-48.

    \ScreenLayoutItem

    \ScreenLayoutFigure{figures/03_02/new_2026_09_10/admin_ai_config.png}{Update AI Rules and Configuration}

\end{itemize}

\subsubsection{View AI Configuration Logs}
\label{subsubsec:view-ai-config-logs}

\begin{itemize}

    \item \textbf{Function trigger:}

    Quản trị viên chọn lịch sử thay đổi cấu hình AI.

    \item \textbf{Function description:}

    \begin{itemize}

        \item \textbf{Actors/Roles:}
        \begin{itemize}
            \item Quản trị viên học vụ.
        \end{itemize}

        \item \textbf{Purpose:}
        \begin{enumerate}
            \item Theo dõi ai thay đổi cấu hình, thay đổi gì và khi nào.
            \item Hỗ trợ audit.
        \end{enumerate}

        \item \textbf{Data processing:}
        \begin{enumerate}
            \item Hệ thống tải \texttt{LogCauHinhAI} theo cấu hình hoặc khoảng thời gian.
            \item Hệ thống kết hợp với \texttt{NguoiDung} để xác định Quản trị viên thực hiện.
            \item Mỗi bản ghi hiển thị hành động, giá trị cũ, giá trị mới và thời gian.
            \item Log là dữ liệu audit; người dùng thông thường không được chỉnh sửa trực tiếp.
            \item Quản trị viên có thể lọc/tìm kiếm theo khả năng triển khai.
        \end{enumerate}

    \end{itemize}

    \item \textbf{Business rules:}

    BR-34, BR-48.

    \ScreenLayoutItem

    \ScreenLayoutFigure{figures/03_02/admin_ai_configuration.png}{View AI Configuration Logs}

\end{itemize}


% =========================================================
% ADMINISTRATION DASHBOARD & REPORTS
% =========================================================

\subsection{Administration Dashboard \& Reports}
\label{subsec:administration-dashboard-reports}

Phân hệ báo cáo quản trị tổng hợp số liệu người dùng, học tập, vận hành, bảo mật và hiệu suất AI.

\subsubsection{Use Reporting Hub, Filters and Data Visualization}
\label{subsubsec:use-reporting-hub}

\begin{itemize}

    \item \textbf{Function trigger:}

    Người dùng có quyền báo cáo chọn ``Báo cáo/Thống kê'' từ Dashboard hoặc thanh điều hướng.

    \item \textbf{Function description:}

    \begin{itemize}

        \item \textbf{Actors/Roles:}
        \begin{itemize}
            \item Học sinh / Sinh viên.
            \item Giáo viên / Giảng viên.
            \item Phụ huynh.
            \item Quản trị viên học vụ.
        \end{itemize}

        \item \textbf{Purpose:}
        \begin{enumerate}
            \item Cung cấp một điểm tập trung để xem các báo cáo phù hợp với vai trò.
            \item Cho phép lọc dữ liệu theo các tiêu chí như khoảng thời gian, học kỳ, lớp và môn học khi tiêu chí đó áp dụng.
            \item Trực quan hóa dữ liệu bằng bảng/biểu đồ theo Screen Flow Reporting Hub.
        \end{enumerate}

        \item \textbf{Data processing:}
        \begin{enumerate}
            \item Hệ thống xác định vai trò và phạm vi dữ liệu được phép xem.
            \item Học sinh chỉ truy cập báo cáo cá nhân; Giáo viên chỉ truy cập báo cáo lớp/phạm vi phụ trách; Phụ huynh chỉ truy cập dữ liệu con đã liên kết; Quản trị viên truy cập báo cáo toàn hệ thống theo quyền.
            \item Hệ thống hiển thị danh sách/nhóm báo cáo khả dụng cho vai trò hiện tại.
            \item Người dùng chọn bộ lọc như khoảng thời gian, học kỳ, lớp hoặc môn/chủ đề khi báo cáo hỗ trợ.
            \item Hệ thống kiểm tra lại phạm vi bộ lọc; việc thay đổi tham số không được mở rộng quyền truy cập.
            \item Hệ thống tổng hợp dữ liệu từ các nguồn học tập/vận hành tương ứng và hiển thị dưới dạng số liệu, bảng, biểu đồ cột, biểu đồ tròn hoặc dạng trực quan phù hợp.
            \item Nếu không có dữ liệu trong phạm vi lọc, hệ thống hiển thị trạng thái không có dữ liệu thay vì tạo số liệu giả.
        \end{enumerate}

    \end{itemize}

    \item \textbf{Business rules:}

    BR-07, BR-11, BR-13, BR-24, BR-48.

    \ScreenLayoutItem

    \begin{figure}[H]
        \centering
        \SafeIncludeGraphics[width=0.94\textwidth,height=0.78\textheight,keepaspectratio]{figures/03_02/admin_analytics.png}
        \caption{Màn hình Tổng quan Phân tích và AI Insights của hệ thống.}
        \label{fig:screen-admin-analytics}
    \end{figure}

\end{itemize}


\subsubsection{Customize and Export a Report}
\label{subsubsec:customize-export-report}

\begin{itemize}

    \item \textbf{Function trigger:}

    Từ Reporting Hub, người dùng chọn chức năng tùy chỉnh báo cáo hoặc xuất tệp.

    \item \textbf{Function description:}

    \begin{itemize}

        \item \textbf{Actors/Roles:}
        \begin{itemize}
            \item Học sinh / Sinh viên.
            \item Giáo viên / Giảng viên.
            \item Phụ huynh.
            \item Quản trị viên học vụ.
        \end{itemize}

        \item \textbf{Purpose:}
        \begin{enumerate}
            \item Cho phép chọn các trường dữ liệu cần hiển thị trong báo cáo.
            \item Xuất báo cáo theo định dạng PDF, Excel hoặc CSV khi được hệ thống hỗ trợ.
        \end{enumerate}

        \item \textbf{Data processing:}
        \begin{enumerate}
            \item Hệ thống duy trì phạm vi dữ liệu và bộ lọc đã được xác thực từ Reporting Hub.
            \item Người dùng chọn các trường/cột được phép đưa vào báo cáo.
            \item Hệ thống loại bỏ hoặc vô hiệu hóa trường dữ liệu nằm ngoài quyền của vai trò hiện tại.
            \item Người dùng chọn định dạng xuất như PDF, Excel hoặc CSV theo Screen Flow.
            \item Hệ thống tạo tệp từ chính tập dữ liệu đã lọc và các trường đã chọn.
            \item Tệp xuất phải tuân theo cùng phạm vi Role Authorization như dữ liệu trên màn hình.
            \item Nếu quá trình tạo tệp thất bại, hệ thống thông báo lỗi và không phát hành tệp không hoàn chỉnh.
        \end{enumerate}

    \end{itemize}

    \item \textbf{Business rules:}

    BR-07, BR-11, BR-13, BR-24, BR-48.

    \ScreenLayoutItem

    \ScreenLayoutFigure{figures/03_02/admin_analytics.png}{Customize and Export a Report}

\end{itemize}


\subsubsection{View Administration Dashboard}
\label{subsubsec:view-admin-dashboard}

\begin{itemize}

    \item \textbf{Function trigger:}

    Quản trị viên đăng nhập thành công hoặc chọn ``Dashboard tổng quan''.

    \item \textbf{Function description:}

    \begin{itemize}

        \item \textbf{Actors/Roles:}
        \begin{itemize}
            \item Quản trị viên học vụ.
        \end{itemize}

        \item \textbf{Purpose:}
        \begin{enumerate}
            \item Hiển thị số liệu tổng quan về người dùng, nội dung, học tập và dịch vụ AI.
            \item Cung cấp điểm truy cập nhanh tới báo cáo.
        \end{enumerate}

        \item \textbf{Data processing:}
        \begin{enumerate}
            \item Hệ thống xác minh quyền Quản trị viên.
            \item Hệ thống tổng hợp số liệu từ tài khoản, lớp học, lộ trình, bài làm, nội dung, thanh toán, thông báo và nguồn phù hợp.
            \item Các chỉ số được tính theo khoảng thời gian/bộ lọc được Dashboard hỗ trợ.
            \item Hệ thống hiển thị số liệu tổng hợp và cảnh báo vận hành nếu có.
            \item Dữ liệu chi tiết chỉ được mở khi người dùng có quyền tương ứng.
        \end{enumerate}

    \end{itemize}

    \item \textbf{Business rules:}

    BR-06, BR-48.

    \ScreenLayoutItem

    \begin{figure}[H]
        \centering
        \SafeIncludeGraphics[width=0.94\textwidth,height=0.78\textheight,keepaspectratio]{figures/03_02/admin_dashboard.png}
        \caption{Màn hình Dashboard tổng quan của Quản trị viên học vụ.}
        \label{fig:screen-admin-dashboard}
    \end{figure}

\end{itemize}

\subsubsection{View Trend and Learning Analytics Reports}
\label{subsubsec:view-trend-reports}

\begin{itemize}

    \item \textbf{Function trigger:}

    Quản trị viên hoặc Giáo viên chọn báo cáo phân tích phù hợp với phạm vi.

    \item \textbf{Function description:}

    \begin{itemize}

        \item \textbf{Actors/Roles:}
        \begin{itemize}
            \item Quản trị viên học vụ.
            \item Giáo viên / Giảng viên.
        \end{itemize}

        \item \textbf{Purpose:}
        \begin{enumerate}
            \item Phân tích xu hướng sử dụng, kết quả học tập và hiệu quả giảng dạy.
            \item Hỗ trợ đánh giá lớp học/nội dung ở mức tổng hợp.
        \end{enumerate}

        \item \textbf{Data processing:}
        \begin{enumerate}
            \item Hệ thống xác định phạm vi dữ liệu theo vai trò.
            \item Người dùng chọn bộ lọc như thời gian, lớp, môn/chủ đề theo chức năng triển khai.
            \item Hệ thống tổng hợp tiến độ, bài làm, kết quả và dữ liệu sử dụng nội dung.
            \item Giáo viên chỉ nhận báo cáo lớp/phạm vi phụ trách; Quản trị viên có thể xem toàn hệ thống.
            \item Hệ thống hiển thị số liệu, bảng và biểu đồ theo thiết kế giao diện sau này.
        \end{enumerate}

    \end{itemize}

    \item \textbf{Business rules:}

    BR-11, BR-24, BR-42.

    \ScreenLayoutItem

    \ScreenLayoutFigure{figures/03_02/admin_analytics.png}{View Trend and Learning Analytics Reports}

\end{itemize}

\subsubsection{View Security and Access Logs}
\label{subsubsec:view-security-access-logs}

\begin{itemize}

    \item \textbf{Function trigger:}

    Quản trị viên chọn ``Nhật ký truy cập và bảo mật''.

    \item \textbf{Function description:}

    \begin{itemize}

        \item \textbf{Actors/Roles:}
        \begin{itemize}
            \item Quản trị viên học vụ.
        \end{itemize}

        \item \textbf{Purpose:}
        \begin{enumerate}
            \item Kiểm tra lịch sử đăng nhập, hành động và sự kiện bảo mật.
            \item Hỗ trợ phát hiện truy cập bất thường.
        \end{enumerate}

        \item \textbf{Data processing:}
        \begin{enumerate}
            \item Hệ thống xác minh quyền truy cập báo cáo bảo mật.
            \item Hệ thống tải dữ liệu log truy cập/bảo mật từ thành phần lưu trữ log được triển khai.
            \item Quản trị viên có thể lọc theo tài khoản, loại sự kiện, thời gian hoặc trạng thái khi hỗ trợ.
            \item Hệ thống hiển thị thông tin audit nhưng không tiết lộ bí mật xác thực.
            \item Log bảo mật không được cho phép người dùng thông thường chỉnh sửa.
        \end{enumerate}

    \end{itemize}

    \item \textbf{Business rules:}

    BR-06, BR-48.

    \ScreenLayoutItem

    \ScreenLayoutFigure{figures/03_02/admin_analytics.png}{View Security and Access Logs}

\end{itemize}

\subsubsection{View AI Model Performance and Export Reports}
\label{subsubsec:view-ai-performance-export}

\begin{itemize}

    \item \textbf{Function trigger:}

    Quản trị viên chọn báo cáo hiệu suất AI hoặc chức năng xuất báo cáo.

    \item \textbf{Function description:}

    \begin{itemize}

        \item \textbf{Actors/Roles:}
        \begin{itemize}
            \item Quản trị viên học vụ.
        \end{itemize}

        \item \textbf{Purpose:}
        \begin{enumerate}
            \item Theo dõi hiệu suất chức năng/mô hình AI theo dữ liệu thu thập.
            \item Xuất báo cáo theo định dạng hệ thống hỗ trợ.
        \end{enumerate}

        \item \textbf{Data processing:}
        \begin{enumerate}
            \item Hệ thống tải các chỉ số hiệu suất AI được ghi nhận theo phiên bản mô hình/cấu hình.
            \item Quản trị viên chọn khoảng thời gian và phạm vi dữ liệu.
            \item Hệ thống tổng hợp chỉ số và hiển thị báo cáo.
            \item Khi chọn xuất, hệ thống tạo tệp theo định dạng được triển khai như PDF, Excel hoặc CSV theo Screen Flow báo cáo.
            \item Tệp xuất chỉ chứa dữ liệu trong phạm vi quyền.
            \item Nếu không đủ dữ liệu, hệ thống thể hiện trạng thái không có dữ liệu thay vì tự tạo số liệu.
        \end{enumerate}

    \end{itemize}

    \item \textbf{Business rules:}

    BR-33, BR-34, BR-48.

    \ScreenLayoutItem

    \ScreenLayoutFigure{figures/03_02/admin_analytics.png}{View AI Model Performance and Export Reports}

\end{itemize}


% =========================================================
% SERVICE PACKAGES & PAYMENT
% =========================================================

\subsection{Service Packages \& Payment}
\label{subsec:service-packages-payment}

Phân hệ thương mại cho phép người dùng xem gói dịch vụ, tạo đơn hàng, áp dụng mã giảm giá, thanh toán, xem lịch sử và hóa đơn điện tử.

\subsubsection{View Service Package List and Package Detail}
\label{subsubsec:view-service-packages}

\begin{itemize}

    \item \textbf{Function trigger:}

    Học sinh/Sinh viên hoặc Phụ huynh mở khu vực gói dịch vụ.

    \item \textbf{Function description:}

    \begin{itemize}

        \item \textbf{Actors/Roles:}
        \begin{itemize}
            \item Học sinh / Sinh viên.
            \item Phụ huynh.
        \end{itemize}

        \item \textbf{Purpose:}
        \begin{enumerate}
            \item Hiển thị các gói dịch vụ đang khả dụng.
            \item Cho phép xem quyền lợi, thời hạn và giá trước khi thanh toán.
        \end{enumerate}

        \item \textbf{Data processing:}
        \begin{enumerate}
            \item Hệ thống tải các \texttt{GoiDichVu} đang ở trạng thái cho phép sử dụng.
            \item Danh sách hiển thị tên gói, giá, thời hạn và mô tả quyền lợi.
            \item Người dùng chọn một gói để xem chi tiết.
            \item Hệ thống hiển thị đầy đủ thông tin gói trước khi chuyển sang Checkout.
            \item Gói không còn hoạt động không được phép tạo giao dịch mới.
        \end{enumerate}

    \end{itemize}

    \item \textbf{Business rules:}

    BR-49.

    \ScreenLayoutItem

    \begin{figure}[H]
        \centering
        \SafeIncludeGraphics[width=0.94\textwidth,height=0.78\textheight,keepaspectratio]{figures/03_02/payment_packages.png}
        \caption{Màn hình danh sách gói dịch vụ và thao tác Checkout.}
        \label{fig:screen-payment-invoice}
    \end{figure}

\end{itemize}

\subsubsection{Checkout and Apply Voucher}
\label{subsubsec:checkout-apply-voucher}

\begin{itemize}

    \item \textbf{Function trigger:}

    Người dùng chọn mua một hoặc nhiều gói và chuyển tới Checkout.

    \item \textbf{Function description:}

    \begin{itemize}

        \item \textbf{Actors/Roles:}
        \begin{itemize}
            \item Học sinh / Sinh viên.
            \item Phụ huynh.
        \end{itemize}

        \item \textbf{Purpose:}
        \begin{enumerate}
            \item Tạo thông tin đơn hàng trước thanh toán.
            \item Cho phép áp dụng mã giảm giá hợp lệ.
        \end{enumerate}

        \item \textbf{Data processing:}
        \begin{enumerate}
            \item Hệ thống xác định người dùng và các gói đã chọn.
            \item Hệ thống tính số lượng/đơn giá hiện tại để chuẩn bị chi tiết đơn hàng.
            \item Nếu nhập mã giảm giá, hệ thống tìm bản ghi \texttt{MaGiamGia} theo \texttt{CodeVoucher}, xác định khóa chính \texttt{MaVoucher}, sau đó kiểm tra thời hạn và điều kiện áp dụng.
            \item Nếu voucher hợp lệ, hệ thống tính giá trị giảm; voucher không hợp lệ không được áp dụng.
            \item Hệ thống hiển thị tổng tiền cuối cùng trước khi xác nhận.
            \item Khi tạo đơn hàng, \texttt{ChiTietDonHang} phải lưu \texttt{SoLuong} và \texttt{DonGiaAtTime}.
        \end{enumerate}

    \end{itemize}

    \item \textbf{Business rules:}

    BR-49, BR-50, BR-51, BR-52.

    \ScreenLayoutItem

    \begin{figure}[H]
        \centering
        \SafeIncludeGraphics[width=0.94\textwidth,height=0.78\textheight,keepaspectratio]{figures/03_02/payment_checkout.png}
        \caption{Màn hình Checkout, áp dụng mã giảm giá và lựa chọn phương thức thanh toán.}
        \label{fig:screen-payment-checkout}
    \end{figure}

\end{itemize}

\subsubsection{Select Payment Method and Confirm Order}
\label{subsubsec:select-payment-method}

\begin{itemize}

    \item \textbf{Function trigger:}

    Người dùng tiếp tục từ Checkout sau khi kiểm tra đơn hàng.

    \item \textbf{Function description:}

    \begin{itemize}

        \item \textbf{Actors/Roles:}
        \begin{itemize}
            \item Học sinh / Sinh viên.
            \item Phụ huynh.
        \end{itemize}

        \item \textbf{Purpose:}
        \begin{enumerate}
            \item Cho phép chọn phương thức thanh toán.
            \item Yêu cầu xác nhận đơn hàng/điều khoản trước khi gửi cổng thanh toán.
        \end{enumerate}

        \item \textbf{Data processing:}
        \begin{enumerate}
            \item Hệ thống hiển thị tổng tiền, gói và voucher đã áp dụng.
            \item Người dùng chọn phương thức thanh toán.
            \item Người dùng xác nhận đơn hàng và điều khoản cần thiết theo Screen Flow.
            \item Hệ thống tạo/cập nhật \texttt{DonHangThanhToan} với người dùng, voucher nếu có, tổng tiền, phương thức và \texttt{TrangThaiGiaoDich = ChoXuLy}.
            \item Hệ thống tạo các \texttt{ChiTietDonHang}.
            \item Hệ thống chuyển yêu cầu sang quy trình/cổng thanh toán.
        \end{enumerate}

    \end{itemize}

    \item \textbf{Business rules:}

    BR-50, BR-51, BR-52, BR-54.

    \ScreenLayoutItem

    \ScreenLayoutFigure{figures/03_02/payment_checkout.png}{Select Payment Method and Confirm Order}

\end{itemize}

\subsubsection{Process and Confirm Payment}
\label{subsubsec:process-confirm-payment}

\begin{itemize}

    \item \textbf{Function trigger:}

    Hệ thống gửi yêu cầu thanh toán và nhận kết quả giao dịch.

    \item \textbf{Function description:}

    \begin{itemize}

        \item \textbf{Actors/Roles:}
        \begin{itemize}
            \item Payment Service.
            \item Học sinh / Sinh viên.
            \item Phụ huynh.
        \end{itemize}

        \item \textbf{Purpose:}
        \begin{enumerate}
            \item Xác nhận trạng thái giao dịch từ cổng thanh toán.
            \item Chỉ kích hoạt quyền lợi khi giao dịch được xác nhận thành công.
        \end{enumerate}

        \item \textbf{Data processing:}
        \begin{enumerate}
            \item Hệ thống gửi mã đơn hàng, số tiền và thông tin cần thiết tới cổng thanh toán.
            \item Người dùng thực hiện bước xác thực như QR/OTP/mật khẩu theo phương thức của cổng.
            \item Hệ thống nhận kết quả và đối chiếu đơn hàng.
            \item Chỉ khi xác nhận thành công hợp lệ, hệ thống cập nhật \texttt{TrangThaiGiaoDich = ThanhCong}.
            \item Nếu thất bại, trạng thái được cập nhật Thất bại và không kích hoạt quyền lợi.
            \item Giao dịch chưa có kết quả cuối cùng giữ trạng thái Chờ xử lý.
            \item Sau khi thành công, hệ thống kích hoạt quyền lợi/gói theo cơ chế triển khai và chuyển sang phát hành hóa đơn.
        \end{enumerate}

    \end{itemize}

    \item \textbf{Business rules:}

    BR-53, BR-54.

    \ScreenLayoutItem

    \ScreenLayoutFigure{figures/03_02/payment_checkout.png}{Process and Confirm Payment}

\end{itemize}

\subsubsection{View Payment Result and Transaction History}
\label{subsubsec:view-payment-history}

\begin{itemize}

    \item \textbf{Function trigger:}

    Người dùng hoàn tất thanh toán hoặc mở ``Lịch sử giao dịch''.

    \item \textbf{Function description:}

    \begin{itemize}

        \item \textbf{Actors/Roles:}
        \begin{itemize}
            \item Học sinh / Sinh viên.
            \item Phụ huynh.
        \end{itemize}

        \item \textbf{Purpose:}
        \begin{enumerate}
            \item Hiển thị kết quả giao dịch hiện tại.
            \item Cho phép xem lại giao dịch của chính tài khoản.
        \end{enumerate}

        \item \textbf{Data processing:}
        \begin{enumerate}
            \item Sau thanh toán, hệ thống hiển thị mã đơn hàng, tổng tiền, phương thức và trạng thái.
            \item Khi mở lịch sử, hệ thống truy vấn \texttt{DonHangThanhToan} theo người dùng hiện tại.
            \item Người dùng chỉ xem giao dịch của chính mình.
            \item Mỗi giao dịch có thể mở chi tiết để xem \texttt{ChiTietDonHang}.
            \item Trạng thái hiển thị tối thiểu gồm Chờ xử lý, Thành công và Thất bại.
        \end{enumerate}

    \end{itemize}

    \item \textbf{Business rules:}

    BR-50, BR-54.

    \ScreenLayoutItem

    \ScreenLayoutFigure{figures/03_02/parent_payment_invoice.png}{View Payment Result and Transaction History}

\end{itemize}

\subsubsection{Access Purchased Service Benefits}
\label{subsubsec:access-purchased-service-benefits}

\begin{itemize}

    \item \textbf{Function trigger:}

    Sau khi một giao dịch được xác nhận thành công, người dùng chọn truy cập nội dung hoặc quyền lợi thuộc gói dịch vụ đã mua.

    \item \textbf{Function description:}

    \begin{itemize}

        \item \textbf{Actors/Roles:}
        \begin{itemize}
            \item Học sinh / Sinh viên.
            \item Phụ huynh.
        \end{itemize}

        \item \textbf{Purpose:}
        \begin{enumerate}
            \item Cho phép sử dụng các quyền lợi đã được kích hoạt từ gói dịch vụ thanh toán thành công.
            \item Ngăn truy cập quyền lợi khi giao dịch chưa thành công hoặc gói không còn hiệu lực.
        \end{enumerate}

        \item \textbf{Data processing:}
        \begin{enumerate}
            \item Hệ thống xác minh người dùng hiện tại là chủ của đơn hàng/gói đã mua.
            \item Hệ thống kiểm tra \texttt{DonHangThanhToan.TrangThaiGiaoDich = ThanhCong}.
            \item Hệ thống xác định các \texttt{GoiDichVu} trong \texttt{ChiTietDonHang} và quyền lợi tương ứng theo cấu hình sản phẩm được triển khai.
            \item Nếu quyền lợi đã được kích hoạt và còn hiệu lực, hệ thống cho phép điều hướng tới nội dung/chức năng tương ứng.
            \item Nếu giao dịch đang Chờ xử lý/Thất bại hoặc quyền lợi hết hạn, hệ thống không cho phép truy cập theo gói đó.
            \item Việc cấp quyền lợi phải được lưu trong cơ chế entitlement/subscription được triển khai; nếu ERD hiện tại chưa có bảng tương ứng, thiết kế dữ liệu cần được bổ sung trước khi hiện thực đầy đủ.
            \item Người dùng không thể sử dụng mã đơn hàng của tài khoản khác để nhận quyền truy cập.
        \end{enumerate}

    \end{itemize}

    \item \textbf{Business rules:}

    BR-49, BR-53, BR-54.

    \ScreenLayoutItem

    \ScreenLayoutFigure{figures/03_02/payment_packages.png}{Access Purchased Service Benefits}

\end{itemize}


\subsubsection{View or Download Electronic Invoice}
\label{subsubsec:view-download-invoice}

\begin{itemize}

    \item \textbf{Function trigger:}

    Người dùng mở một đơn hàng thanh toán thành công và chọn hóa đơn điện tử.

    \item \textbf{Function description:}

    \begin{itemize}

        \item \textbf{Actors/Roles:}
        \begin{itemize}
            \item Học sinh / Sinh viên.
            \item Phụ huynh.
        \end{itemize}

        \item \textbf{Purpose:}
        \begin{enumerate}
            \item Hiển thị và cho phép tải hóa đơn của giao dịch thành công.
            \item Bảo đảm hóa đơn thuộc đúng chủ đơn hàng.
        \end{enumerate}

        \item \textbf{Data processing:}
        \begin{enumerate}
            \item Hệ thống kiểm tra \texttt{DonHangThanhToan.MaNguoiDung} trùng người dùng hiện tại.
            \item Hệ thống kiểm tra giao dịch đã Thành công.
            \item Nếu chưa có hóa đơn và quy trình phát hành hợp lệ, hệ thống tạo \texttt{HoaDonDienTu} gắn \texttt{MaDonHang}.
            \item Hóa đơn lưu số hóa đơn, ngày phát hành, tổng tiền, mã giao dịch, trạng thái và đường dẫn tệp theo thiết kế đã chốt.
            \item Hệ thống hiển thị chi tiết hóa đơn và cung cấp tệp/đường dẫn tải nếu có.
            \item Người dùng không thể thay đổi mã hóa đơn trong request để tải hóa đơn của tài khoản khác.
        \end{enumerate}

    \end{itemize}

    \item \textbf{Business rules:}

    BR-53, BR-69.

    \ScreenLayoutItem

    \ScreenLayoutFigure{figures/03_02/parent_payment_invoice.png}{View or Download Electronic Invoice}

\end{itemize}

\subsubsection{Request a Refund}
\label{subsubsec:request-refund}

\begin{itemize}

    \item \textbf{Function trigger:}

    Người dùng chọn yêu cầu hoàn tiền cho một giao dịch đủ điều kiện.

    \item \textbf{Function description:}

    \begin{itemize}

        \item \textbf{Actors/Roles:}
        \begin{itemize}
            \item Học sinh / Sinh viên.
            \item Phụ huynh.
            \item Payment Service.
        \end{itemize}

        \item \textbf{Purpose:}
        \begin{enumerate}
            \item Ghi nhận yêu cầu hoàn tiền theo chính sách.
            \item Theo dõi kết quả xử lý hoàn tiền.
        \end{enumerate}

        \item \textbf{Data processing:}
        \begin{enumerate}
            \item Hệ thống kiểm tra đơn hàng thuộc người dùng và đã thanh toán thành công.
            \item Hệ thống kiểm tra điều kiện hoàn tiền theo chính sách hiện hành.
            \item Người dùng nhập lý do/thông tin cần thiết theo quy trình triển khai.
            \item Nếu đủ điều kiện, yêu cầu được chuyển tới thành phần thanh toán/quản trị để xử lý.
            \item Kết quả hoàn tiền được phản ánh trong dữ liệu giao dịch theo mô hình trạng thái được mở rộng khi triển khai.
            \item Nếu không đủ điều kiện, hệ thống từ chối và hiển thị lý do phù hợp.
        \end{enumerate}

    \end{itemize}

    \item \textbf{Business rules:}

    BR-55.

    \ScreenLayoutItem

    \ScreenLayoutFigure{figures/03_02/parent_payment_invoice.png}{Request a Refund}

\end{itemize}


% =========================================================
% NOTIFICATION MANAGEMENT
% =========================================================

\subsection{Notification Management}
\label{subsec:notification-management}

Phân hệ thông báo cung cấp Notification Center, cấu hình kênh nhận, gửi/lập lịch và theo dõi trạng thái theo từng người nhận.

\subsubsection{View Notification Center}
\label{subsubsec:view-notification-center}

\begin{itemize}

    \item \textbf{Function trigger:}

    Người dùng đăng nhập chọn ``Thông báo'' từ Dashboard hoặc thanh điều hướng.

    \item \textbf{Function description:}

    \begin{itemize}

        \item \textbf{Actors/Roles:}
        \begin{itemize}
            \item Học sinh / Sinh viên.
            \item Giáo viên / Giảng viên.
            \item Phụ huynh.
            \item Quản trị viên học vụ.
        \end{itemize}

        \item \textbf{Purpose:}
        \begin{enumerate}
            \item Hiển thị thông báo dành riêng cho người dùng.
            \item Phân biệt đã đọc/chưa đọc.
        \end{enumerate}

        \item \textbf{Data processing:}
        \begin{enumerate}
            \item Hệ thống lấy \texttt{MaNguoiDung} từ phiên hiện tại.
            \item Hệ thống truy vấn \texttt{ThongBaoNguoiNhan} theo \texttt{MaNguoiNhan}.
            \item Hệ thống kết hợp \texttt{ThongBao} để hiển thị tiêu đề, nội dung tóm tắt, thời gian và đường dẫn.
            \item Trạng thái đọc lấy từ bản ghi người nhận, không dùng chung cho tất cả người nhận.
            \item Người dùng không được xem thông báo không có bản ghi người nhận tương ứng.
            \item Danh sách có thể phân loại theo lịch học, điểm số, hệ thống hoặc nhóm Screen Flow hỗ trợ.
        \end{enumerate}

    \end{itemize}

    \item \textbf{Business rules:}

    BR-56, BR-58, BR-71.

    \ScreenLayoutItem

    \ScreenLayoutFigure{figures/03_02/new_2026_09_10/notification_inbox.png}{View Notification Center}

\end{itemize}

\subsubsection{View Notification Detail and Mark as Read}
\label{subsubsec:view-notification-detail}

\begin{itemize}

    \item \textbf{Function trigger:}

    Người dùng chọn một thông báo trong Notification Center.

    \item \textbf{Function description:}

    \begin{itemize}

        \item \textbf{Actors/Roles:}
        \begin{itemize}
            \item Học sinh / Sinh viên.
            \item Giáo viên / Giảng viên.
            \item Phụ huynh.
            \item Quản trị viên học vụ.
        \end{itemize}

        \item \textbf{Purpose:}
        \begin{enumerate}
            \item Hiển thị đầy đủ nội dung.
            \item Cập nhật trạng thái đọc của chính người nhận.
        \end{enumerate}

        \item \textbf{Data processing:}
        \begin{enumerate}
            \item Hệ thống xác minh \texttt{ThongBaoNguoiNhan} thuộc người dùng hiện tại.
            \item Hệ thống tải nội dung \texttt{ThongBao}.
            \item Khi mở, hệ thống cập nhật \texttt{TrangThaiDoc = DaDoc} và \texttt{ThoiGianDoc} nếu chưa ghi nhận.
            \item Nếu có \texttt{DuongDanDieuHuong}, người dùng có thể chuyển tới chức năng liên quan.
            \item Trang đích vẫn kiểm tra Role Authorization và phạm vi dữ liệu.
        \end{enumerate}

    \end{itemize}

    \item \textbf{Business rules:}

    BR-58, BR-71, BR-07.

    \ScreenLayoutItem

    \ScreenLayoutFigure{figures/03_02/new_2026_09_10/notification_detail.png}{View Notification Detail and Mark as Read}

\end{itemize}

\subsubsection{Configure Notification Preferences}
\label{subsubsec:configure-notification-preferences}

\begin{itemize}

    \item \textbf{Function trigger:}

    Người dùng mở cài đặt thông báo của tài khoản.

    \item \textbf{Function description:}

    \begin{itemize}

        \item \textbf{Actors/Roles:}
        \begin{itemize}
            \item Học sinh / Sinh viên.
            \item Giáo viên / Giảng viên.
            \item Phụ huynh.
            \item Quản trị viên học vụ.
        \end{itemize}

        \item \textbf{Purpose:}
        \begin{enumerate}
            \item Cho phép bật/tắt Email, SMS và Push App.
            \item Lưu cấu hình riêng cho từng người dùng.
        \end{enumerate}

        \item \textbf{Data processing:}
        \begin{enumerate}
            \item Hệ thống tải \texttt{CauHinhThongBao} của người dùng hiện tại.
            \item Người dùng thay đổi \texttt{NhanEmail}, \texttt{NhanSMS}, \texttt{NhanPushApp}.
            \item Hệ thống xác minh cấu hình thuộc chính tài khoản.
            \item Mỗi người dùng có một cấu hình thông báo theo quan hệ 1--1.
            \item Sau khi lưu, Notification Service sử dụng cấu hình để chọn kênh gửi.
            \item Việc tắt một kênh không xóa thông báo In-app đã tồn tại.
        \end{enumerate}

    \end{itemize}

    \item \textbf{Business rules:}

    BR-59.

    \ScreenLayoutItem

    \ScreenLayoutFigure{figures/03_02/new_2026_09_10/notification_preferences.png}{Configure Notification Preferences}

\end{itemize}

\subsubsection{Create and Send a Notification}
\label{subsubsec:create-send-notification}

\begin{itemize}

    \item \textbf{Function trigger:}

    Quản trị viên/Hệ thống tạo thông báo do sự kiện học tập, cảnh báo hoặc thao tác gửi.

    \item \textbf{Function description:}

    \begin{itemize}

        \item \textbf{Actors/Roles:}
        \begin{itemize}
            \item Quản trị viên học vụ.
            \item Notification Service.
        \end{itemize}

        \item \textbf{Purpose:}
        \begin{enumerate}
            \item Tạo nội dung và phân phối tới người nhận/nhóm.
            \item Theo dõi trạng thái từng người nhận/kênh.
        \end{enumerate}

        \item \textbf{Data processing:}
        \begin{enumerate}
            \item Hệ thống tạo \texttt{ThongBao} với người gửi, đối tượng nhận, tiêu đề, nội dung và đường dẫn nếu có.
            \item Nếu gửi theo nhóm, hệ thống xác định tài khoản phù hợp với \texttt{DoiTuongNhan} như ALL, HOC\_SINH, GIAO\_VIEN, PHU\_HUYNH.
            \item Với mỗi người nhận, hệ thống tạo \texttt{ThongBaoNguoiNhan}.
            \item Hệ thống kiểm tra \texttt{CauHinhThongBao} để xác định Email, SMS, Push được bật.
            \item Mỗi lần gửi theo kênh bắt đầu với trạng thái Chờ xử lý.
            \item Khi thành công chuyển Đã gửi; khi thất bại chuyển Thất bại và ghi \texttt{ThongTinLoi}.
            \item Thông báo In-app vẫn gắn với người nhận để hiển thị trong Notification Center.
        \end{enumerate}

    \end{itemize}

    \item \textbf{Business rules:}

    BR-56, BR-57, BR-58, BR-59, BR-60, BR-70, BR-71.

    \ScreenLayoutItem

    \ScreenLayoutFigure{figures/03_02/new_2026_09_10/teacher_send_notification.png}{Create and Send a Notification}

\end{itemize}

\subsubsection{Schedule a Notification}
\label{subsubsec:schedule-notification}

\begin{itemize}

    \item \textbf{Function trigger:}

    Quản trị viên/Hệ thống tạo thông báo cần gửi tại thời điểm/tần suất xác định.

    \item \textbf{Function description:}

    \begin{itemize}

        \item \textbf{Actors/Roles:}
        \begin{itemize}
            \item Quản trị viên học vụ.
            \item Notification Service.
        \end{itemize}

        \item \textbf{Purpose:}
        \begin{enumerate}
            \item Thiết lập gửi tự động theo lịch.
            \item Hỗ trợ nhắc nhở học tập và sự kiện hệ thống.
        \end{enumerate}

        \item \textbf{Data processing:}
        \begin{enumerate}
            \item Người tạo xác định nội dung, đối tượng nhận, thời điểm và tần suất theo chức năng triển khai.
            \item Hệ thống kiểm tra thời gian lịch gửi hợp lệ.
            \item Thông tin lịch được lưu trong thành phần lập lịch của Notification Service.
            \item Tại thời điểm kích hoạt, hệ thống xác định lại tập người nhận.
            \item Quá trình tạo \texttt{ThongBaoNguoiNhan} và gửi theo kênh thực hiện như chức năng gửi thông thường.
            \item Trạng thái và lỗi của từng lần gửi được ghi nhận.
        \end{enumerate}

    \end{itemize}

    \item \textbf{Business rules:}

    BR-56, BR-59, BR-60, BR-70.

    \ScreenLayoutItem

    \ScreenLayoutFigure{figures/03_02/admin_dashboard.png}{Schedule a Notification}

\end{itemize}

\subsubsection{View Notification Delivery Status and Error Report}
\label{subsubsec:view-notification-delivery-status}

\begin{itemize}

    \item \textbf{Function trigger:}

    Quản trị viên hoặc thành phần giám sát Notification Service mở báo cáo trạng thái/lỗi.

    \item \textbf{Function description:}

    \begin{itemize}

        \item \textbf{Actors/Roles:}
        \begin{itemize}
            \item Quản trị viên học vụ.
            \item Notification Service.
        \end{itemize}

        \item \textbf{Purpose:}
        \begin{enumerate}
            \item Theo dõi số lượng gửi thành công, chờ xử lý và thất bại.
            \item Xem lỗi theo người nhận/kênh.
        \end{enumerate}

        \item \textbf{Data processing:}
        \begin{enumerate}
            \item Hệ thống tải \texttt{ThongBaoNguoiNhan} trong phạm vi thông báo/thời gian được chọn.
            \item Hệ thống tổng hợp theo \texttt{TrangThaiGui} và \texttt{KenhGui}.
            \item Danh sách chi tiết có thể hiển thị người nhận, kênh, thời gian gửi, trạng thái đọc và thông tin lỗi.
            \item Quản trị viên có thể lọc theo trạng thái, kênh, vai trò hoặc thời gian khi hỗ trợ.
            \item Lỗi gửi lấy từ \texttt{ThongTinLoi}; hệ thống không thay lỗi thật bằng dữ liệu suy đoán.
            \item Trạng thái đọc được theo dõi độc lập với trạng thái gửi.
        \end{enumerate}

    \end{itemize}

    \item \textbf{Business rules:}

    BR-60, BR-70, BR-71.

    \ScreenLayoutItem

    \ScreenLayoutFigure{figures/03_02/admin_analytics.png}{View Notification Delivery Status and Error Report}

\end{itemize}


% =========================================================
% III. SOFTWARE REQUIREMENT SPECIFICATION - PHẦN 4 VÀ 5
% Cấu trúc bám theo file mẫu PDCMS.
% Nội dung chuẩn hóa từ Product Overview, Use Case, ERD và Functional Requirements 3.1/3.2 hiện có.
% =========================================================

\clearpage

% =========================================================
% 4. NON-FUNCTIONAL REQUIREMENTS
% =========================================================
\section*{4. Non-Functional Requirements}
\addcontentsline{toc}{section}{4. Non-Functional Requirements}

\subsection*{4.1 External Interfaces}
\addcontentsline{toc}{subsection}{4.1 External Interfaces}

Phần này xác định các giao diện giữa hệ thống \textbf{AI-Edu Pathways} với người dùng và các dịch vụ/hệ thống bên ngoài. Các yêu cầu được tổng hợp từ Product Overview, System Actors, Use Case và các luồng xử lý đã đặc tả trong phần Functional Requirements.

\subsubsection*{4.1.1 User Interfaces (UI)}
\addcontentsline{toc}{subsubsection}{4.1.1 User Interfaces (UI)}

\begin{itemize}
    \item \textbf{Web Application:} Hệ thống cung cấp giao diện Web responsive và hỗ trợ các trình duyệt hiện đại như Chrome, Edge, Firefox và Safari. Phiên bản hiện tại không yêu cầu chế độ offline.
    \item \textbf{Role-based Interfaces:} Sau khi xác thực, người dùng được điều hướng tới Dashboard tương ứng với vai trò Học sinh/Sinh viên, Giáo viên/Giảng viên, Phụ huynh hoặc Quản trị viên học vụ. Giao diện và dữ liệu phải tuân theo Role Authorization và phạm vi sở hữu/liên kết.
    \item \textbf{Student Experience:} Giao diện người học phải cung cấp các điểm truy cập rõ ràng tới lộ trình học cá nhân hóa, tài liệu, luyện tập/thi, AI Tutor, kết quả/phân tích năng lực, thông báo và cài đặt cá nhân.
    \item \textbf{Teacher Experience:} Giao diện Giáo viên phải hỗ trợ lớp học, ngân hàng câu hỏi, đề thi, giao bài, chấm bài, đánh giá tiến độ và Learning Analytics trong phạm vi phụ trách.
    \item \textbf{Parent Experience:} Giao diện Phụ huynh phải tập trung vào hồ sơ con đã liên kết, báo cáo tiến độ, kết quả, lịch học/chuyên cần, cảnh báo, khuyến nghị, giới hạn thời gian học, trao đổi với Giáo viên và thanh toán khi được phép.
    \item \textbf{Administrator Experience:} Giao diện Quản trị viên phải cung cấp quản lý tài khoản/phân quyền, nội dung học thuật, cấu hình AI, báo cáo, log bảo mật, thanh toán và thông báo theo quyền quản trị.
    \item \textbf{Consistency and Feedback:} Các nút, biểu mẫu, bảng, trạng thái, thông báo lỗi/thành công và cách điều hướng phải nhất quán giữa các module. Khi chưa có dữ liệu, đang xử lý hoặc dịch vụ ngoài gặp lỗi, hệ thống phải thể hiện trạng thái rõ ràng thay vì hiển thị dữ liệu giả.
\end{itemize}

\subsubsection*{4.1.2 Hardware Interfaces}
\addcontentsline{toc}{subsubsection}{4.1.2 Hardware Interfaces}

\begin{itemize}
    \item Hệ thống là Web Application và \textbf{không yêu cầu giao tiếp trực tiếp với phần cứng chuyên dụng} trong phạm vi hiện tại.
    \item Người dùng truy cập bằng máy tính để bàn, laptop, máy tính bảng hoặc thiết bị có trình duyệt Web hiện đại và kết nối Internet ổn định.
    \item Các thiết bị ngoại vi như camera, microphone hoặc máy in không phải là thành phần bắt buộc của các luồng nghiệp vụ đã đặc tả hiện tại.
\end{itemize}

\subsubsection*{4.1.3 Software Interfaces (APIs \& Integrations)}
\addcontentsline{toc}{subsubsection}{4.1.3 Software Interfaces (APIs \& Integrations)}

\begin{itemize}
    \item \textbf{AI / LLM Service:} Hệ thống tích hợp dịch vụ AI như OpenAI hoặc Gemini để xử lý prompt, sinh phản hồi AI Tutor, hỗ trợ sinh câu hỏi, giải thích lời giải và các chức năng phân tích/dự đoán. Hệ thống phải xử lý có kiểm soát các trường hợp timeout, rate limit hoặc lỗi dịch vụ và không ghi phản hồi giả khi lời gọi AI thất bại.
    \item \textbf{Vector Database / RAG:} Hệ thống tích hợp cơ sở dữ liệu Vector như Qdrant hoặc Pinecone để tạo/truy vấn embedding, thực hiện Semantic Search và cung cấp context cho RAG trước khi gọi LLM. Ngữ cảnh truy xuất phải thuộc nguồn dữ liệu được hệ thống cho phép sử dụng.
    \item \textbf{Payment Gateway:} Hệ thống tích hợp cổng thanh toán nội địa như VNPay, MoMo hoặc ZaloPay. Yêu cầu giao dịch phải gắn với đơn hàng, số tiền và người dùng tương ứng; kết quả từ Gateway/Webhook/IPN phải được xác minh trước khi cập nhật trạng thái và kích hoạt quyền lợi.
    \item \textbf{Notification Service:} Hệ thống tích hợp các kênh Email, SMS và Push Notification để gửi nhắc nhở, cảnh báo, thông báo học tập và thông báo hệ thống. Trạng thái gửi/lỗi phải được theo dõi theo từng người nhận và kênh gửi.
    \item \textbf{Education Organization Interface:} Hệ thống có khả năng đồng bộ với nguồn dữ liệu của Tổ chức Giáo dục để nhận đề thi ĐGNL, khung năng lực, ngân hàng câu hỏi và metadata/knowledge tags; đồng thời hỗ trợ trao đổi báo cáo thống kê tổng hợp theo phạm vi tích hợp được triển khai.
\end{itemize}

\subsection*{4.2 Quality Attributes}
\addcontentsline{toc}{subsection}{4.2 Quality Attributes}

\subsubsection*{4.2.1 Maintainability}
\addcontentsline{toc}{subsubsection}{4.2.1 Maintainability}

\begin{itemize}
    \item Hệ thống phải được tổ chức theo các module nghiệp vụ tách biệt như Authentication, Student, Teacher, Parent, Learning/Examination, AI/RAG, Payment, Notification và Administration để giảm ảnh hưởng lan truyền khi thay đổi.
    \item Giao diện, API, xử lý nghiệp vụ, dữ liệu và tích hợp dịch vụ ngoài phải có ranh giới trách nhiệm rõ ràng; cấu hình nhà cung cấp bên ngoài không được rải rác trong logic nghiệp vụ.
    \item Các API và quy tắc nghiệp vụ quan trọng phải được tài liệu hóa đủ để hỗ trợ phát triển, review và kiểm thử.
    \item Mã nguồn và tài liệu phải được quản lý bằng Git/GitHub, sử dụng branch và Pull Request theo quy trình quản lý dự án đã xác định; thay đổi quan trọng phải có khả năng truy vết.
\end{itemize}

\subsubsection*{4.2.2 Usability}
\addcontentsline{toc}{subsubsection}{4.2.2 Usability}

\begin{itemize}
    \item Giao diện phải dễ hiểu đối với các nhóm người dùng có mức độ kỹ thuật khác nhau, đặc biệt là Học sinh/Sinh viên và Phụ huynh.
    \item Dashboard phải ưu tiên thông tin và thao tác liên quan trực tiếp tới vai trò hiện tại, hạn chế hiển thị chức năng ngoài quyền hoặc ngoài ngữ cảnh.
    \item Lộ trình học, tiến độ, kết quả thi và phân tích năng lực phải được trình bày theo cách trực quan, cho phép người dùng nhận biết trạng thái hiện tại và bước tiếp theo.
    \item AI Tutor phải sử dụng luồng hội thoại rõ ràng, phân biệt tin nhắn của người dùng và AI, đồng thời duy trì lịch sử trong phạm vi hội thoại của người học.
    \item Các thao tác có hậu quả quan trọng như nộp bài, thay đổi cấu hình, khóa tài khoản hoặc xác nhận thanh toán phải có trạng thái/xác nhận phù hợp trước khi hoàn tất.
    \item Thông báo lỗi phải thân thiện, không hiển thị stack trace, secret, thông tin xác thực hoặc chi tiết kỹ thuật không cần thiết cho người dùng cuối.
\end{itemize}

\subsubsection*{4.2.3 Reliability}
\addcontentsline{toc}{subsubsection}{4.2.3 Reliability}

\begin{itemize}
    \item Hệ thống phải bảo toàn dữ liệu bài làm đã lưu, kết quả đã nộp, tiến độ lộ trình và lịch sử giao dịch khi có thao tác cập nhật đồng thời hoặc lỗi tạm thời.
    \item Khi hết thời gian làm bài, hệ thống phải kết thúc/nộp bài theo dữ liệu đã lưu; sau khi nộp, câu trả lời không được tiếp tục chỉnh sửa.
    \item Một giao dịch thanh toán chỉ được kích hoạt quyền lợi sau khi kết quả thành công được xác nhận hợp lệ; trạng thái Chờ xử lý hoặc Thất bại không được cấp quyền lợi.
    \item Trạng thái gửi thông báo và trạng thái đọc phải được theo dõi độc lập theo từng người nhận; lỗi giao hàng phải phản ánh lỗi thực tế được ghi nhận.
    \item Khi AI, RAG, Payment Gateway hoặc Notification Service không khả dụng, hệ thống phải xử lý lỗi có kiểm soát, giữ dữ liệu nhất quán và không tự tạo kết quả thành công giả.
\end{itemize}

\subsubsection*{4.2.4 Performance}
\addcontentsline{toc}{subsubsection}{4.2.4 Performance}

\begin{itemize}
    \item Các thao tác Web thông thường như điều hướng Dashboard, tải danh sách, xem hồ sơ, xem lộ trình và truy vấn báo cáo phải phản hồi đủ nhanh để không làm gián đoạn luồng sử dụng trong điều kiện mạng bình thường.
    \item Các tác vụ phụ thuộc AI/RAG có thể có độ trễ cao hơn các truy vấn dữ liệu thông thường; trong thời gian xử lý, giao diện phải hiển thị trạng thái đang xử lý và cho phép xử lý lỗi/timeout rõ ràng.
    \item Các luồng nộp bài, thanh toán và cập nhật trạng thái quan trọng phải ưu tiên tính đúng đắn và tính nhất quán dữ liệu, không được đánh đổi tính toàn vẹn để giảm thời gian phản hồi.
    \item Các báo cáo/tệp xuất có tập dữ liệu lớn phải được tạo theo cơ chế không làm treo giao diện và chỉ phát hành tệp khi quá trình tạo hoàn tất thành công.
\end{itemize}

\noindent\textit{Ghi chú: Các tài liệu hiện có của dự án chưa chốt SLA định lượng như thời gian phản hồi tối đa, uptime hoặc số người dùng đồng thời. Vì vậy phiên bản SRS này không tự đặt các con số giả định; các ngưỡng đo lường sẽ được bổ sung sau khi kiến trúc triển khai và môi trường kiểm thử được chốt.}

\subsubsection*{4.2.5 Design Constraints}
\addcontentsline{toc}{subsubsection}{4.2.5 Design Constraints}

\begin{itemize}
    \item Hệ thống được định hướng là Web-based, responsive và yêu cầu kết nối Internet liên tục; phiên bản hiện tại không hỗ trợ offline.
    \item Cá nhân hóa lộ trình và phân tích năng lực phải dựa trên dữ liệu học tập, khung năng lực và ngân hàng câu hỏi; hệ thống không được tạo kết luận không có dữ liệu hỗ trợ.
    \item AI chỉ đóng vai trò hỗ trợ học tập và đề xuất; quyết định học thuật/chuyên môn cuối cùng thuộc Giáo viên hoặc quy trình phê duyệt được cấu hình.
    \item Hệ thống phụ thuộc vào tính sẵn sàng, latency và rate limits của AI API, Vector Database, Payment Gateway và Notification Service; thiết kế phải có cơ chế xử lý khi phụ thuộc ngoài gặp lỗi.
    \item PMP hiện ghi nhận các lựa chọn công nghệ \textit{dự kiến} như React/Next.js, Python/FastAPI, PostgreSQL và Docker. Đây chưa được xem là ràng buộc công nghệ cuối cùng cho SDD cho tới khi nhóm chốt kiến trúc và source triển khai.
\end{itemize}

\subsubsection*{4.2.6 Security}
\addcontentsline{toc}{subsubsection}{4.2.6 Security}

\begin{itemize}
    \item Hệ thống phải áp dụng Role-Based Access Control (RBAC) ở cả giao diện và API; việc ẩn menu trên UI không được thay thế kiểm tra quyền phía server.
    \item Học sinh chỉ được truy cập dữ liệu của chính mình; Giáo viên chỉ được truy cập lớp/học sinh thuộc phạm vi phụ trách; Phụ huynh chỉ được truy cập học sinh có liên kết hợp lệ; các chức năng quản trị toàn hệ thống chỉ dành cho vai trò có quyền.
    \item Mật khẩu không được lưu hoặc xử lý dưới dạng văn bản thuần. Thông báo đăng nhập thất bại không được tiết lộ tài khoản nào tồn tại hoặc thông tin nhạy cảm khác.
    \item Dữ liệu giao dịch thanh toán phải được truyền qua kênh bảo mật và trạng thái giao dịch phải được xác minh từ cổng thanh toán trước khi cấp quyền lợi.
    \item Secret/API key của AI, Payment và Notification Service không được đưa ra phía client hoặc ghi trực tiếp vào thông báo/log dành cho người dùng cuối.
    \item Các thay đổi cấu hình AI và các sự kiện quản trị/bảo mật quan trọng phải có dữ liệu audit phù hợp để truy vết người thực hiện, thời gian và nội dung thay đổi.
    \item Dữ liệu gửi vào AI/RAG phải tuân theo cùng phạm vi truy cập của người dùng và nguồn dữ liệu được phép; nội dung truy xuất không được dùng để mở rộng quyền truy cập ngoài phạm vi đã xác thực.
\end{itemize}

% =========================================================
% 5. REQUIREMENT APPENDIX
% =========================================================
\clearpage
\section*{5. Requirement Appendix}
\addcontentsline{toc}{section}{5. Requirement Appendix}

\subsection*{5.1 Business Rules}
\addcontentsline{toc}{subsection}{5.1 Business Rules}

Các Business Rule dưới đây chuẩn hóa các mã \texttt{BR-01}--\texttt{BR-71} đã được tham chiếu trong phần 3.2 Web Application. Định nghĩa được hợp nhất từ chính luồng xử lý, phạm vi dữ liệu, ERD và điều kiện nghiệp vụ đã mô tả trong các chức năng tương ứng; không bổ sung giới hạn số học chưa có nguồn trong tài liệu dự án.

\renewcommand{\arraystretch}{1.15}
\begin{longtable}{|p{1.6cm}|p{12.7cm}|}
\hline
\rowcolor{orange!15}
\textbf{ID} & \textbf{Rule Definition} \\
\hline
\endfirsthead
\hline
\rowcolor{orange!15}
\textbf{ID} & \textbf{Rule Definition} \\
\hline
\endhead
BR-01 & Chỉ phiên đăng nhập hợp lệ mới được sử dụng để truy cập chức năng bảo vệ; thao tác đăng xuất phải kết thúc/vô hiệu hóa phiên hiện tại. \\ \hline
BR-02 & Tài khoản bị khóa hoặc không hoạt động không được tạo phiên đăng nhập mới cho tới khi trạng thái được khôi phục bởi người có quyền. \\ \hline
BR-03 & Mật khẩu phải được xác minh bằng cơ chế bảo mật phù hợp, không xử lý/lưu dưới dạng văn bản thuần và không để lỗi đăng nhập tiết lộ thông tin nhạy cảm. \\ \hline
BR-04 & Mỗi phiên người dùng phải gắn với một vai trò hợp lệ của tài khoản; tài khoản Học sinh đăng ký mới được khởi tạo với vai trò Học sinh và người dùng không được tự nâng quyền. \\ \hline
BR-05 & Sau đăng nhập thành công, hệ thống điều hướng người dùng tới Dashboard tương ứng với vai trò và chỉ hiển thị chức năng thuộc phạm vi quyền đó. \\ \hline
BR-06 & Quản lý tài khoản toàn hệ thống, khóa/mở khóa, gán vai trò và xem log bảo mật là chức năng quản trị, chỉ người có quyền Quản trị viên học vụ mới được thực hiện. \\ \hline
BR-07 & Mọi truy vấn và thao tác dữ liệu phải kiểm tra cả vai trò lẫn phạm vi sở hữu/liên kết; việc thay đổi ID/tham số request không được mở rộng quyền truy cập. \\ \hline
BR-08 & Email và số điện thoại dùng cho tài khoản phải đúng định dạng và không trùng trong phạm vi yêu cầu duy nhất; Học sinh chỉ được chỉnh các trường hồ sơ thuộc quyền của chính mình. \\ \hline
BR-09 & Giáo viên chỉ được xem và thao tác lớp/buổi học thuộc phạm vi mình phụ trách, trừ khi có quyền quản trị rộng hơn. \\ \hline
BR-10 & Danh sách học sinh của lớp phải được xác định từ quan hệ tham gia lớp hợp lệ; học sinh không có quan hệ với lớp không được xuất hiện như thành viên hiện hành. \\ \hline
BR-11 & Dữ liệu học sinh, bài giao, bài nộp, nhận xét và báo cáo mà Giáo viên truy cập phải giới hạn trong lớp/phạm vi phụ trách. \\ \hline
BR-12 & Phụ huynh chỉ được khởi tạo chức năng giám sát hoặc trao đổi về một học sinh khi tồn tại liên kết Phụ huynh--Học sinh hợp lệ. \\ \hline
BR-13 & Mọi báo cáo, kết quả, cảnh báo, lịch học và phân tích dành cho Phụ huynh phải được giới hạn theo học sinh đã liên kết; không được truy cập hồ sơ con của tài khoản khác. \\ \hline
BR-14 & Dashboard và báo cáo Phụ huynh phải xác định hồ sơ con từ liên kết hợp lệ và chỉ tổng hợp dữ liệu thuộc hồ sơ được chọn. \\ \hline
BR-15 & Mỗi liên kết Phụ huynh--Học sinh chỉ có một cấu hình giới hạn thời gian học đang áp dụng theo quan hệ dữ liệu 1--1. \\ \hline
BR-16 & Khung giờ bắt đầu/kết thúc và tổng thời gian tối đa trong cấu hình giới hạn phải có định dạng và giá trị hợp lệ trước khi lưu. \\ \hline
BR-17 & Mục tiêu học tập/kỳ thi đích phải hợp lệ và được dùng làm đầu vào cho dự đoán kết quả cũng như đánh giá mức độ đạt mục tiêu. \\ \hline
BR-18 & Việc tạo hoặc điều chỉnh Roadmap phải kết hợp mục tiêu hiện tại với dữ liệu năng lực/kết quả/tiến độ của chính người học. \\ \hline
BR-19 & Roadmap ban đầu chỉ được tạo hoàn chỉnh khi có đủ dữ liệu đánh giá đầu vào và mục tiêu/kỳ thi đích cần thiết; nếu thiếu phải yêu cầu bổ sung thay vì suy đoán. \\ \hline
BR-20 & Roadmap phải thuộc đúng người học và được tạo từ kết quả đánh giá/mục tiêu của người đó; người học không được truy cập Roadmap của tài khoản khác. \\ \hline
BR-21 & Roadmap phải duy trì thứ tự/nội dung học và trạng thái tiến độ; thay đổi chương trình/lộ trình không được làm mất lịch sử tiến độ đã phát sinh. \\ \hline
BR-22 & Thay đổi mục tiêu hoặc chương trình có thể kích hoạt đánh giá/điều chỉnh Roadmap, nhưng phải bảo toàn dữ liệu lịch sử đã hoàn thành. \\ \hline
BR-23 & Dữ liệu tiến độ, thời gian học, kết quả và trạng thái chưa đạt mục tiêu là đầu vào cho báo cáo, nhắc nhở và điều chỉnh Roadmap. \\ \hline
BR-24 & Báo cáo, khuyến nghị và phân tích học tập phải được tạo từ dữ liệu thực tế trong phạm vi được phép; khi không có dữ liệu hệ thống phải thể hiện trạng thái rỗng thay vì tự tạo số liệu. \\ \hline
BR-25 & Bài học/tài liệu được đề xuất hoặc mở từ Roadmap phải tồn tại, còn khả dụng và có quan hệ phù hợp với nội dung/lộ trình đang sử dụng. \\ \hline
BR-26 & Mỗi hội thoại AI Tutor phải gắn với đúng người học và chỉ người học đó được mở lịch sử hội thoại của mình. \\ \hline
BR-27 & Khi chức năng AI Tutor sử dụng RAG, hệ thống phải thực hiện truy xuất ngữ cảnh/tài liệu liên quan trước khi gọi LLM theo cấu hình đang áp dụng. \\ \hline
BR-28 & Hội thoại AI Tutor phải được lưu và duy trì theo phiên/hội thoại để người học có thể mở lại lịch sử trong phạm vi của mình. \\ \hline
BR-29 & Prompt AI Tutor chỉ được kết hợp câu hỏi với ngữ cảnh được phép như lịch sử hội thoại, môn/chủ đề, cấp độ và dữ liệu học tập liên quan. \\ \hline
BR-30 & Nếu RAG hoặc LLM thất bại, hệ thống không được ghi phản hồi AI giả là kết quả thành công và phải cho phép xử lý/thử lại phù hợp. \\ \hline
BR-31 & Đánh giá/phản hồi của người học phải gắn với đúng nội dung hoặc phản hồi AI đang được đánh giá và không tự động sửa nội dung học thuật chính thức. \\ \hline
BR-32 & Tin nhắn AI Tutor phải phân biệt rõ người gửi User/AI và được lưu trong đúng hội thoại để duy trì lịch sử. \\ \hline
BR-33 & Kết quả/đề xuất do AI sinh phải tuân theo cấu hình AI đang áp dụng và chỉ đóng vai trò hỗ trợ; nội dung cần kiểm duyệt không được tự động trở thành nội dung chính thức. \\ \hline
BR-34 & Mọi thay đổi cấu hình AI quan trọng phải được ghi log với người thực hiện, hành động, giá trị cũ, giá trị mới và thời gian. \\ \hline
BR-35 & Bài luyện thích ứng/đề thi hiển thị cho người học phải thuộc phạm vi được phép và việc tạo bài thích ứng phải dựa trên năng lực/mục tiêu hiện tại hoặc mức khởi đầu có căn cứ. \\ \hline
BR-36 & Cấu trúc đề/bài luyện, lựa chọn câu hỏi, độ khó, ma trận kiến thức, thứ tự và điểm thành phần phải tuân theo cấu hình đề và dữ liệu câu hỏi hợp lệ. \\ \hline
BR-37 & Một bài làm phát sinh từ assignment phải gắn đúng học sinh, đề thi và \texttt{MaGiaoBaiTap} để Giáo viên có thể đối chiếu bài nộp. \\ \hline
BR-38 & Người học chỉ được bắt đầu/làm đề khi đề còn khả dụng và trong điều kiện thời gian cho phép; sau khi hết thời gian hoặc bài đã nộp không được tiếp tục ghi câu trả lời. \\ \hline
BR-39 & \texttt{BaiLam} chỉ được tạo khi người học đủ điều kiện bắt đầu; thời điểm bắt đầu/nộp và trạng thái bài làm là nguồn xác định hợp lệ cho danh sách submission. \\ \hline
BR-40 & Phần bài cần chấm thủ công chỉ được Giáo viên có quyền chấm/nhận xét và điểm được nhập phải nằm trong phạm vi điểm thành phần. \\ \hline
BR-41 & Sau khi nộp, câu trả lời bị khóa; điểm chi tiết và tổng điểm phải được tính từ dữ liệu bài làm/đáp án/điểm chấm hợp lệ trước khi công bố kết quả. \\ \hline
BR-42 & Kết quả bài làm sau khi chấm là nguồn dữ liệu để cập nhật phân tích năng lực, Knowledge Tracing, dự đoán, báo cáo và điều chỉnh cá nhân hóa. \\ \hline
BR-43 & Kết quả chi tiết chỉ được hiển thị khi bài đã đạt trạng thái chấm/kết quả phù hợp; dữ liệu chưa hoàn tất phải hiển thị trạng thái chờ thay vì kết quả giả. \\ \hline
BR-44 & Quản lý bài học, tài liệu, ma trận kiến thức, câu hỏi và nội dung học thuật chỉ được thực hiện bởi vai trò có quyền và trong phạm vi được phân công. \\ \hline
BR-45 & Cấu hình đề thi phải nhất quán về câu hỏi, đáp án, điểm, thứ tự và thời gian; thay đổi sau khi đã có bài làm phải không làm sai lệch dữ liệu lịch sử. \\ \hline
BR-46 & Đánh giá/nhận xét cuối cùng về học sinh phải do Giáo viên có quyền thực hiện; AI chỉ cung cấp dữ liệu hỗ trợ và không tự thay thế quyết định chuyên môn. \\ \hline
BR-47 & Nội dung/câu hỏi/đề thuộc quy trình kiểm duyệt chỉ được xuất bản hoặc cho người học sử dụng sau khi đạt trạng thái phê duyệt theo chính sách. \\ \hline
BR-48 & Quản lý danh mục dùng chung, cấu hình hệ thống/AI, báo cáo quản trị, log và các thao tác quản trị toàn hệ thống phải tuân theo quyền Quản trị viên và kiểm tra tính toàn vẹn tham chiếu trước khi thay đổi. \\ \hline
BR-49 & Chỉ gói dịch vụ đang hoạt động mới được tạo giao dịch mới; quyền lợi được cấp phải tương ứng với gói đã mua và còn hiệu lực. \\ \hline
BR-50 & Đơn hàng, tổng tiền và lịch sử giao dịch phải gắn với đúng tài khoản; người dùng chỉ được xem giao dịch của chính mình. \\ \hline
BR-51 & Khi áp dụng voucher, hệ thống phải tìm thực thể \texttt{MaGiamGia} theo \texttt{CodeVoucher}, xác định khóa chính \texttt{MaVoucher} và chỉ áp dụng khi thời hạn/điều kiện hợp lệ. \\ \hline
BR-52 & Chi tiết đơn hàng phải lưu số lượng và \texttt{DonGiaAtTime}; tổng tiền cuối cùng phải phản ánh giá tại thời điểm Checkout và voucher hợp lệ đã áp dụng. \\ \hline
BR-53 & Chỉ giao dịch được xác nhận thành công mới được kích hoạt quyền lợi và phát hành/cho phép xem hóa đơn điện tử. \\ \hline
BR-54 & Trạng thái giao dịch phải được cập nhật từ kết quả cổng thanh toán thành Chờ xử lý/Thành công/Thất bại; không được kích hoạt quyền lợi từ kết quả chưa xác nhận hoặc thất bại. \\ \hline
BR-55 & Yêu cầu hoàn tiền chỉ được tạo cho giao dịch thuộc người dùng, đã thanh toán thành công và đáp ứng chính sách hoàn tiền đang áp dụng. \\ \hline
BR-56 & Thông báo học tập/cảnh báo/tin nhắn chỉ được tạo cho đúng người nhận hoặc nhóm người nhận hợp lệ theo sự kiện và quan hệ dữ liệu liên quan. \\ \hline
BR-57 & Khi gửi theo nhóm, hệ thống phải xác định tập người nhận từ đối tượng nhận được cấu hình trước khi tạo bản ghi người nhận và phân phối. \\ \hline
BR-58 & Mỗi người nhận phải có bản ghi thông báo riêng; người dùng chỉ được xem/cập nhật trạng thái đọc của thông báo thuộc chính mình. \\ \hline
BR-59 & Việc gửi Email/SMS/Push phải tôn trọng cấu hình thông báo riêng của người dùng; tắt một kênh không xóa thông báo In-app đã tồn tại. \\ \hline
BR-60 & Trạng thái gửi và lỗi phải được ghi nhận theo từng người nhận/kênh; Chờ xử lý, Đã gửi và Thất bại phải được phân biệt rõ. \\ \hline
BR-61 & Bài được giao và danh sách bài nộp phải liên kết nhất quán giữa \texttt{GiaoBaiTap}, lớp, đề thi, học sinh và \texttt{BaiLam}. \\ \hline
BR-62 & Ngày giao và hạn nộp của assignment phải hợp lệ; hạn nộp không được mâu thuẫn với thời điểm giao. \\ \hline
BR-63 & Chỉ học sinh thuộc lớp được giao và trong thời gian/điều kiện khả dụng mới được truy cập assignment/đề tương ứng. \\ \hline
BR-64 & Buổi học phải gắn với lớp và Giáo viên/phạm vi hợp lệ; thao tác điểm danh chỉ được thực hiện trên buổi học mà người dùng có quyền quản lý. \\ \hline
BR-65 & Điểm danh phải được ghi nhận theo từng học sinh và buổi học; báo cáo chuyên cần phải tổng hợp từ các bản ghi điểm danh thực tế. \\ \hline
BR-66 & Phụ huynh chỉ có quyền xem lịch học/chuyên cần của học sinh đã liên kết và không được sửa dữ liệu điểm danh. \\ \hline
BR-67 & Hội thoại Phụ huynh--Giáo viên phải gắn với một học sinh cụ thể và chỉ các thành viên hợp lệ của hội thoại được xem/gửi tin nhắn. \\ \hline
BR-68 & Tin nhắn trao đổi phải gắn với đúng hội thoại và người gửi hợp lệ; lịch sử được sắp theo thời gian và trạng thái đọc được cập nhật cho tin nhắn nhận được. \\ \hline
BR-69 & Hóa đơn điện tử phải gắn với đúng đơn hàng thanh toán thành công và chỉ chủ đơn hàng được xem/tải hóa đơn của mình. \\ \hline
BR-70 & Thông báo theo lịch hoặc cảnh báo tự động phải được kích hoạt từ tiêu chí/thời gian được cấu hình và xác định lại tập người nhận tại thời điểm gửi khi luồng yêu cầu. \\ \hline
BR-71 & Trạng thái đọc và trạng thái gửi là hai trạng thái độc lập; báo cáo thông báo phải sử dụng dữ liệu thực tế của từng trạng thái và không thay lỗi gửi bằng dữ liệu suy đoán. \\ \hline
\end{longtable}

\subsection*{5.2 Common Requirements}
\addcontentsline{toc}{subsection}{5.2 Common Requirements}

\begin{longtable}{|p{1.6cm}|p{3.2cm}|p{9.2cm}|}
\hline
\rowcolor{orange!15}
\textbf{Req. ID} & \textbf{Requirement Name} & \textbf{Description} \\
\hline
\endfirsthead
\hline
\rowcolor{orange!15}
\textbf{Req. ID} & \textbf{Requirement Name} & \textbf{Description} \\
\hline
\endhead
CR-01 & Data Validation & Tất cả trường bắt buộc trên biểu mẫu phải được kiểm tra trước khi lưu/gửi. Các API ghi dữ liệu phải kiểm tra lại dữ liệu phía server, không chỉ dựa vào validation của UI. \\ \hline
CR-02 & Data Format \& Domain Validation & Email, số điện thoại, ngày/giờ, điểm số, thời lượng, trạng thái, ID tham chiếu và các giá trị enum phải đúng định dạng/miền hợp lệ trước khi xử lý. \\ \hline
CR-03 & Authentication \& Authorization & Mọi route/API bảo vệ phải xác minh phiên/token và vai trò; yêu cầu không hợp lệ hoặc hết hạn phải bị từ chối và điều hướng/xử lý theo cơ chế xác thực. \\ \hline
CR-04 & Data Scope / Ownership & Ngoài kiểm tra vai trò, hệ thống phải kiểm tra quyền sở hữu/phạm vi dữ liệu: Student--self, Teacher--assigned classes/students, Parent--linked child, Admin--authorized administrative scope. \\ \hline
CR-05 & Audit Logging & Các thay đổi cấu hình AI, hành động quản trị, sự kiện truy cập/bảo mật và thao tác nhạy cảm phải có log phù hợp để xác định người thực hiện, thời gian và nội dung thay đổi. \\ \hline
CR-06 & Error Handling & Hệ thống phải hiển thị thông báo lỗi dễ hiểu; không hiển thị stack trace, secret, mật khẩu, token hoặc thông tin nội bộ không cần thiết cho người dùng cuối. \\ \hline
CR-07 & Data Integrity \& Transaction Safety & Các luồng nộp bài, chấm điểm, giao bài, thanh toán, kích hoạt quyền lợi và cập nhật trạng thái phải tránh ghi dữ liệu mâu thuẫn/nhân đôi khi request được gửi lặp hoặc xử lý đồng thời. \\ \hline
CR-08 & Date \& Time Consistency & Thời gian bắt đầu/kết thúc, hạn nộp, thời gian nộp, lịch thông báo, thời gian đọc/gửi và ngày phát hành hóa đơn phải được lưu/xử lý nhất quán theo cấu hình thời gian của hệ thống. \\ \hline
CR-09 & Loading / Empty / Failure States & Khi dữ liệu chưa có, đang tải hoặc phụ thuộc ngoài thất bại, giao diện phải thể hiện đúng trạng thái và không tự tạo dữ liệu để lấp chỗ trống. \\ \hline
CR-10 & External Service Failure Handling & Lỗi/timeout/rate limit từ AI, Vector DB, Payment Gateway và Notification Service phải được ghi nhận và xử lý có kiểm soát; các luồng quan trọng không được chuyển sang trạng thái thành công chỉ vì request đã được gửi. \\ \hline
CR-11 & Export Authorization & Báo cáo/tệp PDF, Excel hoặc CSV chỉ được chứa dữ liệu mà người dùng đã được phép xem trên màn hình; việc đổi tham số export không được mở rộng phạm vi dữ liệu. \\ \hline
CR-12 & Referential Integrity & Trước thao tác cập nhật/xóa/vô hiệu hóa dữ liệu đang được tham chiếu bởi lớp, đề thi, bài làm, lộ trình, tài liệu, đơn hàng hoặc thông báo, hệ thống phải kiểm tra ảnh hưởng để bảo toàn tính nhất quán lịch sử. \\ \hline
\end{longtable}

\subsection*{5.3 Application Messages List}
\addcontentsline{toc}{subsection}{5.3 Application Messages List}

Danh sách dưới đây chuẩn hóa các thông báo cần thiết từ các trạng thái validation/thành công/thất bại đã xuất hiện trong phần 3.2. Nội dung câu chữ có thể được đồng bộ lại với UI implementation sau này, nhưng mã và ngữ cảnh nên được giữ ổn định để phục vụ kiểm thử.

{\scriptsize
\renewcommand{\arraystretch}{1.2}
\begin{longtable}{|p{0.7cm}|p{1.25cm}|p{2.0cm}|p{2.6cm}|p{6.4cm}|}
\hline
\rowcolor{orange!15}
\textbf{\#} & \textbf{Code} & \textbf{Message Type} & \textbf{Context} & \textbf{Content} \\
\hline
\endfirsthead
\hline
\rowcolor{orange!15}
\textbf{\#} & \textbf{Code} & \textbf{Message Type} & \textbf{Context} & \textbf{Content} \\
\hline
\endhead
1 & MSG01 & Inline & Empty result & Không có dữ liệu phù hợp. \\ \hline
2 & MSG02 & Inline (Error) & Required field & Trường \{field\_name\} là bắt buộc. \\ \hline
3 & MSG03 & Inline (Error) & Format validation & Giá trị của \{field\_name\} không đúng định dạng. \\ \hline
4 & MSG04 & Toast (Error) & Authorization & Bạn không có quyền thực hiện thao tác này. \\ \hline
5 & MSG05 & Toast (Info) & Session & Phiên đăng nhập đã hết hạn. Vui lòng đăng nhập lại. \\ \hline
6 & MSG06 & Inline (Error) & Login & Email/Số điện thoại hoặc mật khẩu không chính xác. \\ \hline
7 & MSG07 & Toast (Error) & Account status & Tài khoản hiện đang bị khóa hoặc không hoạt động. \\ \hline
8 & MSG08 & Toast (Error) & Role & Vai trò được chọn không phù hợp với tài khoản. \\ \hline
9 & MSG09 & Toast (Success) & Registration & Đăng ký tài khoản thành công. \\ \hline
10 & MSG10 & Inline (Error) & Duplicate account & Email hoặc số điện thoại đã được sử dụng. \\ \hline
11 & MSG11 & Toast (Success) & Profile & Cập nhật hồ sơ thành công. \\ \hline
12 & MSG12 & Toast (Error) & Data scope & Bạn không được phép truy cập dữ liệu này. \\ \hline
13 & MSG13 & Inline & No learning data & Chưa có dữ liệu học tập để hiển thị. \\ \hline
14 & MSG14 & Inline (Error) & Learning goal & Mục tiêu hoặc kỳ thi đích chưa hợp lệ. \\ \hline
15 & MSG15 & Toast (Info) & Roadmap prerequisite & Vui lòng hoàn thành đánh giá đầu vào và thiết lập mục tiêu trước. \\ \hline
16 & MSG16 & Toast (Success) & Roadmap & Đã tạo lộ trình học cá nhân hóa. \\ \hline
17 & MSG17 & Toast (Success) & Roadmap update & Lộ trình đã được cập nhật theo dữ liệu mới nhất. \\ \hline
18 & MSG18 & Inline (Info) & AI Tutor & AI Tutor đang xử lý yêu cầu... \\ \hline
19 & MSG19 & Toast (Error) & AI/RAG failure & Không thể nhận phản hồi AI lúc này. Vui lòng thử lại. \\ \hline
20 & MSG20 & Toast (Success) & Feedback & Đã ghi nhận phản hồi của bạn. \\ \hline
21 & MSG21 & Toast (Error) & Exam availability & Bài luyện/đề thi hiện không khả dụng hoặc ngoài thời gian cho phép. \\ \hline
22 & MSG22 & Toast (Success) & Exam attempt & Đã bắt đầu bài làm. \\ \hline
23 & MSG23 & Toast (Info) & Exam timeout & Đã hết thời gian. Hệ thống đang nộp bài theo dữ liệu đã lưu. \\ \hline
24 & MSG24 & Confirm & Submit exam & Bạn có chắc chắn muốn nộp bài? Sau khi nộp sẽ không thể sửa câu trả lời. \\ \hline
25 & MSG25 & Toast (Success) & Submit exam & Nộp bài thành công. \\ \hline
26 & MSG26 & Inline (Info) & Grading & Bài làm đang chờ chấm/hoàn tất chấm điểm. \\ \hline
27 & MSG27 & Inline (Info) & Result & Kết quả chưa sẵn sàng. \\ \hline
28 & MSG28 & Toast (Success) & Learning content & Lưu nội dung học tập thành công. \\ \hline
29 & MSG29 & Toast (Info) & Approval & Nội dung đang chờ kiểm duyệt trước khi công bố. \\ \hline
30 & MSG30 & Toast (Success) & Assignment & Giao bài thành công. \\ \hline
31 & MSG31 & Inline (Error) & Assignment deadline & Hạn nộp không hợp lệ so với thời điểm giao. \\ \hline
32 & MSG32 & Toast (Success) & Grading & Đã lưu điểm và nhận xét. \\ \hline
33 & MSG33 & Toast (Error) & Parent link & Không có liên kết Phụ huynh--Học sinh hợp lệ cho hồ sơ này. \\ \hline
34 & MSG34 & Toast (Success) & Time limit & Đã lưu giới hạn thời gian học. \\ \hline
35 & MSG35 & Toast (Success) & Parent-Teacher message & Gửi tin nhắn thành công. \\ \hline
36 & MSG36 & Toast (Success) & Account status & Cập nhật trạng thái tài khoản thành công. \\ \hline
37 & MSG37 & Toast (Success) & Configuration & Cập nhật cấu hình thành công. \\ \hline
38 & MSG38 & Inline & Report & Không có dữ liệu trong phạm vi/bộ lọc đã chọn. \\ \hline
39 & MSG39 & Toast (Info) & Export & Đang tạo tệp báo cáo... \\ \hline
40 & MSG40 & Toast (Error) & Export & Không thể tạo tệp báo cáo. Vui lòng thử lại. \\ \hline
41 & MSG41 & Inline (Error) & Voucher & Mã giảm giá không hợp lệ, hết hạn hoặc không đáp ứng điều kiện áp dụng. \\ \hline
42 & MSG42 & Toast (Success) & Voucher & Áp dụng mã giảm giá thành công. \\ \hline
43 & MSG43 & Toast (Info) & Payment & Giao dịch đang chờ xử lý. \\ \hline
44 & MSG44 & Toast (Success) & Payment & Thanh toán thành công. Quyền lợi đang được kích hoạt. \\ \hline
45 & MSG45 & Toast (Error) & Payment & Thanh toán không thành công. Quyền lợi chưa được kích hoạt. \\ \hline
46 & MSG46 & Toast (Error) & Entitlement & Gói/quyền lợi chưa được kích hoạt hoặc đã hết hiệu lực. \\ \hline
47 & MSG47 & Inline (Info) & Invoice & Hóa đơn điện tử chưa sẵn sàng cho giao dịch này. \\ \hline
48 & MSG48 & Toast (Error) & Refund & Giao dịch không đáp ứng điều kiện hoàn tiền hiện tại. \\ \hline
49 & MSG49 & Toast (Success) & Notification preferences & Đã lưu cài đặt thông báo. \\ \hline
50 & MSG50 & Toast (Success) & Notification & Đã tạo/gửi thông báo theo cấu hình. \\ \hline
51 & MSG51 & Inline (Error) & Notification schedule & Thời điểm hoặc tần suất gửi thông báo không hợp lệ. \\ \hline
52 & MSG52 & Toast (Error) & Notification delivery & Gửi thông báo thất bại cho một hoặc nhiều người nhận. Vui lòng kiểm tra báo cáo lỗi. \\ \hline
\end{longtable}
}


% =========================================================
% IV. SOFTWARE DESIGN DESCRIPTION
% Giữ phần System Architecture / Package / Database Design từ bản nền,
% sau đó thay Detailed Design bằng các artifact draw.io do nhóm cung cấp.
% =========================================================
\clearpage
\includepdf[pages=153-159,fitpaper=true,pagecommand={\thispagestyle{empty}}]{base/AI_EDU_PATHWAYS_FINAL_289P_DIAGRAM_QA.pdf}

\input{chapters/04_detailed_design_from_drawio}

% =========================================================
% V. SOFTWARE TESTING DOCUMENTATION
% =========================================================
\clearpage
\includepdf[pages=273-289,fitpaper=true,pagecommand={\thispagestyle{empty}}]{base/AI_EDU_PATHWAYS_FINAL_289P_DIAGRAM_QA.pdf}

\end{document}
